\documentclass[11pt]{article}

\usepackage[T1]{fontenc}
\usepackage{lmodern}
\usepackage{amsmath,amssymb,amsthm,mathtools}
\usepackage{mathrsfs}
\usepackage{enumitem}
\usepackage{geometry}
\usepackage[colorlinks=true,linkcolor=blue,citecolor=blue,urlcolor=blue]{hyperref}

\geometry{a4paper,margin=28mm}
\setlength{\emergencystretch}{2em}

\newtheorem{theorem}{Theorem}[section]
\newtheorem{proposition}[theorem]{Proposition}
\newtheorem{lemma}[theorem]{Lemma}
\newtheorem{corollary}[theorem]{Corollary}
\theoremstyle{definition}
\newtheorem{definition}[theorem]{Definition}
\newtheorem{assumption}[theorem]{Assumption}
\theoremstyle{remark}
\newtheorem{remark}[theorem]{Remark}
\newtheorem{example}[theorem]{Example}

\newcommand{\Qell}{\mathbf Q_\ell}
\newcommand{\Spec}{\operatorname{Spec}}
\newcommand{\Spf}{\operatorname{Spf}}
\newcommand{\an}{\mathrm{an}}

\title{Local Kummer Realizations and Verdier Duality\\
\large Continuous cohomology, nearby cycles, and finite-root traces}
\author{YINGFENG JIANG}
\date{16 September 2026}

\begin{document}
\maketitle

\begin{abstract}
For a fixed semistable node with smooth disc directions, we construct
a geometric trace inverse system from genuine finite-root shriek
pushforwards.  Its integral limit is the derived dual of continuous
inertia cochains, with the geometric shift and Tate twist, and agrees
with actual nearby costalks and analytic compact support.
The construction retains full finite group bars, coefficient reduction,
module actions and integral torsion.  Its inputs are comparisons of
Kummer coefficients on the relative infinite-root fibre, tame nearby
fibre and analytic tube; on the selected perfect cartesian adic source,
the analytic realization is fully faithful and the three ordinary
images are continuous inertia cohomology.  All comparisons use the
fixed chart, base-root neutralization, compactification and carriers.
No shriek functor on the bare infinite-root fibre is defined.
\end{abstract}

\section{Introduction}
\label{frag:explicit-three-realization}

Let $xy=\pi$ be a semistable node.  Its vanishing tube is an open
annulus, and its relative pro-$\ell$ inertia is $\mathbf Z_\ell(1)$.
This paper compares Kummer coefficients on that tube, on the
logarithmic nearby fibre, and on the relative infinite-root fibre.
The ordinary images are continuous inertia cohomology.  The main
question is whether their supported partners agree as well.

The difficulty lies at finite root level.  A finite classifying
stack has unbounded cohomology, even though continuous inertia
cohomology admits a two-term complex.  We identify the geometric
trace transitions on the full finite bars and show that their
inverse limit agrees with analytic compact support and the nearby
costalk.  This construction retains the coefficient maps, module
actions and integral torsion.

Here $R_r=\mathbf Z/\ell^r$.  A source coefficient is a cartesian
system of discrete continuous inertia complexes, perfect over each
$R_r$ with one common Tor interval; the analytic target uses the
analogous lisse-perfect systems on the ordinary etale site.
Definitions~\ref{def:frag-local-adic-source-target}
and~\ref{def:frag-compatible-presenting-tails} specify the categories
and fixed-precision presentations; Lemma~\ref{lem:frag-presenting-system-assembly}
constructs their trace functors and coefficient maps.  The appendix proofs
of root-site descent and integral recovery are part of the argument.

\begin{theorem}[Local star comparison and finite-carrier trace]
\label{thm:frag-node-complete-local-datum}
Let $R$ be a complete strictly henselian discrete valuation ring with
algebraically closed residue field $k$, uniformizer $\pi$, and
$\ell\ne\operatorname{char}k$.  Fix $C=\widehat{\overline{\operatorname{Frac}R}}$.
Put
\[
 X=\operatorname{Spec}R[x,y,z_1,\ldots,z_e]/(xy-\pi),\quad
 Y=V(x,y,z_1,\ldots,z_e),\quad N=1+e.
\]
Use the vertical log structure, compatible base roots, and the
relative Kummer tower of $x$; these fix the splitting
$P^{\mathrm{gp}}=\mathbf Z\rho\oplus\mathbf Z e_x$.
Fix the compactification
$\overline X:\ X_0Y_0=\pi X_1Y_1$ in
$(\mathbf P^1_R)^2\times(\mathbf P^1_R)^e$ and the carriers
$g_n:Z_n=\mathbf A^N_{BK_n}\to\operatorname{Spec}k$,
where $K_n=\mu_{\ell^n}$ and $\Gamma=\mathbf Z_\ell(1)$.
Write $p_n:Z_n\to BK_n$ for the bundle projection and
$q_{m,n}:Z_m\to Z_n$ for the quotient-induced map with the same carrier.

Let $V_\bullet\in\mathcal S_{\mathrm{ad}}(\Gamma)$, with notation
as in Definition~\ref{def:frag-local-adic-source-target}.
The coefficient source is the associated-sheaf construction on these
Kummer torsors, with its induced tame descent and no additional wild
base-inertia representation.  Its finite-precision extensions and
coefficient duals are log-lisse and underlying perfect near $Y$.
The compactification boundary misses $Y$ and
\[
 T=]Y[_C=\{\,|\pi|<|x|<1\,\}\times(\mathbb D^\circ_C)^e
\]
is separated, smooth and boundaryless of dimension $N$.
These properties follow from this explicit model and construction.
Write $\mathcal E_r$ for the generic realization,
$j_\eta:X_{\bar\eta}\hookrightarrow\overline X_{\bar\eta}$,
$i:Y\hookrightarrow\overline X_s$, and $D_\ell=R\operatorname{Hom}_{\mathbf Z_\ell}(-,\mathbf Z_\ell)$.

\begin{enumerate}[label=\textup{(\roman*)}]
\item The realization
$\mathcal R_{\mathrm{ad}}:\mathcal S_{\mathrm{ad}}(\Gamma)
\to\mathcal P_{\mathrm{ad}}(T)$ is fully faithful.  Its ordinary images satisfy
\[
 Rg_{*,\mathrm{ad}}^{\mathrm{act}}V_\bullet
 \simeq N_{*,\mathrm{ad}}(V_\bullet)
 \simeq R\Gamma_{\mathrm{ad}}(T,\mathcal R_{\mathrm{ad}}V_\bullet)
 \simeq K_\Gamma(V_\bullet),
\]
where $N_{*,r}=i^*R\Psi^t_{\overline X}Rj_{\eta,*}\mathcal E_r$.
\item At each precision choose a finite presentation $W_{n_0(r),r}$
of $V_r$ as in Definition~\ref{def:frag-compatible-presenting-tails},
and inflate it for $n\ge n_0(r)$, setting $E_{n,r}=p_n^*W_{n,r}$.
The fixed-precision functor of
Lemma~\ref{lem:frag-presenting-system-assembly} is computed by
\[
 H_{!,r}(V_r)=R\varprojlim_{n\ge n_0(r),\operatorname{Tr}^!}Rg_{n,!}E_{n,r},
 \qquad
 \mathsf G_{!,\mathrm{ad}}^{\mathrm{tr}}(V_\bullet)=R\varprojlim_rH_{!,r}(V_r).
\]
The root transitions are geometric counits.  The coefficient limit
uses that lemma's natural maps $b_r$, followed by the given source
reductions.  Lemma~\ref{lem:frag-point-trace-diagram-descent} proves
their precise perfect-dual reduction square and cartesianness.
This construction is presentation-independent and satisfies
\[
 \mathsf G_{!,\mathrm{ad}}^{\mathrm{tr}}(V_\bullet)
 \simeq N^c_{\mathrm{ad}}(V_\bullet)
 \simeq R\Gamma_{c,\mathrm{ad}}(T,\mathcal R_{\mathrm{ad}}V_\bullet)
 \simeq D_\ell K_\Gamma(V_\bullet^\vee)(-N)[-2N],
\]
where $N^c_r=i^!R\Psi^t_{\overline X}Rj_{\eta,!}\mathcal E_r$.
It is not a shriek functor on the bare infinite-root fibre.
\end{enumerate}
All limits here are derived; the integral point outputs are recovered
from their cartesian finite-precision systems.  The finite operations are those of
Assumptions~\ref{ass:frag-ambient-six-functor-domain}
and~\ref{ass:frag-finite-trace-domain};
Proposition~\ref{prop:frag-point-finite-domain-verified} verifies
their use on these point carriers.  The comparisons
preserve the star-algebra module actions and the specified coefficient
reductions.  For input Tor-amplitude $[a,b]$, their perfect output
amplitudes are $[a,b+1]$ and $[a+2N-1,b+2N]$.
For the unit, $A=K_\Gamma(\mathbf Z_\ell)$ has trace partner
$D_\ell A(-N)[-2N]$ with its evaluation module structure.
Canonicity is relative to the stated model, neutralization,
compactification, carrier and enhanced coefficient categories.
\end{theorem}

The geometric comparison uses logarithmic nearby cycles and formal
completion.  The root comparison requires finite-presentation descent
on the representable-etale site and a map-level identification of
finite-stack traces.  These are proved below; the continuous bar and
Iwasawa resolutions serve as computational tools.


\paragraph{Previous work and scope.}
Logarithmic nearby cycles and formal completion are the geometric
inputs of \cite{NakayamaNearby,FujiwaraKatoNakayama,BerkovichFormalII,BerkovichEtale}.
Infinite-root constructions and the absolute Kummer comparison are
prior work \cite{TalpoVistoliInfiniteRoot,TalpoVistoliKN,ScherotzkeSibillaTalpoMcKay,CarchediEtAlLogHomotopy}.
Talpo--Vistoli already prove an equivalence between the Kummer etale
topos and the representable etale topos of the infinite root stack
\cite[Theorem 6.22, arXiv v2 numbering]{TalpoVistoliInfiniteRoot}.
The published actual-stack/pro-system and
representable-etale comparisons of
\cite[\S\S5--6]{CarchediEtAlLogHomotopy} concern profinite shapes.
Our actual-site theorem proves the finite-presentation descent needed
on the specified relative, base-root-neutralized pro-$\ell$ chart;
it also tracks the structural inverse image, hence continuous
invariants with coefficients.  It is not a new absolute infinite-root
or Kummer-shape comparison.
Brumer supplies the classical pseudocompact-module framework
\cite{BrumerPseudocompact}; the rank-one Iwasawa twist in
\cite[Appendix A.2]{EsnaultKerzPencils} is a related construction.
Our multidimensional adic action-sheaf comparison is proved separately.
The finite-stack duality input is also prior work: for a finite
$G$, \cite[Example 4.9.2]{LaszloOlssonFinite} computes compact
support on $BG$ as the dual of the full group-cochain complex.
Once compatible maps are identified, duality turning a homotopy
colimit into a homotopy limit is formal; the resulting continuous
dual formula alone is not a novelty claim.
The contribution here is assembling the specified geometric carriers
and their counits, identifying the original coefficient arrows $b_r$
by the evaluation pairing, and comparing that same system with
actual nearby costalks and analytic compact support, including module
actions and integral reduction.  Relative coefficient and Iwasawa
comparisons provide the inputs needed for this assembly.
All comparisons concern the fixed model and carrier.  Model independence,
incidence/Gysin compatibility and convolution require further results.
The root-chart functoriality proved here concerns supplied augmented
log-point maps.  Cross-stratum specialization, links and Gysin maps
require actual open/closed immersions and nearby-cycle constructions,
as in the refinement paper; they do not follow from relative lattice
maps.  Supported realization over general bases, rank-changing
comparisons and descent of integral models are separate extensions.
They are not conclusions of the fixed-node theorem.
\label{frag:literature-status}

The comparison maps use three additional ingredients whose proofs
are included here.  Proposition~\ref{prop:frag-finite-quotient-enhancement}
identifies the enhanced finite-stack adjunctions and their scalar
mates; Lemma~\ref{lem:frag-point-trace-diagram-descent} then identifies
the actual continuous trace reduction.  Proposition~
\ref{prop:frag-completion-comparison-monoidal} constructs the
multiplicative completion map before invoking the external
cohomological comparison theorem.  For the general-base coefficient
model, Lemmas~\ref{lem:frag-filtered-enhanced-localization}
and~\ref{lem:frag-etale-sections-torsion} supply the categorical
and sheaf-theoretic inputs.  The appendices also treat supplied finite
semilinear descent diagrams; the node theorem has no such descent
group.  Global weak base-inertia extensions are not claimed here.

We first describe the coefficients and the root structural functor,
then construct the analytic and nearby realizations.  Verdier duality
and the finite-root trace calculation give the supported comparison.
The proof of the theorem and examples follow.  Appendices contain
the general-base descent arguments, standard comparison machinery,
and two obstructions to larger coefficient sources.  Extensions to
finite symmetries and proper images are recorded separately.


\paragraph{Dependency map for the node theorem.}
The arrows below denote dependencies of proofs.  At finite precision:
\[
\begin{array}{c@{\qquad}c}
\text{root site; log specialization}&\text{finite quotient adjunctions}\\
\text{completion; analytic sites}&\Downarrow\\
\Downarrow&\text{full bars; geometric traces}\\
\text{star and full faithfulness}&\Downarrow\\
\Downarrow\ \text{Verdier duality}&\text{continuous dual; reduction square}\\
\text{compact support; costalk}&\text{finite-carrier trace limit}
\end{array}
\]
The continuous dual and reduction-square step on the right uses
star perfectness from the left.  Their pairing identifies the two
bottom outputs.  Integral recovery then applies to the resulting
compatible coefficient systems on both sides.
The left branch uses
Lemma~\ref{lem:frag-log-limit-slices},
Proposition~\ref{prop:frag-log-specialization-kummer-localization},
Corollary~\ref{cor:frag-point-actual-root-bridge}, and
Lemma~\ref{lem:frag-derived-analytic-site-model}.
The right branch uses
Proposition~\ref{prop:frag-finite-quotient-enhancement},
Lemma~\ref{lem:frag-point-trace-diagram-descent}, and
Theorem~\ref{thm:frag-point-C1}.
Point-source recovery is
Proposition~\ref{prop:frag-point-source-iwasawa-admissibility} and
Lemma~\ref{lem:frag-point-perfect-recovery}.
The general-base appendices supply these inputs by specialization;
their finite-action and torsion lemmas remain part of the proof.
General-base sheaf globalization, finite semilinear descent, and the
classical pseudocompact Hom interpretation are additional conclusions,
not further hypotheses of the node theorem.  The unipotent refinement
interface is a subsequent corollary.  None of these extensions supplies
model independence or a global supported functor.

\section{The node and its Kummer coefficients}
\label{sec:local-node-coefficients}

Fix the data of Theorem~\ref{thm:frag-node-complete-local-datum}.
The chart is $\mathbf N\to P=\mathbf N^2$, $1\mapsto\rho=(1,1)$.
The coordinate $x$ identifies $P^{\mathrm{gp}}/\mathbf Z\rho$ with
$\mathbf Z$.  Fixing compatible roots of $\pi$ leaves the relative
groups $K_n=\mu_{\ell^n}$ and $\Gamma=\varprojlim_nK_n$.
The cover $a^{\ell^n}=x$ on the generic fibre is a $K_n$-torsor;
its associated-sheaf functor constructs the coefficients used below.

A point of the generic fibre specializes to $Y$ precisely when
$|x|,|y|,|z_i|<1$.  Thus
\[
 T=\{\,|\pi|<|x|<1\,\}\times(\mathbb D^\circ)^e.
\]
We compactify by $X_0Y_0=\pi X_1Y_1$ in
$(\mathbf P^1_R)^2\times(\mathbf P^1_R)^e$; the boundary misses $Y$.
The finite root carriers are $\mathbf A^N_{BK_n}$, with $N=1+e$.
Their vector-bundle trace fixes the factor $(-N)[-2N]$.

Throughout, $n$ is root level (root order $\ell^n$) and $r$ is
coefficient precision.
We take the continuous root (co)limit at fixed $r$, recover the
integral coefficient system, and finally invert $\ell$ if desired.
Write $V^\vee=R\operatorname{Hom}_{\mathbf Z_\ell}(V,\mathbf Z_\ell)$
with contragredient action.  This coefficient dual differs from the
geometric source dual $V^\vee(N)[2N]$.

\begin{definition}[Integral Kummer coefficients]
\label{def:frag-adic-coefficients}
Put $R_r=\mathbf Z/\ell^r$.  On a coefficient topos $X$ let
$\mathcal D(X,R_r)$ be the stable localization of module-sheaf
complexes at quasi-isomorphisms.  Define
\begin{equation}
 \mathcal D_{\mathrm{ad}}(X,\mathbf Z_\ell)
 =\varprojlim_{r,\,-\otimes^L_{R_{r+1}}R_r}\mathcal D(X,R_r),
 \qquad \mathcal D^b_{c,\mathrm{ad}}(X,\mathbf Z_\ell)
 \subset\mathcal D_{\mathrm{ad}}(X,\mathbf Z_\ell).
 \label{eq:frag-adic-infinity-category}
\end{equation}
The full subcategory consists of derived cartesian systems with
uniform bounded amplitude and the stated constructible perfect
restrictions on a finite stratification.  Mapping spectra are those
of this limit category.  Reduction is taken in the ambient category:
$\mathbf F_\ell\otimes^L_{\mathbf Z/\ell^2}\mathbf F_\ell$ shows why
it is not a functor on all bounded constructible complexes.
Rational coefficients mean inversion of $\ell$ in this integral category.

On finite Artin stacks we use the enhanced six operations
\cite{LaszloOlssonFinite,LaszloOlsson,LiuZhengEnhanced,LiuZhengAdic}.
Individual star and shriek images lie in their respective $+$ and
$-$ constructible domains.  All root (co)limits are formed in the
ambient unbounded category: finite classifying-stack cohomology is
unbounded \cite[Example 4.9.2]{LaszloOlssonFinite}.  Boundedness of the
continuous outputs is proved by the Iwasawa calculation below.

For $\Gamma_M=M\otimes\mathbf Z_\ell(1)$ put
$\Lambda_M^{\mathrm{Iw}}=\mathbf Z_\ell[[\Gamma_M]]$ and
$\Lambda_{M,r}^{\mathrm{Iw}}=\Lambda_M^{\mathrm{Iw}}/\ell^r$; in particular
\begin{equation}
 \Lambda_{M,r}^{\mathrm{Iw}}
 \ne R_r[\Gamma_M/\ell^r\Gamma_M].
 \label{eq:frag-iwasawa-not-finite-group-ring}
\end{equation}
At a point, $\mathsf{Mod}^{\mathrm{pc}}_{\Lambda_M^{\mathrm{Iw}}}$ denotes
pseudocompact modules, derived from inverse systems of finite-length
modules for $(\ell^r,I_M^a)$, with completed tensor and continuous Hom.
An \emph{Iwasawa-admissible coefficient} is a bounded, locally
pseudo-coherent object of finite Tor-amplitude over the Iwasawa
algebra whose underlying $\mathbf Z_\ell$-complex is perfect.
This is the classical point terminology.  Over a stratified base,
``selected'' or ``admissible'' means the uniformly underlying-perfect
cartesian class in Definition~\ref{def:frag-adic-iwasawa-sheaves},
without an implicit global pseudocompact sheaf condition.  Torsion cohomology
is allowed; the regular Iwasawa module is excluded when $M\ne0$.

For sheaf coefficients the independent adic module model is
Definition~\ref{def:frag-adic-iwasawa-sheaves}.
Theorem~\ref{thm:frag-F-adic-iwasawa} identifies it with the selected
continuous-action systems.  A separate global pseudocompact-sheaf
interpretation is not used.  At precision $r$, finite-action descent
supplies a root level $n(r)$, including derived maps; these levels
need not be bounded as $r$ varies.  Ambient coefficient recovery is
Lemma~\ref{lem:frag-adic-completion}, and perfect recovery at a point
is Lemma~\ref{lem:frag-point-perfect-recovery}.

We write $V=V_{\mathrm{cont}}[1/\ell]$ and
$\Lambda_{M,\mathbf Q_\ell}^{\mathrm{Iw}}=\Lambda_M^{\mathrm{Iw}}[1/\ell]$.
This is the bounded-denominator algebra, rather than
$\mathbf Q_\ell[[T_1,\ldots,T_d]]$; coefficients throughout admit
the specified perfect integral model.
\end{definition}


\begin{definition}[The source and analytic target of adic realization]
\label{def:frag-local-adic-source-target}
Let $\mathcal A_{\Gamma,r}$ be discrete continuous $\Gamma$-modules
with coefficients $R_r$, and set
\[
 \mathcal S_r=\{V\in D(\mathcal A_{\Gamma,r}):
                    V\text{ is perfect over }R_r\},\qquad
 \mathcal P_r(T)=\operatorname{Perf}_{\mathrm{lisse}}(T_{\acute et},R_r).
\]
The latter is the full subcategory of module-sheaf complexes locally
represented by bounded finite free complexes, with finite locally
constant cohomology.  Both are full subcategories of the stable
localizations at quasi-isomorphisms.  Scalar reduction is the
\emph{derived} functor $-\otimes^L_{R_{r+1}}R_r$ in the ambient categories.
Write $\mathcal S_{\mathrm{ad}}(\Gamma)$ and
$\mathcal P_{\mathrm{ad}}(T)$ for the full subcategories of their
ambient category limits consisting of cartesian systems in the
displayed classes with one common finite Tor interval.  This
includes \emph{all} such systems.  Proposition~\ref{prop:frag-point-source-iwasawa-admissibility}
proves their classical Iwasawa-admissibility at a point.  No common
finite root level or trivializing cover across all $r$ is required.
The point realization of the underlying system of $V_\bullet$ is
$V=R\varprojlim_rV_r$; its perfectness is a conclusion of
Lemma~\ref{lem:frag-point-perfect-recovery}.

Tensor, the coefficient dual, and the internal coefficient Hom are
formed levelwise: $V_r^\vee=R\operatorname{Hom}_{R_r}(V_r,R_r)$ with
contragredient action, and
$\underline{\operatorname{Hom}}(V_r,W_r)=V_r^\vee\otimes^L_{R_r}W_r$.
These are distinct from the global equivariant mapping complex.
The same conventions apply on $T$.  Completed tensor means the
resulting cartesian system; no ordinary integral sheaf is substituted.
Set
\[
 K_\Gamma(V_\bullet)=R\varprojlim_r C^*_{\mathrm{cts}}(\Gamma,V_r).
\]
We also write this as
$R\operatorname{Hom}^{\mathrm{cts}}_{\mathbf Z_\ell[[\Gamma]]}
(\mathbf Z_\ell,V_\bullet)$, computed by the finite augmentation
Koszul complex in the adic module model of
Theorem~\ref{thm:frag-F-adic-iwasawa}.  Proposition~\ref{prop:frag-point-source-iwasawa-admissibility} constructs
classical pseudocompact models for the entire point source, and
Proposition~\ref{prop:frag-point-pc-Hom} identifies the resulting
calculation with classical continuous derived Hom.  The definition
of the source and its mapping spectra is the cartesian-system one.
\end{definition}

\section{Root coefficients and continuous cohomology}
\label{sec:local-root-cohomology}

On the chosen logarithmic fibre, finite root reductions are $BK_n$.
Finite nilpotent thickening does not change their etale coefficient
topoi.  The following descent results identify the actual root
structural functor, including its maps, with continuous invariants.

\begin{proposition}[Finite derived action descent over a point]
\label{prop:frag-point-finite-action-descent}
Fix $r$.  Let $\mathcal D_{n,r}^{b,\mathrm{fin}}$ be the full subcategory
of $D(R_r[K_n]\text{-}\mathrm{Mod})$ with bounded finite underlying
cohomology, and let $\mathcal D_{\Gamma,r}^{b,\mathrm{fin}}$ be the analogous
subcategory of the derived category of discrete continuous $\Gamma$-modules.
Inflation induces an equivalence of stable categories
\[
\operatorname*{colim}_n\mathcal D_{n,r}^{b,\mathrm{fin}}
 \simeq\mathcal D_{\Gamma,r}^{b,\mathrm{fin}}.
\]
It restricts to the full subcategories whose underlying $R_r$-complex is
perfect.  A finite simplicial diagram, including its specified higher
simplices, descends to a common level.  This does not assert descent of
an entire infinite diagram at one level.
\end{proposition}

\begin{proof}
First consider finite modules $M,N$ with action through $K_{n_0}$.
Resolve $M$ by a bounded-above complex $P$ of finitely generated free
$R_r[K_{n_0}]$-modules.  On forgetting the action this is a resolution
by finite free $R_r$-modules.  The equivariant bar construction applied
to this resolution computes derived equivariant Hom: its double complex
has terms
\[
 C^a\bigl(K_n,\operatorname{Hom}_{R_r}(P^{-b},N)\bigr),
 \qquad a,b\ge0.
\]
Use continuous cochains for $\Gamma$.  This computation follows by taking
the bar resolution relative to the forgetful functor to $R_r$-modules;
the free underlying resolution makes each vertical Hom derived-correct.
Each diagonal $a+b=t$ is finite.  Every continuous map from the compact
space $\Gamma^a$ to the finite coefficient module factors through some
$K_n^a$: its finitely many clopen fibres have a common open normal
subgroup on which they are invariant.  Consequently the filtered colimit
of the finite double complexes is the continuous double complex, with
all differentials, and exactness of that colimit gives
\[
\operatorname*{colim}_{n\ge n_0}
 R\operatorname{Hom}_{R_r[K_n]}(M,N)
 \simeq R\operatorname{Hom}_{\Gamma,\mathrm{disc}}(M,N).
\]
There is no finite-group truncation in this argument.

Finite Postnikov filtrations in both variables extend the identity to
bounded objects with finite cohomology.  In particular it identifies
mapping spectra, not just degree-zero maps.  Every finite continuous
module has action through a finite quotient.  Induct on the number of
nonzero cohomology degrees of an object on the continuous side.  Present
its lower truncation and its highest cohomology module at finite levels;
the mapping identity presents the Postnikov attaching morphism at a
common larger level.  Its fibre presents the object.  This proves essential
surjectivity without asserting that the cohomology modules of a perfect
complex are perfect.  Inflation leaves the underlying $R_r$-complex
unchanged, so the equivalence restricts to the underlying-perfect class.
For the last assertion use the filtered-quasicategory model and
finite-diagram argument in
Lemma~\ref{lem:frag-filtered-enhanced-localization}, applied to the
category equivalence just proved.  It descends the finitely specified
simplices themselves, not only the actions on cohomology.
\end{proof}

\begin{corollary}[The actual root structural functor over a geometric point]
\label{cor:frag-point-actual-root-bridge}
Over an algebraically closed point, the actual relative root site has
the generating site of finite continuous $\Gamma$-sets.  Its sheaf
topos is the topos of discrete continuous $\Gamma$-sets, structural
pullback is trivial action, and actual derived star on the selected
coefficients is $C^*_{\mathrm{cts}}(\Gamma,-)$.
\end{corollary}
\begin{proof}
For the node, at order $a=\ell^n$ the rigidified atlas is explicitly
$\operatorname{Spec}k[u,v]/(u^a,v^a,uv)$ with $\mu_a$-weights
$(1,-1)$.  Its monomial reduction is the point, and refinements
send $u_n,v_n$ to the corresponding powers of $u_m,v_m$.
Lemma~\ref{lem:frag-relative-atlas-torsors} verifies the nerve and
test-torsor properties used in the limit argument (and gives the
formula for the other split log-point charts).
Theorem~\ref{thm:frag-actual-root-fp-descent} descends etale fp objects,
maps and coverings, allowing fpqc-stack sources.  At a finite level,
nil reduction gives $BK_n$.  An etale fp algebraic space over an
algebraically closed field is a finite set of points.  Descent along
the torsor atlas equips it with a $K_n$-action; covering families
are jointly surjective maps.  The filtered site is therefore finite
continuous $\Gamma$-sets.  The generating-site argument, with the
smoothification just specified, identifies its sheaves with smooth
sets.  Structural pullback is trivial action on finite base sets,
which generate the base topos, hence on all base sheaves.
The resulting exact equivalence of module-sheaf abelian categories
identifies injectives and the right adjoint of structural pullback.
Deriving in their complex localizations gives continuous invariants
with its actual unit.  Point finite-action descent and the finite
augmentation resolution prove selected boundedness and reduction.
This point argument uses neither general-base constructible descent
nor $\Delta$-rectification.
\end{proof}


Put $\Lambda=\mathbf Z_\ell[[\Gamma]]$.  A generator $\gamma$ gives
$\Lambda=\mathbf Z_\ell[[\gamma-1]]$ and
\[
 K_\Gamma(V)=R\operatorname{Hom}^{\mathrm{cts}}_\Lambda(\mathbf Z_\ell,V)
 \simeq [\,V\xrightarrow{\gamma-1}V\,].
\]
The two displayed copies have relative degrees $0,1$.  In rank $d$
the corresponding complex is the Koszul complex of
$\gamma_1-1,\ldots,\gamma_d-1$.  Naturality is expressed by derived
continuous cohomology; the coordinates only supply a finite model.

\begin{lemma}[Continuous bar--Iwasawa comparison]
\label{lem:frag-continuous-bar-iwasawa}
There is an augmentation-preserving equivalence
\[
 \mathbf K_M^{\mathrm{Iw}}\xrightarrow{\sim}\mathbf B_M^{\mathrm{cts}}
\]
in the derived category of complete $\Lambda_M^{\mathrm{Iw}}$-modules.  For an
Iwasawa-admissible coefficient $V_{\mathrm{cont}}$, both sides compute
continuous group cohomology:
\begin{equation}
 R\Gamma_{\mathrm{cts}}(\Gamma_M,V_{\mathrm{cont}})
 \simeq
 R\!\operatorname{Hom}^{\mathrm{cts}}_{\Lambda_M^{\mathrm{Iw}}}
      (\mathbf K_M^{\mathrm{Iw}},V_{\mathrm{cont}}).
 \label{eq:frag-continuous-cohomology}
\end{equation}
The comparison is natural in continuous lattice maps and finite groups of
lattice automorphisms.  The cup product is the canonical multiplication on the
derived endomorphism object of the augmentation module; it does not depend on
a strict choice of comparison map.
\end{lemma}

\begin{proof}
After the chosen basis,
$\Lambda_M^{\mathrm{Iw}}=\mathbf Z_\ell[[T_1,\ldots,T_d]]$ and the augmentation
module is the quotient by the regular sequence $(T_1,\ldots,T_d)$.  Hence
\eqref{eq:frag-iwasawa-koszul} is a finite free resolution.  The completed bar
complex is a topologically free resolution of the same augmentation module.
The comparison theorem for projective resolutions gives the displayed
equivalence, uniquely up to contractible homotopy in the enhanced derived
category.  Applying continuous internal Hom gives
\eqref{eq:frag-continuous-cohomology}.  Naturality and equivariance follow by
using the basis-free derived object
$R\!\operatorname{Hom}^{\mathrm{cts}}_{\Lambda_M^{\mathrm{Iw}}}
(\mathbf Z_\ell,-)$; the chosen Koszul complex is only a finite model for it.
\end{proof}


\section{Analytic realization and nearby cycles}
\label{sec:local-analytic-realization}

The Kummer torsors define the generic coefficient before any
cohomology comparison is made.  We prove the comparison for a split
vertical log-smooth chart, then specialize to the node.  The
specialization functor is
$R\Theta=R\nu_*\mu^*$ for
$T_{\acute et}\xleftarrow{\mu}T_{q\acute et}\xrightarrow{\nu}Y_{\acute et}$.
The completion and site comparisons used here are recorded in
Appendix~\ref{sec:frag-analytic-nearby}.

\begin{definition}[The vertical logarithmic coefficient domain]
\label{def:frag-vertical-log-B-domain}
Let $X/R$ be a finite-type fs saturated log-smooth scheme, vertical near
a noetherian locally closed stratum $Y$: the generic log structure is
trivial on that neighbourhood.  On $Y$ fix a split characteristic chart
$\mathbf N\rho\to P$ with $\rho$ primitive and
$L=P^{\mathrm{gp}}/\mathbf Z\rho$ free of rank $d$.  Use the actual
special formal completion $\widehat X_{/Y}$ and its geometric tube,
without assuming an entire-torus product.  For a supplied generic open
$j:U\hookrightarrow X_\eta$, require
$Y\cap\overline{X_\eta\setminus U}=\varnothing$.

A coefficient for the theorem below is a bounded complex $\mathcal L_r$
on $X_{\mathrm{ket}}$, in a neighbourhood of $Y$, with finite locally
constant cohomology and the selected underlying $R_r$-perfectness.  Its
generic restriction, on the ordinary etale site by verticality, is
$\mathcal E_r$.  Its restriction $E_r$ to the geometric relative log
fibre is required to be trivial on the non-$\ell$ relative tame factor.
For adic coefficients require a derived-cartesian system with the
uniform amplitude conditions already specified in Theorem~\ref{thm:frag-F-adic-iwasawa}.
If arithmetic descent is retained, it is the descent supplied by this
actual prime-to-$p$ log-Kummer source and the chosen geometric base-root
lift.  We allow no additional wild base-inertia representation on the
coefficient.  This restriction is imposed on the logarithmic
neighbourhood, not merely on the restriction to $Y$.  The non-$\ell$
relative tame factor is discarded by exact invariants; external tame
base inertia, including its prime-to-$\ell$ part, remains separate
action data and is not absorbed into $\Gamma$.  The same convention
applies to dual coefficients and to log-lisse sources used below.
The extension $\mathcal L_r$ is actual source data, not an extension
asserted to exist for every abstract constructible action coefficient.
The following statements take values on $Y_{\acute et}$ itself;
application of a further $Rf_{Y,*}$ is a separate operation on both sides.
\end{definition}



\begin{lemma}[Slices of the geometric logarithmic limit sites]
\label{lem:frag-log-limit-slices}
Let $X$ be an fs log scheme over the logarithmic trait.  Use the
finite logarithmic base extensions and the limiting Kummer etale
sites defining the geometric total, special and generic topoi
in \cite[\S8.2, (8.2.1)]{FujiwaraKatoNakayama}.
If $q:X'\to X$ is Kummer etale, its base changes determine objects
$Q_\alpha$ of these three topoi $\mathcal X_\alpha$, and canonically
\[
 \mathcal X'_\alpha\simeq\mathcal X_\alpha/Q_\alpha,
 \qquad \alpha\in\{\mathrm{tot},\bar s,\bar\eta\}.
\]
Under these equivalences the generic and special structural maps
for $X'$ are the localizations of those for $X$.
The identifications commute with further base extension, composition
of Kummer etale maps, and transport by geometric base automorphisms.
\end{lemma}
\begin{proof}
Work on quasi-compact chart opens and use the finite-presentation
basis of the Kummer etale sites.  A general Kummer etale object is
covered by these opens, so the resulting statement glues.
Let $\lambda$ index the finite logarithmic base extensions.  At each
level use the fs fibre products and denote the corresponding site
by $\mathcal C_{\lambda,\alpha}$ and the base change of $q$ by
$Q_{\lambda,\alpha}$.  Subcanonicity and stability of Kummer etale
maps under composition and fs base change identify the site over
$X'_\lambda$ with
$\mathcal C_{\lambda,\alpha}/Q_{\lambda,\alpha}$.
Its coverings are precisely the induced Kummer etale coverings.

Here is the passage to the limit explicitly.  The limiting site is
presented by the filtered category $\mathcal C_{\infty,\alpha}$
of finite-level objects, with morphisms taken after further base
extension and the topology generated by finite-level coverings.
This is the limiting-site convention in the cited definition, not
the small site of an unqualified underlying scheme limit.
There is a functor
\[
 \operatorname*{colim}_\lambda
 (\mathcal C_{\lambda,\alpha}/Q_{\lambda,\alpha})
 \longrightarrow \mathcal C_{\infty,\alpha}/Q_\alpha.
\]
An object on the right is a finite-level object $U_\lambda$ together
with an arrow to $Q_\alpha$.  That arrow is represented at a later
level, so it gives an object on the left.  A morphism over $Q_\alpha$
is a finite-level morphism and one triangle equality; both are
represented at a common later level.  The same observation detects
equality of two morphisms.  This proves equivalence of the categories.
For a finite covering family all arrows and their triangles can be
represented at one level.  The induced covering topology is the
same on both sides, since both are generated by pullbacks of the
finite-level coverings.  On the quasi-compact basis finite coverings
suffice; arbitrary opens and coverings are obtained by gluing.
Sheafification identifies the topos of this slice site with the
slice by its representable sheaf, proving the assertion.

The generic and special maps are induced at every level by the
same fs pullbacks.  Pulling back an arrow $U\to Q$ gives the arrow
between the corresponding pullbacks, with its triangle unchanged.
Thus the three site equivalences intertwine those structural maps,
not just the underlying categories.  Successive pullbacks supply
the composition identifications; their associativity follows from
the universal property of fibre product.  A base automorphism
transports these same diagrams (and reindexes the finite extensions),
which gives the asserted geometric-action compatibility.  No formula
interchanging a derived direct image and a filtered colimit is used.
\end{proof}

\begin{proposition}[Kummer localization of logarithmic specialization]
\label{prop:frag-log-specialization-kummer-localization}
Let $X$ be an fs log scheme over the logarithmic trait and
$q:X'\to X$ a Kummer etale morphism.  Use the geometric logarithmic
base change defining nearby cycles, with the integral saturated fibre
products and limiting Kummer etale topoi of
\cite[\S8.2, (8.2.1)]{FujiwaraKatoNakayama}.
For $\mathcal L\in D^+(X_{\mathrm{ket}},R_r)$ there is a canonical
exchange
\[
 q_{\bar s}^*R\Psi_X^{\log}\mathcal L
 \xrightarrow{\sim}
 R\Psi_{X'}^{\log}(q^*\mathcal L).
\]
It makes the specialization-unit square commute:
\[
\begin{array}{ccc}
 q_{\bar s}^*(\mathcal L|_{X_{\bar s}^{\log}})
 &\xrightarrow{\ q_{\bar s}^*\mathrm{sp}_{\mathcal L}\ }&
 q_{\bar s}^*R\Psi_X^{\log}\mathcal L\\
 \big\Vert&&\big\downarrow\sim\\
(q^*\mathcal L)|_{(X')_{\bar s}^{\log}}
 &\xrightarrow{\ \mathrm{sp}_{q^*\mathcal L}\ }&
 R\Psi_{X'}^{\log}(q^*\mathcal L).
\end{array}
\]
For a covering family of such $q$, the pullbacks on the geometric
logarithmic special fibre are jointly conservative.  The exchanges
are compatible with morphisms of coefficients, composition of covers,
and the retained inertia action.
\end{proposition}

\begin{proof}
The Kummer etale topology is subcanonical; composition and fs base
change preserve Kummer etale maps and their coverings
\cite[\S2.1--2.3 and Theorem 2.6]{FujiwaraKatoNakayama}.
Consequently $(X')_{\mathrm{ket}}$ is the slice of
$X_{\mathrm{ket}}$ by the representable object $X'$: its objects are
Kummer etale objects over $X$ with a map to $X'$, and its coverings
are the induced coverings.  This is an equality of the localized
sites, not an ordinary-scheme etale base-change assertion.

Let $\mathcal X$, $\mathcal X_{\bar s}$ and $\mathcal X_{\bar\eta}$
be the logarithmic total, special and generic topoi in the geometric
base change, and write $i,j$ for their structural morphisms.
The defining expression is $R\Psi_X^{\log}\mathcal L
=i^*Rj_*(\mathcal L|_{\mathcal X_{\bar\eta}})$.
Lemma~\ref{lem:frag-log-limit-slices} identifies the three topoi
for $X'$ with the slices by the compatible base changes of $q$,
including their generic and special structural maps.

Apply Lemma~\ref{lem:frag-slice-derived-base-change} to $j$ and this
object of $\mathcal X$.  Commutation of inverse images in the
special-fibre square then gives
\[
 q_{\bar s}^*i^*Rj_*
 \simeq i'^*q_{\mathcal X}^*Rj_*
 \simeq i'^*Rj'_*q_{\bar\eta}^*.
\]
The specialization map is $i^*$ applied to the unit
$\mathcal L|_{\mathcal X}\to Rj_*j^*(\mathcal L|_{\mathcal X})$.
The unit clause of the lemma therefore gives the displayed square
with this very map.  Coverings remain coverings under the geometric
log base change, so exact inverse images along them are jointly
conservative on cohomology sheaves.  This proves the descent assertion.
Composition and inertia compatibility follow from the same site
pullbacks and adjunction units before forgetting the action.
\end{proof}

\begin{lemma}[Logarithmic specialization for the actual lisse source]
\label{lem:frag-log-nearby-geometric-input}
For an fs log-smooth scheme $X$ over the logarithmic trait and a
bounded $R_r$-complex $\mathcal L$ on $X_{\mathrm{ket}}$ with finite
locally constant cohomology, the log-specialization unit
\[
 \mathcal L|_{X_{\bar s}^{\log}}
 \xrightarrow{\sim}R\Psi_X^{\log}\mathcal L
\]
is an equivalence.  In addition there is a natural exchange
\[
 R\Psi_X(R\varepsilon_{\eta,*}\mathcal L_\eta)
 \simeq R\varepsilon_{s,\mathrm{rel},*}R\Psi_X^{\log}\mathcal L.
\]
These use the geometric logarithmic base change, not the special
fibre with its absolute log structure.  In the vertical case the
generic logarithmic forgetful map $\varepsilon_\eta$ is the identity.
The maps are the specialization unit and the forgetful exchange on
the indicated actual topoi.
\end{lemma}

\begin{proof}
The external local-acyclicity input is the constant-coefficient
statement of \cite[Theorem 3.2]{NakayamaNearby}, also stated in
\cite[Theorem 8.3]{FujiwaraKatoNakayama}.  It applies to
$\mathbf Z/\ell^s$ for every $s$.  A finite constant $R_r$-module
is a finite direct sum of such cyclic modules, so its unit is an
equivalence as well.  This is the only coefficient case taken
from that theorem.

There are three separate steps for the lisse extension of this input.
First choose a Kummer-etale covering family trivializing a finite
locally constant coefficient.  Finitely many such coefficients or
specified maps can be handled on the common refinement obtained by
finite fibre products of covering families; no single cover for an
infinite coefficient system is used.  Second, each covering fs log
scheme remains log smooth over the trait, since Kummer etale is log
etale and log smoothness is stable under composition.  Saturation
over the trait is not an extra hypothesis of the cited theorem.
Third, Proposition~\ref{prop:frag-log-specialization-kummer-localization}
identifies the pullback of the actual specialization unit with the
unit on that covering scheme.  The base-changed Kummer-etale covering
of the geometric logarithmic special fibre is jointly conservative,
so the unit on the original coefficient is invertible.
This third step and the source descent, rather than a stronger
coefficient version of the external theorem, provide the extension.

Finite truncation triangles apply this natural unit to a bounded
complex with locally constant cohomology.  No splitting of its
Postnikov extensions or simultaneous trivialization of the whole
complex is required.

The forgetful exchange is the derived natural transformation of
\cite[\S8.2, (8.2.3)]{FujiwaraKatoNakayama}, which is stated there
for $D^+$ coefficients, rather than only the constant coefficient.
Its construction uses strict closed base change for the forgetful
map and the geometric logarithmic base limit.  Apply that exchange
to the present source.  Its unit and pullback naturality are retained
under the localizations above.  Thus neither the lisse specialization
nor its map-level compatibility has been added as a hypothesis.
\end{proof}

\begin{theorem}[Tube cohomology of vertical log-lisse Kummer coefficients]
\label{thm:frag-vertical-log-star}
For the actual prime-to-$p$ log-Kummer source of
Definition~\ref{def:frag-vertical-log-B-domain}, with no additional
wild base-inertia action and with the full retained inertia action, write
$s:\,]Y[_{\bar\eta}\to Y$ for geometric specialization and
$\Gamma=\operatorname{Hom}(L,\mathbf Z_\ell(1))$.  There is a natural
zigzag of equivalences
\[
 C^*_{\mathrm{cts}}(\Gamma,E_r)
 \xrightarrow{\ \gamma\ }
 i_Y^*R\Psi_X(Rj_*\mathcal E_r)
 \xrightarrow{\ \beta\ }
 Rs_*\mathcal E_r^{\mathrm{an}},
 \qquad
 i_Y^*R\Psi_X^t(Rj_*\mathcal E_r)
 \xrightarrow{\sim}i_Y^*R\Psi_X(Rj_*\mathcal E_r).
\]
Here $E_r$ is interpreted on the actual root fibre by
Proposition~\ref{prop:frag-relative-log-root-interface}.
The first zigzag respects the commutative algebra of the unit and its
coefficient modules, and the tame comparison preserves these structures.
The maps are natural for algebraized Kummer refinements, strict-etale
restrictions, and supplied finite semilinear descent.  They identify
actual functors on the indicated coefficient source.  Their selected
adic outputs are uniformly bounded and derived cartesian.

This closes the local star comparison for this vertical lisse domain.
It does not assert essential surjectivity from every abstract
cross-stratum action coefficient to $\mathcal L_r$, and supplies no
exceptional or incidence comparison.
\end{theorem}

\begin{proof}
Remove the closure of the generic complement near $Y$, as in Proposition~\ref{prop:frag-ordinary-nearby-boundary-disjoint}.
The source there is the actual generic restriction of $\mathcal L_r$.
Let $E_r$ be its restriction to the relative geometric log fibre.
Apply $R\varepsilon_{Y,\mathrm{rel},*}$ to the specialization unit
\[
 E_r\longrightarrow (R\Psi^{\log}\mathcal L_r)|_{\widetilde Y}.
\]
Lemma~\ref{lem:frag-log-nearby-geometric-input} makes this an
equivalence and identifies its target with the restricted ordinary
nearby complex.  Proposition~\ref{prop:frag-relative-log-root-interface}
identifies its source with continuous invariants and with the actual
root structural star.  This constructs $\gamma$ on the entire face
etale topos, including its maps; it is not a reconstruction from equal
stalk dimensions.  Proposition~\ref{prop:frag-ordinary-nearby-boundary-disjoint}
constructs the actual completion comparison $\beta$ and proves it is
an equivalence.  The rightmost term is therefore cohomology on the
actual tube, even when that tube is an annulus or a higher toric domain.

For tameness, the geometric logarithmic special fibre is constructed
using prime-to-$p$ roots of the base parameter.  Its descent action,
and that of the restriction of the given $X_{\mathrm{ket}}$ coefficient,
factor through tame base inertia.  Consequently wild inertia acts
trivially on the cohomology of
$R\varepsilon_{Y,\mathrm{rel},*}E_r$.  The ordinary comparison just
constructed is equivariant for the full retained inertia action:
the specialization unit and logarithmic forgetful exchange are formed
before forgetting that action and are natural for base automorphisms.
Thus it proves this assertion for the
actual nearby complex.  Lemma~\ref{lem:frag-wild-inertia-criterion}
now makes its tame-invariant inclusion an equivalence.  This argument
uses the log extension and verticality, rather than merely the name
of the relative deck group.

The specialization unit is the unit of a geometric pullback/right-adjoint
pair.  The forgetful exchange in the displayed log/ordinary comparison
is the mate of the strict pullback comparison of those same sites.
Both are morphisms of the induced lax symmetric-monoidal functors:
their adjoints are composites of strong tensor inverse images and the
monoidal counits.  Their invertibility on the present source was proved
before taking the inverse exchange used in $\gamma$.  Proposition~
\ref{prop:frag-completion-comparison-monoidal} gives precisely this
structure for $\beta$.  Invariants have the analogous lax structure
and their inclusion is monoidal.  These constructions prove the
algebra/module claims for the specified zigzag, including higher
multiplication cells, without positing a new common operadic site.
They also give the stated naturality under the geometric maps and
coefficient morphisms.

Finally the continuous augmentation resolution has $d+1$ finite free
terms.  On an input of amplitude $[a,b]$ it gives output amplitude
contained in $[a,b+d]$, independent of $r$.  The construction commutes
with derived coefficient reduction on the selected class.  The natural
comparison zigzag transports this property to the actual analytic and
nearby outputs.  Only after this recovery do we take the inverse limit
in $r$.  No finite group bar has been replaced by a bounded complex
before continuous passage.  A subsequent direct image from $Y$ to another
base requires its own boundedness/preservation assertion.
\end{proof}

\begin{corollary}[Actual tube cohomology and geometric refinement weights]
\label{cor:frag-log-tube-root-weights}
For the unit coefficient in Theorem~\ref{thm:frag-vertical-log-star},
\[
 R^j s_*R_r\simeq\bigwedge^j(L\otimes R_r)(-j),\qquad 0\le j\le d,
\]
with zero in other degrees.  Suppose the supplied compatible fs Kummer
models over chosen tame base roots induce multiplication by
$\ell^{m-n}$ on their relative characteristic lattices.  Their actual
tube pullbacks act on these sheaves by $\ell^{j(m-n)}$.  These are
geometric refinement maps between the covers, not an identification of
a finite-cover direct image with a finite group bar.
\end{corollary}

\begin{proof}
The theorem identifies the actual derived output with continuous
invariants by a natural geometric map.  For the trivial coefficient,
its degree-one cohomology is
$\operatorname{Hom}_{\mathrm{cts}}(\Gamma,R_r)=L\otimes R_r(-1)$.
Cup products compute its exterior powers by the augmentation Koszul
resolution.  The character pairing in the finite Kummer charts sends
a root refinement to restriction along multiplication by
$\ell^{m-n}$ on relative inertia, hence to the stated degree-one
weight and its exterior powers.  Naturality of $\gamma$ and $\beta$
identifies these maps with actual tube pullbacks.  The bounded degree
range and eventual vanishing of positive weights at fixed $r$ are
therefore now geometric assertions in this domain.  They can be used
in the actual finite-torsor filtration of Proposition~\ref{prop:frag-actual-finite-galois-descent}.
\end{proof}

\begin{corollary}[Kummer coefficients on a semistable node]
\label{cor:frag-node-full-Kummer-star}
Let the residue field be algebraically closed, put
\[
 X=\operatorname{Spec}R[x,y,z_1,\ldots,z_e]/(xy-\pi),\qquad
 Y=V(x,y,z_1,\ldots,z_e),
\]
and give $X$ the log structure of its special fibre.  Take the full
generic fibre near $Y$ and no extra descent group.  For every selected
bounded derived continuous $\Gamma=\mathbf Z_\ell(1)$ coefficient,
there is a compatible actual Kummer realization on the geometric tube
\[
 ]Y[_{\bar\eta}=\{\,|\pi|<|x|<1\,\}\times(\mathbb D^\circ)^e.
\]
For this source the actual root star, actual analytic star, and tame
nearby star are identified by Theorem~\ref{thm:frag-vertical-log-star},
including algebra/module structures and selected adic recovery.  The
analytic coefficient realization is fully faithful on the selected
underlying-perfect derived category.  Thus the abstract-source
extension problem for arbitrary constructible sources is absent for
this example and its specified Kummer tower.
\end{corollary}

\begin{proof}
The chart $\mathbf N\to\mathbf N^2$, $1\mapsto(1,1)$, is vertical,
log smooth and saturated, with primitive base element and relative
rank one.  Its boundary-free neighbourhood satisfies
Theorem~\ref{thm:frag-vertical-log-star}.
At precision $r$, Proposition~\ref{prop:frag-point-finite-action-descent}
presents the coefficient and its maps over some $K_n=\mu_{\ell^n}$.
The fs Kummer cover
\[
 \operatorname{Spec}R[a,y,z]/(a^{\ell^n}y-\pi)\longrightarrow X
\]
is a $K_n$-torsor: in its fs self-fibre product the ratio of the
two roots is an $\ell^n$th root of unity, giving the deck copies.
Associated sheaves are exact and refinement agrees with inflation.
They therefore construct the log-lisse extension on the entire
selected derived source.

This covering scheme is fs log smooth but need not be saturated
over $R$; the local-acyclicity input only needs fs log smoothness.
After adjoining compatible $\pi_n$, a saturated model of its selected
geometric component is $ab=\pi_n$, with
$x=a^{\ell^n}$, $y=b^{\ell^n}$ and deck action
$(a,b)\mapsto(\zeta a,\zeta^{-1}b)$.  Refinement raises both
coordinates to $\ell^{m-n}$.  This gives the stated annulus and discs.

Apply Theorem~\ref{thm:frag-vertical-log-star} to the constructed
coefficient.  For perfect $E,F$ their realized internal Hom is the
associated sheaf of $R\operatorname{Hom}_{R_r}(E,F)$ with conjugation
action.  Its star comparison is the realization-induced mapping map,
as verified in Lemma~\ref{lem:frag-B-local-mapping-unit}; its source is
$C^*_{\mathrm{cts}}(\Gamma,R\operatorname{Hom}_{R_r}(E,F))$.
This proves full faithfulness.  Natural reduction and the uniform
Koszul bounds give the cartesian adic statements.
\end{proof}

\begin{remark}[Why the vertical and boundary qualifications are substantive]
\label{rem:frag-horizontal-log-test}
Consider $X=\operatorname{Spec}R[t]$ with the log structure of
$\pi t=0$ and $Y=\{\pi=t=0\}$.  Its chart is
$P=\mathbf N^2$, $\rho=(1,0)$; the base lift is primitive and the
relative characteristic group has rank one.  Nevertheless, for the
constant coefficient on the whole generic fibre, smooth base change
gives $i_Y^*R\Psi_XR_r=R_r$ in degree zero.  Proposition~\ref{prop:frag-ordinary-nearby-boundary-disjoint} identifies this with
the cohomology of the actual open-disc tube.  Continuous invariants
of the rank-one relative log inertia instead have a nonzero degree-one
term $R_r(-1)$.

There is no contradiction with Theorem~\ref{thm:frag-vertical-log-star}: this log structure
is not vertical.  The logarithmic generic coefficient corresponds on
the ordinary generic fibre to $Rj_*R_r$ for $j:\mathbf G_m\to\mathbf A^1$,
not to the constant sheaf on $\mathbf A^1$.  The horizontal generic
boundary also meets $Y$.  Thus neither this forgetful direct image nor
the open-boundary condition may be dropped.  This example tests the
logarithmic chart extension, rather than an assumption of an entire-torus
product, which this open-disc tube does not satisfy.
\end{remark}


\begin{theorem}[Realization of local Kummer coefficients]
\label{thm:frag-theorem-B-local-closed}
Let $Y=\operatorname{Spec}k$ be a geometric face of a vertical fs saturated
log-smooth algebraization, with the primitive split chart
$\mathbf N\rho\to P$ of Definition~\ref{def:frag-vertical-log-B-domain},
and suppose the generic boundary misses $Y$.  Fix the chart's relative
base-root neutralization, an integral splitting
$P^{\mathrm{gp}}=\mathbf Z\rho\oplus Q$, and no additional descent
group.  Canonicity below is relative to these data.  Every selected
bounded derived continuous coefficient for
$\Gamma=\operatorname{Hom}(P^{\mathrm{gp}}/\mathbf Z\rho,\mathbf Z_\ell(1))$
has an actual lisse Kummer generic realization.  On this source,
\[
 Rg_*^{\mathrm{act}}E_r\simeq C^*_{\mathrm{cts}}(\Gamma,E_r)
 \simeq i_Y^*R\Psi^t_X\mathcal E_r
 \simeq R\Theta(\mathcal E_r^{\mathrm{an}}|_{]Y[}).
\]
The same holds for the selected cartesian adic systems, with uniform
bounded recovery, and after rationalization.  The comparisons preserve
the unit algebra and coefficient modules and all coefficient maps.
The analytic realization is fully faithful on this selected source.
The coefficient extension is constructed in the proof.
\end{theorem}

\begin{proof}
Here is the source construction in arbitrary rank.  Primitivity allows
an integral splitting $P^{\mathrm{gp}}=\mathbf Z\rho\oplus Q$; fix one
compatible with the chosen neutralization.  The absolute chart root
cover adjoining $P/\ell^n$ is a torsor in the Kummer etale topology
with group
\[
 A_n=\operatorname{Hom}(P^{\mathrm{gp}},\mu_{\ell^n})
       =K_n\times\mu_{\ell^n}.
\]
The second factor is the base-root direction.  Extend a $K_n$-module
along the projection $A_n\to K_n$ and form its associated sheaf on
$X_{\mathrm{ket}}$.  Over the chosen geometric base-root lift the
relative action is the original $K_n$-action.  This is an actual
finite Kummer torsor construction, even though its underlying scheme
map is not etale on the special fibre.  Its generic restriction is
ordinary etale because the log structure is vertical.

Associated sheaves are exact in the module variable.  Compatible
chart root covers make this construction commute with inflation.
Pointwise finite derived-action descent therefore extends it to all
selected bounded derived continuous coefficients, with their morphisms
and finite diagrams.  The result is an actual log-lisse extension of
the kind required by Theorem~\ref{thm:frag-vertical-log-star}.
Apply that theorem, Proposition~\ref{prop:frag-relative-log-root-interface} and Corollary~\ref{cor:frag-actual-root-derived-realization} to obtain the displayed comparisons.
The analytic term means $R\nu_*\mu^*$, as specified in
Definition~\ref{def:frag-analytic-specialization-sites}.

For full faithfulness apply the same comparison to the realized
internal Hom of two underlying-perfect coefficients, as in
Corollary~\ref{cor:frag-node-full-Kummer-star}.  Its source is the
continuous equivariant mapping complex; the comparison is the map
induced by the same associated-sheaf realization; the classifying-map
and unit verification is spelled out in
Lemma~\ref{lem:frag-B-local-mapping-unit}.
Proposition~\ref{prop:frag-local-adic-full-faithfulness} gives the
adic functor and full faithfulness; uniform Koszul bounds give its
point-valued cohomology recovery.  The construction
is relative to the fixed neutralization; descent between arbitrary
choices or different characteristic strata is not asserted here.
\end{proof}



\begin{lemma}[The local mapping comparison is induced by realization]
\label{lem:frag-B-local-mapping-unit}
Fix the geometric-face chart, relative splitting and base-root lift
of Theorem~\ref{thm:frag-theorem-B-local-closed}, and write $\mathcal R_r$ for its actual
associated-sheaf realization on $T=]Y[_{\bar\eta}$.
For underlying-perfect continuous coefficients $V,W$, the canonical
map induced by that functor is an equivalence
\[
 R\operatorname{Hom}_{\Gamma,\mathrm{disc}}(V,W)
 \longrightarrow
 R\operatorname{Hom}_{T_{\acute et}}(\mathcal R_rV,\mathcal R_rW).
\]
This identifies the mapping map itself, not just its source and
target up to an unspecified isomorphism.  Its finite-precision reduction squares induce the adic assertion
of Proposition~\ref{prop:frag-local-adic-full-faithfulness}.
\end{lemma}

\begin{proof}
Use the derived model of Lemma~\ref{lem:frag-derived-analytic-site-model}.
First make the analytic site comparison explicit.  Over a geometric
face, Theorem~\ref{thm:frag-vertical-log-star} takes values in
$R\Gamma(T_{q\acute et},\mu^*F)$.  For the bounded constructible $R_r$-torsion complexes used here,
start with their ordinary etale cohomology sheaves $F$.  The canonical map
\[
 R\Gamma(T_{\acute et},F)\longrightarrow
 R\Gamma(T_{q\acute et},\mu^*F)
\]
is an equivalence by
\cite[Theorem 3.3(i)--(ii), with $U=X=T$]{BerkovichFormalI}.
Finite truncation triangles extend this to our bounded complexes.
This is the natural map of the tensor inverse image $\mu^*$, so it
respects evaluation and tensor pairings.  More explicitly, $\mu^*$
is exact and strong symmetric monoidal; the displayed global-sections
map is its adjunction-unit map.  The ordinary/quasi-etale site square
for a finite etale Kummer torsor commutes on inverse images.
Naturality of the unit in that square gives compatibility with
torsor pullback.  Thus these compatibilities concern the canonical
map whose invertibility is cited, rather than additional choices
of cohomological isomorphisms.  It justifies
using ordinary analytic etale global sections in this proof; it
neither supplies $\mu^!$ nor realizes an arbitrary quasi-etale sheaf
on the ordinary etale site.

At a finite level, the associated-sheaf functor for the actual Kummer
torsor is an exact tensor inverse image.  It preserves derived tensor
products: equivariant free resolutions are underlying-flat, and their
associated sheaves are locally flat.  Its compatible inflation maps
therefore give a tensor realization after finite derived-action
descent.  In particular, for underlying-perfect $V,W$, it identifies
\[
 \mathcal R_r\bigl(R\operatorname{Hom}_{R_r}(V,W)\bigr)
 \simeq R\mathcal Hom_{T}(\mathcal R_rV,\mathcal R_rW),
\]
with the conjugation action on the coefficient on the left.
Here perfectness supplies the evaluation identity $V^\vee\otimes W$;
it is not inferred from a cohomology-dimension calculation.

Let $U=R\operatorname{Hom}_{R_r}(V,W)$ and present it at a finite
level $K_n$.  Its actual torsor classifying map supplies
\[
 \kappa_n:C^*(K_n,U_n)\longrightarrow R\Gamma(T,\mathcal R_rU),
\]
by the pullback/direct-image adjunction, as in
Proposition~\ref{prop:frag-actual-finite-galois-descent}.
These maps commute with inflation.  In the relative logarithmic
fibre the same classifying map is induced by $\Gamma\to K_n$:
the absolute chart torsor has first been restricted along the fixed
geometric base-root lift.  Consequently its map there is the
canonical inflation
$C^*(K_n,U_n)\to C^*_{\mathrm{cts}}(\Gamma,U)$.

Naturality of the log-specialization unit and of completion for
this actual torsor classifying map gives the commuting triangle
\[
 \begin{array}{ccc}
 C^*(K_n,U_n)&\longrightarrow&C^*_{\mathrm{cts}}(\Gamma,U)\\
 &\searrow\kappa_n&\downarrow\beta\gamma\\
 &&R\Gamma(T,\mathcal R_rU).
 \end{array}
\]
To see the map-level assertion, construct both diagonals by inserting
the adjunction unit for the finite torsor and then the specialization
unit.  Transitivity of these units gives the two composites, while
the completion map is their mate.  Thus the triangle compares the
same classifying morphism on the two actual sites.  It does not
follow from merely having the same deck group, and does not assert
that $\kappa_n$ is an equivalence at finite $n$.

The middle vertex and the homotopy in that triangle can be retained
explicitly.  Put $L(U)=R\Gamma(Y,i_Y^*R\Psi_XRj_*\mathcal R_{\eta,r}U)$.
For each finite presentation the diagram is
\[
\begin{array}{ccccc}
 C^*(K_n,U_n)&\xrightarrow{\mathrm{inf}}&C^*_{\mathrm{cts}}(\Gamma,U)
       &\xrightarrow{\gamma}&L(U)\\
 \big\downarrow\kappa_n&&\big\downarrow\beta\gamma&&\big\downarrow\beta\\
 R\Gamma(T,\mathcal R_rU)&=&R\Gamma(T,\mathcal R_rU)
       &=&R\Gamma(T,\mathcal R_rU).
\end{array}
\]
Construct its left rectangle on the augmented Cech nerve of the
actual torsor.  In degree $q$ use the $(q+1)$-fold fibre product of
that torsor, its log restriction, and its formal completion.  All
face and degeneracy maps are the pullback projections and diagonals
of this one torsor.  The specialization unit is natural for these
maps; the completion mate is obtained from the same inverse-image
square in each degree.  The unit identities therefore give a map
of augmented cosimplicial diagrams, not a choice of equalities on
cohomology.  Totalization supplies the displayed homotopy, with
compatibility under refinements from the morphisms of these nerves.
The coherent composition cells are those of these actual functors
and their adjunctions.  Applying this construction to evaluation
and tensor products proves compatibility with composition of
internal Homs.  This describes the homotopy used above and does
not identify any finite-level bar with the continuous output.

Take the filtered colimit of the source of this triangle at fixed
$r$.  Point finite-action descent identifies it with continuous
cochains; hence its diagonal is precisely the realization-induced
map.  Theorem~\ref{thm:frag-vertical-log-star} proves that $\beta\gamma$ is invertible.
Applying the internal-Hom identity above now proves the displayed
mapping equivalence, compatibly with evaluation, identities and
composition.  The finite group bars have remained full throughout.

For adic systems, use the coefficient-compatible tensor realizations
and their actual cartesian outputs.  Mapping spectra in the limit
category of such systems are the inverse limits of the finite
precision mapping spectra, so the just proved equivalences give
adic full faithfulness.  This is a statement about adic systems,
not an identification with an ordinary inverse-limit sheaf on the
small etale site.  No common finite action level for all $r$ is used.
\end{proof}

\begin{proposition}[Adic realization and its mapping spectra]
\label{prop:frag-local-adic-full-faithfulness}
The finite Kummer associated-sheaf functors induce an exact tensor
functor $\mathcal R_{\mathrm{ad}}:\mathcal S_{\mathrm{ad}}(\Gamma)
\to\mathcal P_{\mathrm{ad}}(T)$.  It preserves coefficient duals
and is fully faithful.  For two systems its mapping map is
\[
 \operatorname{Map}_{\mathcal S_{\mathrm{ad}}}(V_\bullet,W_\bullet)
 =\varprojlim_r\operatorname{Map}_{\mathcal S_r}(V_r,W_r)
 \xrightarrow{\sim}
 \varprojlim_r\operatorname{Map}_{\mathcal P_r(T)}(\mathcal R_rV_r,\mathcal R_rW_r).
\]
Both limits are limits of spectra, including their derived information.
\end{proposition}
\begin{proof}
On a finite torsor the associated-sheaf functor is an exact strong
tensor inverse image.  Hence its derived reduction square commutes
canonically, including iterated reductions.  Finite-action descent
extends these squares to the selected continuous source.  At each
precision a finite presentation is a bounded complex of finite
modules with action through a finite quotient; its underlying
perfectness makes its realized complex locally perfect on that
torsor.  Thus it has the same local Tor bounds and lies in
$\mathcal P_r(T)$.

For perfect objects, derived scalar extension preserves duals and
identifies internal Hom with dual tensor.  Tensor adds Tor intervals,
and duality sends $[a,b]$ to $[-b,-a]$; these bounds are uniform.
This proves closure of both selected categories under these
coefficient operations and proves cartesianness of their results.
It does not assert closure under arbitrary geometric pushforwards.

The universal property of a limit of stable categories computes
mapping spectra by the displayed limits for cartesian sections.
Passing to a full selected subcategory leaves those maps unchanged.
The finite-precision assertion of
Lemma~\ref{lem:frag-B-local-mapping-unit} identifies each map by
realization, and its torsor/reduction naturality identifies their
transition squares.  Taking the limit proves full faithfulness.
No interchange of homotopy groups with limits, and hence no
Mittag--Leffler assertion, is used here.  Perfect recovery for point
outputs is the separate result of
Lemma~\ref{lem:frag-point-perfect-recovery}.  Rationalization is
performed on these integral categories and mapping spectra afterwards.
\end{proof}

\begin{corollary}[Local essential image and the scope of closure]
\label{cor:frag-B-local-essential-image}
The fixed-source realization of Theorem~\ref{thm:frag-theorem-B-local-closed} is an equivalence
onto its full analytic essential image.  That image is closed under
finite cones of arbitrary analytic morphisms between its objects
and under retracts.  Thus there is no further mapping or attaching-map
existence problem inside this one local Kummer image.
\end{corollary}

\begin{proof}
The preceding lemma lifts every analytic mapping spectrum between
two realized objects to the source.  Exactness then realizes its
cones.  The selected underlying-perfect source is idempotent complete,
so full faithfulness also lifts and splits retracts.  This argument
is confined to a fixed geometric face and source; it neither glues
different characteristic strata nor identifies the distinct
proper-image realizations in
Proposition~\ref{prop:frag-B-proper-morphism-obstruction}.
\end{proof}


\section{Compact support and Verdier duality}
\label{sec:frag-supported-realization}

We compare two independently defined geometric supported objects:
analytic compact support and algebraic nearby costalks.  Geometric
duality on each side precedes dualization of the star comparison.
This order both fixes the shift and avoids defining compact support
by the desired answer.

\begin{definition}[Absolute source duals and coefficient duals]
\label{def:frag-source-duality-ledger}
Fix the compatible base dualizing objects $\omega_{\mathscr S,r}$.
For a geometric source $X$ with structural map $p$ in the finite formalism,
write $\omega_{X,r}=p^!\omega_{\mathscr S,r}$ and
\[
\mathbb D_{X,r}F:=R\underline{\operatorname{Hom}}_{R_r}
(F,\omega_{X,r}).
\]
In particular $\mathbb D_{\mathrm{an},r}$ is the absolute dual on the
geometric analytic source, and $\mathbb D_{\mathrm{near},r}$ denotes the
absolute dual on the generic algebraic source of $E_r$; it is not an
undefined dual on the value of $R\Psi^t$.
For the finite enhanced root face put
\[
\mathbb D_{\mathrm{root},n,r}E_{n,r}
:=R\underline{\operatorname{Hom}}_{R_r}
(E_{n,r},g_n^!\omega_{\mathscr S,r}).
\]
The notation $\mathbb D_{\mathrm{root}}E$ denotes its compatible continuous
system only under Assumption~\ref{ass:frag-finite-trace-domain} and the adic
recovery statement.  On a product coefficient, with $N=d+e$,
\[
\mathbb D_{\mathrm{root},n,r}(\mathcal L_r\boxtimes V_{n,r})
\simeq (\mathbb D_{\mathscr S,r}\mathcal L_r\boxtimes
V_{n,r}^{\vee_0})(N)[2N],
\quad V_{n,r}^{\vee_0}=R\underline{\operatorname{Hom}}_{R_r}(V_{n,r},R_r),
\]
with contragredient action.  We abbreviate
$\mathcal L^\dagger=\mathbb D_{\mathscr S}\mathcal L$.
The factor $(N)[2N]$ is already inside the source dual; it is inserted only
once.  $A_M^*$ and $A_{M,e}^!$ below use base-linear coefficient duality,
not a second copy of base Verdier duality.
\end{definition}


\begin{lemma}[Open-neighbourhood locality of the nearby costalk]
\label{lem:frag-nearby-costalk-open-locality}
Let $k:W\hookrightarrow\overline X$ be an open neighbourhood of $Y$,
with $i_Y=k_s i_{Y,W}$.  Canonical open base change gives
\[
 k_s^*R\Psi_{\overline X}F\simeq R\Psi_W k_{\bar\eta}^*F,
 \qquad
 i_Y^!R\Psi_{\overline X}F\simeq
 i_{Y,W}^!R\Psi_W k_{\bar\eta}^*F.
\]
These also hold for tame nearby cycles and retain coefficient maps.
\end{lemma}
\begin{proof}
Open restriction is localization of the ordinary etale site.  Apply
Lemma~\ref{lem:frag-slice-derived-base-change} to the generic direct
image in nearby cycles and commute the special inverse images.
This proves the first identity with its canonical map.  The second
uses transitivity of exceptional pullback and $k_s^!=k_s^*$.
Wild invariants commute with restriction, giving the tame assertion.
For $W$ obtained by deleting the generic boundary closure,
$W_{\bar\eta}\subset U_{\bar\eta}$; both restrictions of
$Rj_{\eta,*}E$ and $Rj_{\eta,!}E$ equal $E|_{W_{\bar\eta}}$.
Thus duality is applied on a separated finite-type neighbourhood to
the bounded finite-Tor lisse source and its smooth source dual.
The displayed identities return this local costalk to the global
notation, without a global wild-triviality hypothesis.
\end{proof}

\begin{proposition}[Nearby-cycle duality]
\label{prop:frag-S-nearby-duality}
Let $\overline X/R$ be separated of finite type and use geometric nearby
cycles with $R_r$-coefficients.  With dualizing complexes relative to the
geometric generic and special points, the canonical comparison is
\[
 R\Psi_{\overline X}\mathbb D_{\overline X_{\bar\eta},r}F
 \xrightarrow{\sim}
 \mathbb D_{\overline X_s,r}R\Psi_{\overline X}F
\]
for the bounded constructible finite-Tor coefficient class used here.
The duality map is \cite[Corollary 3.8]{LuZhengNearbyDuality};
Corollary 3.9 supplies bounded constructibility and finite-Tor
preservation.  These are assertions for finite-type schemes.
The duality comparison preserves its evaluation pairing; see also
\cite[\S3.1]{LuZhengNearbyDuality}.

Suppose that, near a selected face $Y$, the model is vertical log smooth,
its generic fibre is smooth of pure dimension $N$, the generic boundary
misses $Y$, and $F$ is the generic restriction of a lisse Kummer complex
with the tame descent convention of
Definition~\ref{def:frag-vertical-log-B-domain}.
Then both $R\Psi F$ and $R\Psi\mathbb D F$ have trivial wild action on
their cohomology in an open neighbourhood of $Y$.  Consequently the
ordinary duality identity supplies the tame identity needed for the
costalk formula on that neighbourhood.
\end{proposition}

\begin{proof}
In \cite[\S3.1]{LuZhengNearbyDuality}, the base is a Henselian
valuation trait and the coefficient ring is Noetherian torsion,
annihilated by an integer invertible on the base.  Here it is
$\operatorname{Spec}R$ and $R_r$, and our bounded finite-Tor
constructible input belongs to their $D^-_c$ and $D_{cft}$.
Their target is $\overline X_s\bar\times_s\eta$; restriction to the
chosen geometric generic point gives ordinary geometric nearby
cycles.  Under this restriction their relative dualizing objects
become those over $\bar\eta$ and $\bar s$ (Remark 3.6).
Thus a smooth generic $N$-fold has dualizing object $R_r(N)[2N]$;
there is no additional trait shift.  Pairing compatibility is the
one following their Corollary 3.9.  We use no Artin-stack extension
of that result; the passage to tame cycles and costalks is below.

For the second assertion use Lemma~\ref{lem:frag-nearby-costalk-open-locality}
and remove the closure of the generic boundary.
All ensuing tame identities, $j_*$/$j_!$ identifications and their
Verdier mates are taken on this restricted open neighbourhood; only
afterwards is $i_Y^!$ applied.  They are not assertions of wild
triviality for the unlocalized global $Rj_*F$.
On the remaining neighbourhood both $j_*$ and $j_!$ of the coefficient
restrict to that generic coefficient.  Since the generic fibre is smooth,
its absolute dual is $F^{\vee_0}(N)[2N]$; dualizing the log-lisse extension
and adding this single twist and shift gives a log-lisse extension of
the dual.  Log local acyclicity and the log/ordinary exchange of
Lemma~\ref{lem:frag-log-nearby-geometric-input} apply to both extensions.
Write $W$ for this one neighbourhood, $G=\operatorname{Gal}(\bar K/K)$
and $P\subset G$ for wild inertia.  The action of $G$ on the geometric
log special fibre factors through $G/P$
\cite[\S8.1]{FujiwaraKatoNakayama}.  The chosen log-lisse extension
$\mathcal L$ is pulled back from the specified Kummer coefficient;
its transport is the induced geometric transport, with no additional
wild representation on its coefficients.  Consequently $P$ acts
trivially on that log special fibre and its restricted coefficient.
The local-acyclicity unit and the log/ordinary exchange give
\[
 C_F:=R\Psi_W F\simeq
 R\varepsilon_{s,*}(\mathcal L|_{W_{\bar s}^{\log}}).
\]
Here the comparison is the exchange map in the category with smooth
$G$-action, as in \cite[\S8.2, (8.2.3)]{FujiwaraKatoNakayama},
not an isomorphism after forgetting the action.  Functoriality of
$R\varepsilon_{s,*}$ therefore makes $P$ act trivially on every
cohomology sheaf of $C_F$.  This is the argument of
\cite[Corollary 8.4, (8.4.1)--(8.4.2)]{FujiwaraKatoNakayama};
the extension from constant to the selected log-lisse coefficients
is supplied by Lemma~\ref{lem:frag-log-nearby-geometric-input}.
It applies also to the dual extension on the same $W$: dualization
preserves the trivial wild transport, and the prime-to-$p$ Tate
factor has trivial wild action.

For either nearby complex $C$, the canonical map from its smooth
wild invariants to $C$ has cohomology maps
$H^j(C)^P\to H^j(C)$.  Exactness of these invariants, as proved in
Lemma~\ref{lem:frag-wild-inertia-criterion}, makes all these maps
isomorphisms.  Thus tame and ordinary nearby cycles agree on $W$,
before applying $i_Y^!$.  This order is
essential: trivial wild action on $i_Y^*C$ alone would not identify
$i_Y^!C$ with the costalk of its tame part.  The actual algebraic
identities $\mathbb D i_Y^*=i_Y^!\mathbb D$ and
$\mathbb D Rj_*=Rj_!\mathbb D$ now give the asserted tame costalk
duality on this domain.  No invariants under the remaining tame base
inertia are taken; its action remains external.
\end{proof}

\begin{definition}[The local supported domain]
\label{def:frag-S-local-domain}
Use the actual source of Theorem~\ref{thm:frag-theorem-B-local-closed},
with $Y$ a closed geometric point.  Let $N=d+e$ be the pure generic
dimension, and choose the supplied proper algebraic compactification
$\overline X/R$, with $Y$ disjoint from the generic-boundary closure.
Set $T=]Y[_{\bar\eta}$.  Require the selected coefficients to have the
uniform perfectness conditions already used in Theorem~\ref{thm:frag-theorem-B-local-closed}.
Write $j_\eta:U_{\bar\eta}\hookrightarrow\overline X_{\bar\eta}$ for
the given generic open and $i_Y:Y\hookrightarrow\overline X_s$.
The two actual supported objects are
\[
 A^c_r(E)=R\Gamma_c(T,\mathcal E_r^{\mathrm{an}}),\qquad
 N^c_r(E)=i_Y^!R\Psi^t_{\overline X}(Rj_{\eta,!}\mathcal E_r).
\]
The first uses compact support on the ordinary analytic etale site.
The second uses the algebraic costalk; its outer face direct image is
the identity because the face base is the geometric point.
Neither object is defined here as the dual of a star object.
\end{definition}

\begin{proposition}[actual analytic compact-support duality]
\label{prop:frag-S-analytic-duality}
In Definition~\ref{def:frag-S-local-domain}, $T$ is a separated smooth
boundaryless analytic space of pure dimension $N$.  For a selected lisse
perfect $F$ on $T$, its actual compact-support trace gives
\[
 R\Gamma_c(T,F)\xrightarrow{\sim}
 D_r R\Gamma(T,F^{\vee_0}(N)[2N]),\qquad
 D_r=R\operatorname{Hom}_{R_r}(-,R_r).
\]
For the realized Kummer coefficients the two sides are perfect.
\end{proposition}

\begin{proof}
The closed-point tube $T$ is an open analytic domain in
$\overline X_{\bar\eta}^{\mathrm{an}}$.  The boundary-disjoint hypothesis
places it in the smooth generic neighbourhood, so its structural map
is separated, smooth and boundaryless of dimension $N$.  With
$j_T:T\hookrightarrow\overline X_{\bar\eta}^{\mathrm{an}}$,
$R\Gamma_c(T,F)=R\Gamma(\overline X_{\bar\eta}^{\mathrm{an}},j_{T,!}F)$.

Apply \cite[Theorem 7.3.1]{BerkovichEtale} to this structural map,
the source coefficient $F$, and the target coefficient $R_r$.
Its hypotheses are $G\in D^-$, target in $D^+$ and $\ell\ne p$;
no quasi-compactness of the source is required.  The actual trace gives
\[
 D_rR\Gamma_c(T,F)\simeq R\Gamma(T,F^{\vee_0}(N)[2N]).
\]
For realized coefficients the right side is perfect by
Theorem~\ref{thm:frag-theorem-B-local-closed}.
The ring $R_r$ is an injective cogenerator over itself: duality is
exact and faithful, and a module whose dual is finite embeds in
its finite double dual.  The actual compact-support complex therefore
has bounded finite cohomology.  Finite $R_r$-modules are reflexive,
so, writing $C=R\Gamma_c(T,F)$, biduality gives $C\simeq D_rD_rC$.
Since $D_rC$ is already perfect by the displayed star comparison,
and $D_r$ takes bounded finite projective complexes to such complexes,
$C$ is perfect.  This proves the asserted trace-induced equivalence
without presupposing compact-support finiteness.
\end{proof}

\begin{theorem}[Compact support and nearby costalks]
\label{thm:frag-S-local-closed}
In Definition~\ref{def:frag-S-local-domain}, there is a canonical
comparison on the realized coefficient source
\[
 \theta^!_r:A^c_r(E)\xrightarrow{\sim}N^c_r(E),
\]
and both objects are naturally identified with
\[
 H^!_r(E)=D_r C^*_{\mathrm{cts}}(\Gamma,E_r^{\vee_0})(-N)[-2N].
\label{eq:frag-S-local-dual-formula}
\]
The comparison respects the common star-algebra action on coefficients.
It is compatible with coefficient reduction and has uniformly bounded,
derived-cartesian adic recovery on the actual coefficient source of
Theorem~\ref{thm:frag-theorem-B-local-closed}.
\end{theorem}

\begin{proof}
Let $\mathcal D\mathcal E_r=\mathcal E_r^{\vee_0}(N)[2N]$.
Its actual log/generic realization is obtained by dualizing the source
construction of Theorem~\ref{thm:frag-theorem-B-local-closed}, so the already proved star comparison
applies to this coefficient and to its maps.  Denote that comparison by
\[
 \beta_{\mathcal D E}:i_Y^*R\Psi^t Rj_{\eta,*}\mathcal D\mathcal E_r
          \xrightarrow{\sim}R\Gamma(T,\mathcal D\mathcal E_r^{\mathrm{an}}).
\]
On the left, Proposition~\ref{prop:frag-S-nearby-duality} and the actual finite algebraic duality identities give
\[
 \begin{aligned}
 D_r i_Y^*R\Psi^t Rj_{\eta,*}\mathcal D\mathcal E_r
 &\simeq i_Y^!\mathbb D_{\overline X_s,r}
                          R\Psi^t Rj_{\eta,*}\mathcal D\mathcal E_r\\
 &\simeq i_Y^!R\Psi^t Rj_{\eta,!}\mathcal E_r=N^c_r(E).
 \end{aligned}
\]
The last step uses source biduality.  Proposition~\ref{prop:frag-S-nearby-duality} justifies the tame exchange
on a neighbourhood before the costalk is taken.  Proposition~\ref{prop:frag-S-analytic-duality} identifies the
dual of the right side of $\beta_{\mathcal D E}$ with $A^c_r(E)$.
Define $\theta^!_r$ to be $D_r(\beta_{\mathcal D E})$ under these two
geometrically proved identifications.  This constructs a map between
the two independently defined supported objects and proves it invertible.

Applying Theorem~\ref{thm:frag-theorem-B-local-closed} to $\mathcal D\mathcal E_r$ gives
\eqref{eq:frag-S-local-dual-formula}.  The factor $(N)[2N]$ belongs to
the source dual and becomes $(-N)[-2N]$ exactly once under $D_r$.
The ordinary star comparison preserves algebra/module structures.
Its dual therefore preserves their transpose actions.  Analytic Poincare
duality and nearby duality identify those actions with the actual
cup-product/evaluation actions on compact support and costalks.
This proves the module assertion for the constructed map, with its
coherent cells.
For coefficient reduction, use the canonical coefficient-change maps
of the actual operations and their adjunction mates.  The comparison
maps above are constructed from pullback, evaluation and trace, so
these maps form the corresponding natural coefficient-change diagram.
On the star side Theorem~\ref{thm:frag-theorem-B-local-closed} has already proved it cartesian.
On perfect point-valued outputs, the canonical base-change map for
$D_r$ is an isomorphism.  Dualizing that diagram consequently proves
cartesianness for the actual supported outputs and compatibility for
$\theta^!_r$.
If the input has a common perfect amplitude $[a,b]$, the formula gives
supported amplitude in $[a+2N-d,b+2N]$, uniformly in $r$.  The same
finite construction gives the finite-constructibility bounds.  The
proved cartesian system therefore admits the adic recovery used in this
paper.  Take its coefficient inverse limit only after this verification,
and rationalize afterwards.  \end{proof}

\begin{proposition}[refinement traces and the scope of equivariance]
\label{prop:frag-S-trace-naturality}
The comparison of Theorem~\ref{thm:frag-S-local-closed} is natural for
coefficient maps and for compatible algebraized finite Kummer covers
whose selected geometric tubes are finite etale over the original tube
and whose chosen compactified models admit compatible finite morphisms
after the specified tame base change.
Under the supported/star duality pairings, actual finite-cover trace is
the adjoint of actual star pullback.  These adjunctions respect
composition.  If an actual finite semilinear action and its star-source
comparison are supplied, the supported comparison is equivariant for
that action; prime-to-$\ell$ finite descent preserves its recovery bounds.
\end{proposition}

\begin{proof}
For a finite etale cover $q:T'\to T$, the counit
$q_*q^*\to\mathrm{id}$ is the trace, locally the sum over its finite
sheets.  The equality between pairing after trace and pairing after
pullback follows from the projection formula and transitivity of the
structural trace.  The generic algebraic cover has the same finite-etale
adjunction; nearby-duality and the completion comparison carry its
natural maps to the corresponding analytic diagram.  Dualizing this
star diagram is exactly the supported diagram defining $\theta^!$.
Composition follows from composition of the counits, including their
existing adjunction coherence.  No claim is made that a finite root
Artin-stack bar equals the finite analytic covering complex.

An actual automorphism preserving the model, face and coefficient data
preserves these trace and evaluation maps.  Thus the same construction
is equivariant on the supplied finite action diagram.  Exactness of
prime-to-$\ell$ group invariants gives the last assertion.  This does
not manufacture equivariant source realizations for arbitrary choices
of a root neutralization, nor does it establish chart-change incidence.
\end{proof}

\begin{corollary}[Geometric realization of the continuous root trace]
\label{cor:frag-C2-local-closed}
On the domain of Theorem~\ref{thm:frag-S-local-closed}, with the fixed
rank-$N$ root carrier and the finite duality formalism of Theorem~\ref{thm:frag-point-C1}, there is
an integral adic comparison
\[
 A^c_{\mathrm{ad}}(E)\simeq N^c_{\mathrm{ad}}(E)
                  \simeq\mathsf G_{!,\mathrm{ad}}^{\mathrm{tr}}(E).
\]
It identifies the unit with the base-linear dualizing module of the
continuous star algebra.  Together with Theorem~\ref{thm:frag-theorem-B-local-closed} this realizes
the local star/trace datum on this source.  It does not construct a
shriek direct image on the actual infinite-root stack.
\end{corollary}

\begin{proof}
Theorem~\ref{thm:frag-point-C1} computes the continuous finite-root trace by the
right side of \eqref{eq:frag-S-local-dual-formula}, with the same
contragredient coefficient and the same single rank-$N$ source twist.
Theorem~\ref{thm:frag-S-local-closed} identifies both actual supported objects with that
formula, through the already identified continuous star functor.
Perfect duality fixes the associated unit module and its action.
The comparison is made after continuous trace assembly, not between
individual finite Artin-stack and analytic shriek images.  All adic
transitions have been verified before taking their inverse limit.
\end{proof}

\begin{corollary}[The supported degree and twist on the unit]
\label{cor:frag-S-unit-degrees}
In Theorem~\ref{thm:frag-S-local-closed}, for the unit coefficient and $L=M^\vee$,
\[
 H_c^{2N-j}(T,R_r)
 \simeq \bigwedge^j(M\otimes R_r)(j-N),\qquad 0\le j\le d,
\]
and the other cohomology groups vanish.  The same formula holds for
the nearby costalk.  In particular an open smooth disc contributes
only $R_r(-1)$ in degree $2$, whereas a semistable node annulus
contributes $R_r$ in degree $1$ and $R_r(-1)$ in degree $2$.
\end{corollary}

\begin{proof}
Theorem~\ref{thm:frag-theorem-B-local-closed} identifies the ordinary unit cohomology with
$\bigwedge^j(L\otimes R_r)(-j)$ in degree $j$.  Apply the single
$D_r(-)(-N)[-2N]$ in Theorem~\ref{thm:frag-S-local-closed} and dualize the finite free
exterior factors.  The two examples have respectively $(d,N)=(0,1)$
and $(1,1)$.  This computes the supported cohomology, not a formality
statement for the cochain algebra or a normal-pair incidence map.
\end{proof}



\section{The finite-root trace}
\label{sec:local-finite-root-trace}
\label{sec:frag-point-core}

We now identify the supported root term geometrically.  At fixed
$r$ the star maps are inflation; their duals must be the counits
of the genuine finite-stack shriek adjunctions.  This map-level
identification, followed by the continuous limit, relates the
root trace to the supported comparison of the preceding section.
The base throughout this section is a geometric point, with
$\Gamma=\mathbf Z_\ell^d$, $K_n=\Gamma/\ell^n\Gamma$ and fixed
rank-$N$ vector-bundle carriers.  There is no additional descent group.

\begin{proposition}[The enhanced finite-quotient comparison]
\label{prop:frag-finite-quotient-enhancement}
Let $A=R_r$ and let $G$ be a finite constant etale group over the
algebraically closed field $k$, with $\ell$ invertible in $k$.
Use the unbounded stable etale category $\mathcal D(BG,A)$ of
\cite[\S6.2]{LiuZhengEnhanced}.  There is a natural equivalence
\[
 \Phi_{G,A}:\mathcal D(BG,A)\simeq D(A[G]).
\]
It identifies $D_c^b(BG,A)$ with the full subcategory of bounded
finite-cohomology modules; it does not identify $D_c^b$ with all
of $D(A[G])$.  For $h:BG\to\operatorname{Spec}k$ and a surjection
$G\to K$ inducing $q:BG\to BK$, the functors and their adjunctions
are identified as follows:
\[
\begin{array}{c|c|c|c}
 f&f^*=f^!&f_*&f_!\\ \hline
 h&\text{trivial action}&R\operatorname{Hom}_{A[G]}(A,-)
                         &A\otimes^L_{A[G]}-\\
 q&\operatorname{Res}_{A[G]}^{A[K]}
   &R\operatorname{Hom}_{A[G]}(A[K],-)
   &A[K]\otimes^L_{A[G]}-
\end{array}
\]
The restriction notation in the second row means restriction along
$A[G]\to A[K]$.  These are comparisons of enhanced adjunctions,
including the geometric units, counits and their composite cells.
They commute with coefficient restriction and identify the smooth
shriek scalar-change maps with derived tensor associativity.
\end{proposition}
\begin{proof}
\emph{The category before the operations.}
The finite etale atlas $\operatorname{Spec}k\to BG$ has Cech terms
$G^s$.  Enhanced descent
\cite[Theorem 6.2.13(1)]{LiuZhengEnhanced} gives
\[
 \mathcal D(BG,A)\simeq
 \varprojlim_{[s]\in\Delta}D(G^s,A)
 =\varprojlim_{[s]\in\Delta}\prod_{G^s}D(A).
\]
The theorem applies to this smooth surjection with torsion
coefficients in the unbounded enhancement.  For a Deligne--Mumford
stack this is its etale derived category; its lisse-etale description
uses cartesian cohomology, not the whole lisse-etale category.
The displayed limit is $\operatorname{Fun}(BG,D(A))$, the category
of homotopy-coherent $G$-actions.  Evaluation at the unique object is
conservative and preserves colimits.  Its left adjoint is free action
$M\mapsto\bigoplus_{g\in G}M$, with monad $A[G]\otimes_A-$;
composition of group elements is the monad multiplication.
The free-action bar resolution is augmented to each action object.
After evaluation its augmentation is split by the identity element,
so its realization is that object.  The same bar resolves a module
for the monad.  These resolutions give inverse comparison functors
between actions and $A[G]$-modules, together with their units and
counits.  This proves the enhanced equivalence, rather than lifting a
triangulated equivalence afterwards.  Finite cohomology is detected
at the atlas, which proves the constructible restriction.
The atlas description also identifies tensor with diagonal tensor
of actions and commutes with restriction of coefficient rings.

\emph{Which adjunction is transported.}
The map of nerves for $G\to K$ restricts an action, and the nerve
map to the point gives trivial action.  Thus these are the actual
$f^*$, with their structural comparison cells.  The maps $h,q$ are
smooth of dimension zero: after a torsor base change the fibres of
$q$ are classifying stacks of its finite etale kernel.  Enhanced
smooth purity \cite[Theorem 6.2.9(2)]{LiuZhengEnhanced} identifies
$f^!$ with this same $f^*$ through the geometric trace counit.
This is an assertion about nonrepresentable smooth maps in that
theorem, not representable proper base change.

For $B=A[G]$, $C=A[K]$ (or $C=A$ for $h$), restriction has adjoints
$C\otimes_B^L-$ and $R\operatorname{Hom}_B(C,-)$ on unbounded
module categories.  The comparison with geometric $f_!$ is the
Yoneda comparison induced by the following specified chain:
\[
\begin{aligned}
 \operatorname{Map}_C(\Phi f_!E,N)
 &\simeq\operatorname{Map}_{\mathcal D(BG,A)}(E,f^!\Phi^{-1}N)\\
 &\simeq\operatorname{Map}_B(\Phi E,\operatorname{Res}N)\\
 &\simeq\operatorname{Map}_C(C\otimes_B^L\Phi E,N).
\end{aligned}
\]
The first equivalence is the geometric adjunction, the second uses
its fixed purity map and the atlas equivalence.  Consequently the
geometric counit maps to the tensor adjunction counit, namely
$C\otimes_B^L\operatorname{Res}N\to N$ induced by the action.
The unit is $M\to\operatorname{Res}(C\otimes_B^LM)$.
Using $f^*\dashv f_*$ in the same way identifies its unit and
counit with those of Hom.  These mapping-space equivalences are
natural before Yoneda is applied.  For composable quotients their
pasting is tensor associativity (and right-adjoint composition).
The composite-adjunction identities therefore supply the same
composition cells, including on transformations.

\emph{Coefficient maps and the carrier.}
For $A\to A'$ write $S=-\otimes_A^LA'$ and $U$ for restriction.
For these smooth maps, purity and inverse image give the actual
coefficient-restriction square $f_A^!U\simeq Uf_{A'}^!$.
Its left mate is the geometric map $S f_{A,!}\to f_{A',!}S$.
Under the just constructed comparisons this is precisely
\[
 S\bigl(C\otimes_B^LM\bigr)
   \simeq C'\otimes_{B'}^L(SM).
\]
For example the whole counit square, with $N$ a $C$-module, becomes
\[
\begin{array}{ccc}
 S\Phi_K q_{A,!}q_A^*\Phi_K^{-1}N&\xrightarrow{S\Phi(\epsilon_q)}&SN\\
 \downarrow\wr&&\Vert\\
 C'\otimes_{B'}^L\operatorname{Res}(SN)&\xrightarrow{\mathrm{act}}&SN.
\end{array}
\]
Both horizontal arrows are the action augmentation; the vertical
arrow is the geometric scalar mate followed by $\Phi$.
The same argument applies to $h$.  It does not invoke the
perfect-coefficient-map hypothesis of general coefficient-change
theorems: $R_{r+1}\to R_r$ need not be perfect.

For the fixed carrier $p:\mathbf A^N_{BG}\to BG$, smooth purity
has its single factor $A(N)[2N]$.  The projection formula and the
representable Thom trace identify
$p_!p^*W=W(-N)[-2N]$.  Its trace and purity commute with coefficient
restriction and with pullback of this bundle along $q$.
Consequently inserting $p$ in the preceding square tensors it
with $A'(-N)[-2N]$; the carrier refinement is a bundle pullback
with the identity on the affine-space direction.  Transitivity of
these actual traces and the composite-adjunction identities give
the joint root/scalar/Thom cells used below.  This verifies the
specified arrows, not merely the resulting complexes.
\end{proof}

\begin{proposition}[Finite geometric bars and their duals over a point]
\label{prop:frag-point-geometric-bars}
Let $h_n:BK_n\to\Spec k$, let $p_n:X_n\to BK_n$ be a rank-$N$ vector
bundle, and put $g_n=h_np_n$.  For an underlying-$R_r$-perfect coefficient
$V_n$ on $BK_n$, write $E_n=p_n^*V_n$ and
$V_n^\vee=R\operatorname{Hom}_{R_r}(V_n,R_r)$ with contragredient action.
In the finite six-functor theory with smooth purity, composition, the
projection formula, and their adjunctions, there are natural identities
\[
 Rg_{n,*}E_n\simeq R\operatorname{Hom}_{R_r[K_n]}(R_r,V_n),
 \qquad
 Rg_{n,!}E_n\simeq
 \bigl(R_r\otimes^L_{R_r[K_n]}V_n\bigr)(-N)[-2N],
\]
where $R_r$ is the trivial right module in the second formula.  In
$D(R_r)$ the second term is canonically
\[
 R\operatorname{Hom}_{R_r}
 \bigl(C^*(K_n,V_n^\vee),R_r\bigr)(-N)[-2N].
\]
For compatible vector bundles under $K_m\to K_n$, the geometric trace
is the dual of full-bar inflation.  The same conclusions transport to
finite nilpotent thickenings with these reductions.
\end{proposition}

\begin{proof}
Proposition~\ref{prop:frag-finite-quotient-enhancement} identifies the
actual enhanced adjunctions for $h_n$ with trivial action, derived
invariants and derived coinvariants.  In particular the latter
computes the given geometric $h_{n,!}$ with its counit.
Homotopy invariance and the compactly supported Thom isomorphism for the
vector bundle give $Rp_{n,*}p_n^*V_n=V_n$ and
$Rp_{n,!}p_n^*V_n=V_n(-N)[-2N]$.  Composition proves the two formulas.

Use the full free group-bar resolution.  Its duality pairing with the
opposite bar identifies derived invariants of $V_n^\vee$, dualized over
$R_r$, with derived coinvariants of $V_n$.  This can be checked with a
bounded finite-module representative of $V_n$: $R_r$ is self-injective,
so ordinary $R_r$-linear dual is exact on finite modules, and each bar
diagonal against that bounded representative is finite.  Biduality is
termwise finite-module biduality.  Equivalently the calculation uses the
Frobenius trace on $R_r[K_n]$ and the group inversion on the opposite bar.
Thus the internal Hom here is literally the one in $D(R_r)$; no
identification of two different ambient dualities is being assumed.
The positive unbounded bar becomes a negative unbounded chain complex,
not another cochain complex.

For a quotient of groups, inflation is the adjunction unit for pullback
and derived invariants.  Dualizing its bar pairing gives the counit for
derived coinvariants and trivial relative action.  Since these are the
geometric adjunctions just identified, this is the geometric trace.
Adjunction composition supplies its coherent composition law.
Lemma~\ref{lem:frag-central-refinement-purity} gives transport through
nilpotent reductions.  This proof neither replaces $h_!$ by $h_*$ nor
asserts that every unbounded finite star bar commutes with derived
coefficient reduction.
\end{proof}

\begin{definition}[Finite presentations at a fixed precision]
\label{def:frag-compatible-presenting-tails}
Let $\mathcal C_{n,r}$ be the full stable subcategory of
$D(R_r[K_n]\text{-}\mathrm{Mod})$ on underlying-$R_r$-perfect objects.
Use the stable localization of complexes, with inflation along
$K_m\to K_n$.  Write $\mathcal C_r=\mathcal S_r(\Gamma)$.
Proposition~\ref{prop:frag-point-finite-action-descent} identifies
$\operatorname*{colim}_n\mathcal C_{n,r}$ with $\mathcal C_r$.
A finite presentation of $V\in\mathcal C_r$ is $W\in\mathcal C_{n,r}$
with an equivalence $\operatorname{Inf}_n^\Gamma W\simeq V$.
Its tail uses $W_m=\operatorname{Inf}_n^mW$ and $E_m=p_m^*W_m$.
Root composition cells are those of inflation itself.
When computing a coefficient map, its finite source, target and
comparison can be placed at a common root level.  This convention
does not choose a lift of the entire adic diagram to finite levels.
\end{definition}

\begin{lemma}[Fixed-precision trace functors and coefficient maps]
\label{lem:frag-presenting-system-assembly}
The geometric root counits define an exact functor
$H_{!,r}:\mathcal C_r\to D(R_r)$.  For a finite presentation $W$
at level $n$, its value is naturally
\[
 \widehat H_{n,r}(W)=R\varprojlim_{m\ge n}
          Rg_{m,!}p_m^*\operatorname{Inf}_n^mW.
\]
For each coefficient reduction $S=-\otimes^L_{R_{r+1}}R_r$ there
is a natural transformation
\[
 b_r:S H_{!,r+1}\longrightarrow H_{!,r}S.
\]
These transformations have the composition cells for successive
reductions.  Their construction precedes the proof that they are
isomorphisms.
\end{lemma}
\begin{proof}
At fixed $(n,r)$, the functor into
$\operatorname{Fun}(\{m\ge n\}^{\mathrm{op}},D(R_r))$ sends $W$
to the displayed geometric complexes.  The arrow $m\to l$, $m\ge l$,
is the transformation
\[
 Rg_{m,!}q_{m,l}^*E_l
 \simeq Rg_{l,!}Rq_{m,l,!}q_{m,l}^!E_l
 \xrightarrow{Rg_{l,!}\epsilon_{m,l}}Rg_{l,!}E_l.
\]
Here $q^*=q^!$ and the geometric/module counit comparison are
those of Proposition~\ref{prop:frag-finite-quotient-enhancement}.
Inflation and these derived adjunctions are functors and transformations
on stable module localizations.  Their composite-counit law defines
the diagram on all simplices.  Composing with the limit functor
therefore defines $\widehat H_{n,r}$ on mapping spectra as well.

Restriction to a cofinal tail gives a natural equivalence
$\widehat H_{n,r}\simeq\widehat H_{l,r}\operatorname{Inf}_n^l$.
The comparison for $n\le l\le m$ is the composition of restriction
maps of diagrams, so these equivalences form a coherent cocone.
The universal property
\[
 \operatorname{Fun}(\operatorname*{colim}_n\mathcal C_{n,r},D(R_r))
 \simeq\varprojlim_n\operatorname{Fun}(\mathcal C_{n,r},D(R_r))
\]
constructs $H_{!,r}$, with presentation comparisons and higher maps.
Limits in a stable category are exact, so these are exact functors.

Scalar reduction $S:\mathcal C_{n,r+1}\to\mathcal C_{n,r}$ commutes
with inflation.  The left-adjoint coinvariant calculation and the
carrier Thom trace give finite-level natural maps
\[
 S(Rg_{m,!}p_m^*W_m)\longrightarrow Rg_{m,!}p_m^*(SW_m).
\]
The counit square in Proposition~\ref{prop:frag-finite-quotient-enhancement}
identifies these particular geometric maps with scalar extension of
the coinvariant augmentation and the same Thom trace.  They therefore
commute with the root counits.  Compose them with the canonical comparison
$S\lim Y_m\to\lim SY_m$ to obtain a transformation on
$\mathcal C_{n,r+1}$.  Tail restriction commutes with this
construction.  The preceding equivalence of functor categories,
now also on natural transformations, descends it to $b_r$.

For three coefficient rings, the direct and iterated maps from
scalar extension of a limit are the same map of limit cones;
finite scalar extension and the coinvariant augmentations have
their usual associativity cells.  This proves the composition cells
before any presentation is chosen for an adic object.  Applying
these functors and transformations to its original cartesian section
produces the output diagram.  Thus no inverse system of spaces of
choices of finite presenting diagrams is needed.
\end{proof}

\begin{lemma}[The trace pairing and the specified coefficient arrow]
\label{lem:frag-point-trace-diagram-descent}
The fixed-precision geometric functor has a natural equivalence
\[
 \phi_r:H_{!,r}(V)(N)[2N]\xrightarrow{\sim}
 D_rX_r(V),\qquad X_r(V)=C^*_{\mathrm{cts}}(\Gamma,V^\vee).
\]
Put $A=R_{r+1}$, $B=R_r$, $S=-\otimes_A^LB$ and omit the common
Thom factor.  Let
$\alpha:S X_{r+1}(V)\to X_r(SV)$ be the actual star coefficient map,
and let $\delta:S D_A X\to D_B(SX)$ be the map adjoint to scalar
extension of evaluation.  Then the previously constructed $b_r$
satisfies the following square, whose lower arrow points to the left:
\begin{equation}
\begin{array}{ccc}
 S H_{!,r+1}(V)(N)[2N]&\xrightarrow{\ b_r\ }&H_{!,r}(SV)(N)[2N]\\
 {\scriptstyle\delta\, S\phi_{r+1}}\downarrow&&
                    \downarrow{\scriptstyle\phi_r}\\
 D_B(SX_{r+1}(V))&\xleftarrow{\ D_B\alpha\ }&D_BX_r(SV).
\end{array}
\label{eq:frag-specified-trace-reduction-square}
\end{equation}
In particular it identifies the actual $b_r$ with standard perfect-dual
reduction once continuous star perfectness and base change are used.
\end{lemma}
\begin{proof}
\emph{The fixed-precision pairing.}
For a presentation $W$ put
\[
 X_m=C^*(K_m,W_m^\vee),\qquad
 Y_m=(Rg_{m,!}p_m^*W_m)(N)[2N].
\]
Proposition~\ref{prop:frag-point-geometric-bars} supplies the actual
pairings $e_m:Y_m\otimes X_m\to R_r$ and equivalences $Y_m\simeq D_rX_m$.
For $m\ge l$, write $t_{m,l}:Y_m\to Y_l$ for the geometric counit
and $a_{l,m}:X_l\to X_m$ for inflation.  Their paired identity is
\[
 e_l(t_{m,l}\otimes1)=e_m(1\otimes a_{l,m}).
\]
It follows from the invariants-unit/coinvariants-counit adjunction
on the opposite full bars; Frobenius evaluation takes the identity
coefficient, with group inversion, and never divides by $|K_m|$.
The identity is compatible with composition by the counit law.
Hence $Y=\lim_mY_m$ pairs with $X=\operatorname*{colim}_mX_m$.
Internal Hom sends this colimit to the dual limit, proving $Y\simeq D_rX$.
Finite-action descent, applied to maps from the trivial object to
$V^\vee$, identifies $X$ with $X_r(V)$.  The functor construction in
the preceding lemma descends the pairing itself, so $\phi_r$ is a
natural transformation, not just an isomorphism on objects.

\emph{Coefficient change before taking either limit.}
Choose a finite presentation over $A$ and reduce it to $B$; this
computes the natural transformation on the finite source category.
Use primes for the $B$-complexes.  Let $t_m:SY_m\to Y'_m$ be the
finite shriek coefficient map and let $\alpha_m:SX_m\to X'_m$ be
the usual derived-invariants coefficient map, using
$S(W_m^\vee)\simeq(SW_m)^\vee$.  In $D(B)$ one has
\[
 e'_m(t_m\otimes\alpha_m)=S e_m:
                SY_m\otimes_B SX_m\longrightarrow B.
\tag{*}
\]
For a map-level verification, denote trivial action by $T_A$, its
left adjoint (coinvariants) by $Q_A$, and its right adjoint
(invariants) by $C_A$.  The finite pairing is the composite
\[
\begin{aligned}
 Q_A(W_m)\otimes C_A(W_m^\vee)
 &\simeq Q_A(W_m\otimes T_A C_A(W_m^\vee))\\
 &\longrightarrow Q_A(W_m\otimes W_m^\vee)
 \longrightarrow Q_A(T_AA)\longrightarrow A.
\end{aligned}
\]
The first arrow uses the invariants counit, the second coefficient
evaluation, and the last the coinvariants counit.  The projection
identity follows from tensor associativity for derived coinvariants.
This is the full-bar contraction of
Proposition~\ref{prop:frag-point-geometric-bars}, hence the geometric
pairing already used in $\phi_r$.

By definition, $\alpha_m$ is adjoint to scalar extension of the
invariants counit
$T_A C_A(W_m^\vee)\to W_m^\vee$, using
$S(W_m^\vee)\simeq(SW_m)^\vee$.  Scalar change commutes with
$T_A$ and the left-adjoint coinvariant tensor $Q_A$ and carries
coefficient evaluation to evaluation over $B$.  Applying it to
the displayed composite therefore gives exactly
$e'_m(t_m\otimes\alpha_m)$.  This proves (*) for the specified
adjunction maps, without an interchange of derived Hom and tensor.
The Thom factor is obtained from the same scalar-compatible
vector-bundle trace.  No finite star base-change isomorphism or
truncation of a full bar is used.

\emph{Passage to the continuous pairing.}
Let $\pi_m:Y\to Y_m$ and $\pi'_m:Y'\to Y'_m$ be projections.
The construction of $b_r$ in the preceding lemma says precisely
\[
 \pi'_m b_r=t_m S\pi_m.
\]
Since $S$ preserves colimits, $SX=\operatorname*{colim}_mSX_m$.
Restrict the two pairings on $SY\otimes SX$ to each $SX_m$.
The displayed projection identity and (*) show that both restrictions
are $S(e_m(\pi_m\otimes1))$.  Their compatible equality on this
colimit diagram proves
\[
 e'(b_r\otimes\alpha)=S e.
\]
Tensor--Hom adjunction turns this identity into exactly the displayed
square.  All equalities here are the indicated natural homotopies of
bar, evaluation and limit-cone transformations; they descend through
the functor-category equivalence.  No tensor/inverse-limit interchange
has been used.

Finally the length-$d$ finite free augmentation Koszul resolution over
$R_r[[\Gamma]]$ computes the continuous invariants.  The identification with continuous invariants preserves the augmentation
and its right-adjoint counit: it is the comparison of maps from the
trivial object under finite-action/torsion descent.  Therefore $\alpha$
is carried to the canonical derived-Hom scalar map for this resolution.
Finite-projective base change makes that map invertible; the
underlying-perfect coefficient makes $X$
perfect.  Thus $\delta$ is invertible too, and the square gives
\[
 \phi_r b_r=(D_B\alpha)^{-1}\delta\,S\phi_{r+1}.
\]
This is the standard reduction of a perfect dual.  It proves the
claim for the original geometric coefficient arrow.  Successive
reductions retain the composition cells already constructed for
$b_r$ and for evaluation.
\end{proof}

\begin{proposition}[Six operations on the point carriers]
\label{prop:frag-point-finite-domain-verified}
Let $k$ be algebraically closed of characteristic different from $\ell$,
let $K_n=(\mu_{\ell^n})^d$, and let
$p_n:Z_n=\mathbf A^N_{BK_n}\to BK_n$ be the fixed rank-$N$ carrier.
For $g_n:Z_n\to\operatorname{Spec}k$, the published finite-coefficient
Artin-stack operations supply the point finite trace domain used
in Theorem~\ref{thm:frag-point-C1}.  Its ambient dual is literally
$R\operatorname{Hom}_{R_r}(-,R_r)$ in $D(R_r)$.
Thus Assumption~\ref{ass:frag-finite-trace-domain} is not an additional
hypothesis for this sequence of split point carriers.
This assertion does not verify that assumption on a general base.
\end{proposition}

\begin{proof}
The base field is excellent, and its finite-type schemes have finite
$\ell$-cohomological dimension in each fixed dimension.  The ring
$R_r$ is an Artin local Gorenstein ring of characteristic invertible
on the base.  Thus these finite-type Deligne--Mumford stacks and
coefficients are in the domain of \cite{LaszloOlssonFinite}.
The smooth-purity identification for $BK_n\to\operatorname{Spec}k$
is $h_n^!=h_n^*$: its stack dimension is zero
\cite[Corollary 4.6.2]{LaszloOlssonFinite}.
For the enhancement we use \cite[\S6.2, Theorems 6.2.9 and
6.2.13]{LiuZhengEnhanced}, as instantiated in
Proposition~\ref{prop:frag-finite-quotient-enhancement}.  That
proposition constructs the category comparison from atlas descent
and gives the mapping-space chain transporting each geometric
adjunction.  The classical operations are the ones recovered on
homotopy categories; the enhanced transformations used here are
fixed by this construction.
The published classifying-stack calculation
\cite[Example 4.9.2]{LaszloOlssonFinite} in particular retains the
negative unbounded shriek complex; properness of $BK_n$ does not
identify it with the positive star complex.

Smooth purity and the usual compactly supported Thom trace for the
representable vector-bundle map supply the single factor
$(-N)[-2N]$.  The finite free group bar, paired with its opposite
bar over the self-injective ring $R_r$, identifies its linear dual
with the chain complex for coinvariants on bounded finite-module
coefficients.  This is the actual point Verdier dual in $D(R_r)$,
including on the one-sided unbounded bar outputs.  Inflation and
its paired map are precisely the already identified geometric
adjunction maps; their composition cells are those of these
adjunctions.  The coefficient/counit square and its carrier extension are
those of Proposition~\ref{prop:frag-finite-quotient-enhancement};
no unspecified enhancement of a triangulated equivalence is used.

Coefficient extension commutes with the left-adjoint coinvariant
and Thom constructions.  No corresponding unconditional finite-bar
star reduction is asserted.  Internal Hom in $D(R_r)$ sends the
star colimit to the trace limit, and only the resulting finite
continuous Koszul complex is used for bounded reduction and adic
recovery.  This verifies the ambient requirements actually used
by Theorem~\ref{thm:frag-point-C1}.  If the fixed carriers are replaced by the stipulated
finite nilpotent thickenings, the representable nil-transport lemma
applies; no nonrepresentable proper star/shriek equality is used.
\end{proof}

\begin{theorem}[The continuous trace of finite root carriers]
\label{thm:frag-point-C1}
Use the point carriers of Proposition~\ref{prop:frag-point-geometric-bars};
their six-functor domain is verified in
Proposition~\ref{prop:frag-point-finite-domain-verified}.  Let $V_\bullet$ be a
cartesian source system of Definition~\ref{def:frag-local-adic-source-target},
now with $\Gamma=\mathbf Z_\ell^d$, and write
$V=R\varprojlim_rV_r$ for its underlying perfect point complex.
Put $\Lambda=\mathbf Z_\ell[[\Gamma]]$.  Finite presentations of $V_r$
exist by Proposition~\ref{prop:frag-point-finite-action-descent}.
Use the fixed-precision functor and coefficient transformations of
Lemma~\ref{lem:frag-presenting-system-assembly}; each value may be
computed on any inflated finite presenting tail with its geometric
counit transitions.  No compatible choice across all $r$ is required.
Then
\[
 H_{!,r}(V_r):=R\!\varprojlim_{n,\operatorname{Tr}^!}Rg_{n,!}E_{n,r}
 \simeq
 R\operatorname{Hom}_{R_r}
 \bigl(C^*_{\mathrm{cts}}(\Gamma,V_r^\vee),R_r\bigr)(-N)[-2N].
\]
The inverse-limit output $H_{!,r}$ is presentation-independent
and perfect over $R_r$; the individual $Rg_{n,!}E_{n,r}$ may remain
unbounded.  The coefficient maps $b_r$ constructed above are
equivalences after derived reduction, and
\[
 R\!\varprojlim_r H_{!,r}(V_r)
 \simeq R\operatorname{Hom}_{\mathbf Z_\ell}
 \bigl(K_\Gamma(V_\bullet^\vee),\mathbf Z_\ell\bigr)(-N)[-2N].
\]
Write
\[
 \mathsf G_{!,\mathrm{ad}}^{\mathrm{tr}}(V)
 :=R\!\varprojlim_rH_{!,r}(V_r),\qquad
 \mathsf G_!^{\mathrm{tr}}(V)
 :=\mathsf G_{!,\mathrm{ad}}^{\mathrm{tr}}(V)[1/\ell].
\]
Thus the root supported term is defined by the actual finite-carrier
trace system before it is identified with the continuous dual.
Here $V^\vee$ denotes the coefficient dual system and its recovered
point complex.  The Iwasawa derived Hom in this formula is continuous
in the sense of Definition~\ref{def:frag-local-adic-source-target}.
\end{theorem}

\begin{proof}
Lemma~\ref{lem:frag-presenting-system-assembly} constructs the
geometric functors and their coefficient transformations independently
of the dual formula.  Lemma~\ref{lem:frag-point-trace-diagram-descent}
then proves that formula as an equivalence of functors, and its
square~\eqref{eq:frag-specified-trace-reduction-square} identifies
exactly their previously constructed coefficient maps with perfect-dual
reduction.  Thus the given output diagram is cartesian, with uniform
perfect bounds from the finite continuous Koszul complex.  The finite
root bars themselves need not be bounded.

Apply Lemma~\ref{lem:frag-point-perfect-recovery} to this output
system and to the continuous star system on $V_\bullet^\vee$.
Perfect $\mathbf Z_\ell$-duality has the standard finite reductions;
the proved square identifies their entire diagram with the geometric
one.  The reduction/limit adjunction therefore gives the stated
integral comparison with its original maps.  Rationalize only after
this recovery.  The root limit has been taken first throughout.
\end{proof}

\section{Proof of the theorem and examples}
\label{sec:frag-direct-images}

\begin{proof}[Proof of Theorem~\ref{thm:frag-node-complete-local-datum}]
Compactify
$X$ in $(\mathbf P^1_R)^2\times(\mathbf P^1_R)^e$ by the equation
\[
 X_0Y_0=\pi X_1Y_1,
 \qquad x=X_0/X_1,\quad y=Y_0/Y_1.
\]
Take the standard affine charts for the $z_i$.  This is proper,
and $Y$ lies in its displayed affine open; the closure of the generic
boundary is in the complementary infinity divisors and misses $Y$.
Near $Y$ the chart is $\mathbf N\to\mathbf N^2$,
$1\mapsto(1,1)$, with smooth disc directions.  It is fs, saturated,
vertical and log smooth, with generic dimension $N$ and relative
rank one.  The tube is the actual open annulus times open discs,
so the analytic supported domain of Theorem~\ref{thm:frag-S-local-closed} applies.

The source construction of
Corollary~\ref{cor:frag-node-full-Kummer-star} uses the actual
Kummer torsors $a^{\ell^n}=x$.  Point finite-action descent handles
all bounded derived coefficients and maps at each $r$, and
Theorem~\ref{thm:frag-F-adic-iwasawa} identifies the chosen adic
coefficient systems.  Proposition~\ref{prop:frag-point-source-iwasawa-admissibility}
proves that the Iwasawa computations apply to the entire point
source quantified in the theorem.
Corollary~\ref{cor:frag-point-actual-root-bridge} identifies actual
structural star with continuous invariants.  On the actual Kummer torsors, Proposition~
\ref{prop:frag-log-specialization-kummer-localization} localizes the
specialization unit, and Lemma~\ref{lem:frag-log-nearby-geometric-input}
proves its invertibility from the constant-coefficient theorem.
The prime-to-$p$ log-Kummer source supplies the required tame descent.
Theorem~\ref{thm:frag-vertical-log-star} then gives actual nearby
and tube star by its specialization/completion maps.
Lemma~\ref{lem:frag-B-local-mapping-unit} proves that the induced
mapping map, including the ordinary-etale/quasi-etale bridge, is
an equivalence at finite precision; Proposition~\ref{prop:frag-local-adic-full-faithfulness}
passes that same map to the specified adic categories.

The geometric nearby and analytic trace dualities of Proposition~\ref{prop:frag-S-nearby-duality} and Proposition~\ref{prop:frag-S-analytic-duality}
apply to this compactification and this actual lisse realization.
Dualizing the same star map therefore constructs the actual
compact-support/costalk comparison in Theorem~\ref{thm:frag-S-local-closed}.
Proposition~\ref{prop:frag-point-finite-domain-verified} supplies
the published finite stack domain and its ambient module duality.
Theorem~\ref{thm:frag-point-C1} consequently identifies the genuine inverse trace tower
with the dual of the same continuous star.  This proves the second
row, including its single source twist and shift.  Its module
actions are those induced by geometric evaluation and trace.

At each fixed $r$ use the length-one continuous Koszul construction,
only after the root-bar passage.  It and its perfect linear dual
have the displayed uniform Tor-amplitudes, and the previously
constructed geometric coefficient-change maps are cartesian by
Theorem~\ref{thm:frag-theorem-B-local-closed}, Theorem~\ref{thm:frag-S-local-closed} and Theorem~\ref{thm:frag-point-C1}.
Lemma~\ref{lem:frag-point-perfect-recovery} now recovers actual
perfect $\mathbf Z_\ell$-outputs and their canonical reductions.
Lemma~\ref{lem:frag-adic-completion} identifies the canonical reductions
through its pro-zero Tor tower.  Completing the comparison maps and
then rationalizing proves the remaining assertions.
\end{proof}


\begin{corollary}[The geometric unipotent interface for refinement]
\label{cor:frag-refinement-unipotent-interface}
Set $e=0$ in Theorem~\ref{thm:frag-node-complete-local-datum}, put
$Q=\mathbf Q_\ell$, and retain its coordinate and compatible base roots.
Let $\mathcal U_\Gamma$ be the full subcategory of the rational
coefficient source with bounded finite-dimensional unipotent cohomology
representations (successive quotients have trivial $\Gamma$-action).
Let $\mathcal U^b(\mathbf G_{m,\bar K})$ denote the actual geometric
constructible complexes with unipotent lisse cohomology.  The associated
Kummer sheaf functor and continuous cochains give equivalences
\[
 \mathcal U^b(\mathbf G_{m,\bar K})\simeq\mathcal U_\Gamma
 \xrightarrow[\ K_\Gamma\ ]{\sim}\operatorname{Perf}(A_\Gamma),
 \qquad A_\Gamma=K_\Gamma(Q).
\]
Here $A_\Gamma$ has its continuous-cochain $E_\infty$ structure;
modules use the right-module convention.  The tensor on the last
category is $\otimes_{A_\Gamma}$, corresponding to the diagonal
coefficient tensor over $Q$.

For the actual associated coefficient $F$, the star comparison is
an $A_\Gamma$-module comparison, and its supported companion is
\[
 R\Gamma(Y,i^!R\Psi^t_{\overline X}Rj_{\eta,!}F)
 \simeq D_QK_\Gamma(V^\vee)(-1)[-2],
 \qquad D_Q=R\operatorname{Hom}_Q(-,Q).
\]
The action on the right is dual to the cochain action on
$K_\Gamma(V^\vee)$; the identification uses the actual evaluation
and trace pairing of the theorem.
\end{corollary}
\begin{proof}
In either coefficient category, unipotent sheaves are obtained from
constant finite-dimensional sheaves by their actual extension triangles.
Finite Postnikov towers then show that every displayed object is in
the thick closure of the unit.  Conversely kernels, cokernels,
extensions and summands of unipotent representations remain unipotent,
so this thick closure has exactly the stated cohomology condition.
The associated-sheaf functor on the torsors $z_n^{\ell^n}=x$ is exact
and tensor.  Its map on unit endomorphisms is the actual Kummer map
\[
 C_{\mathrm{cts}}(\Gamma,Q)\longrightarrow
 R\Gamma(\mathbf G_{m,\bar K},Q).
\]
It is an equivalence: both sides have $Q$ in degree zero,
$Q(-1)$ in degree one and no other cohomology, and the map sends
the universal Kummer class to that of $x$.  The geometric torus
calculation is the rank-one Kummer calculation; extension of the
algebraically closed field to $C$ preserves it.  Restriction to the
annulus gives the same class and the actual star map already proved
in the theorem.  Exactness and retracts extend the unit mapping
comparison in each variable to full faithfulness; unit generation
then gives essential surjectivity onto the actual geometric category.
This argument uses derived extension maps, not just representations
on the cohomology of an object.

On $\mathcal U_\Gamma=\operatorname{thick}(Q)$, cochains are
$R\operatorname{Hom}(Q,-)$.  The unit mapping comparison and finite
cones and retracts identify this functor with the one-object Morita
equivalence onto $\operatorname{Perf}(A_\Gamma)$.
The lax tensor structure supplies the balanced natural map
\[
 K_\Gamma(V)\otimes_{A_\Gamma}K_\Gamma(W)
       \longrightarrow K_\Gamma(V\otimes_Q W).
\]
It is the unit identification on $(Q,Q)$ and hence an equivalence
on both thick closures, by exactness in each variable.  Its symmetry
and associativity are inherited from this same cochain transformation.
The supported assertion and its module action are the theorem's
trace comparison, restricted to this source and rationalized.
\end{proof}

\begin{remark}[Shared coordinates and the boundary of the interface]
\label{rem:frag-refinement-coordinate-interface}
Write $t=\pi$, and fix $t_n\in\bar K$ with
$t_n^{\ell^n}=t$ and $t_{n+1}^{\ell}=t_n$.
On any generic overlap with $x=t^a u$, the maps of actual torsors are
\[
 z_n^{\ell^n}=x\ \longmapsto\ w_n=z_n/t_n^a,
 \qquad w_n^{\ell^n}=u.
\]
They commute with root transition and deck action.  For $v=u^{-1}$
the torsor map is $w_n\mapsto w_n^{-1}$ and the deck map is
$\zeta\mapsto\zeta^{-1}$.  Thus translation by $t^a$ identifies the
relative action, while inversion pulls it back along
$\gamma\mapsto\gamma^{-1}$.  Applying the associated-sheaf functor
and the bar cochain functor gives the specified maps of cochain
algebras and their coefficient modules, including all composition
maps: division by $t_n^{a+b}$ is successive division by $t_n^a,t_n^b$,
and inversion squares to the identity.  The familiar signs on
$H^1$ are consequences of these maps.

Evaluation commutes with these torsor identifications and group
automorphisms.  The induced finite quotient equivalences preserve
adjunction counits; hence the finite-carrier trace pairings and their
coefficient reductions commute as well.  On a geometric overlap,
analytic trace has the same compatibility by functoriality under
isomorphism.  Dualizing the same star map therefore transports the
supported comparison, with its Tate twist, rather than choosing a
new comparison from its cohomology.  All these assertions use the
fixed compatible base roots; no base-Galois-equivariant choice is
asserted.

This is the coefficient interface with the refinement paper: on a
thickness-one point chart its actual star algebra and coefficient
module are $A_\Gamma$ and $K_\Gamma(V)$, and its supported direct image
is the displayed costalk.  For a general stratum $S$, identifying
$H_S$, $N_S(F)$ and the direct image of $\kappa_S^!R\Psi F$ with
these local data still uses that paper's actual specialization,
restriction and gluing maps.  In particular, the coordinate isomorphisms
above compare coefficients on generic overlaps; cross-stratum links
and Gysin maps must come from the actual open/closed geometry.
They are not supplied by root-chart functoriality or lattice projections.
Neither a thickness-$m$ node nor a
positive-dimensional exceptional support is covered by this
corollary.  The Morita assertion concerns only the geometric
unipotent subsource, not the entire source of the node theorem.
\end{remark}

\begin{remark}[A nontrivial character and its integral torsion]
\label{rem:frag-node-nontrivial-character-test}
In the node theorem choose a generator $\gamma$ of relative inertia
and a rank-one coefficient $V=\mathbf Z_\ell$ on which $\gamma$ acts
by $u=1+\ell$ for odd $\ell$, or $u=1+4$ for $\ell=2$.
These give continuous characters of $\mathbf Z_\ell$; their reductions
factor through finite root quotients.  Put $a=u-1$.  The star complex
is $[\mathbf Z_\ell\xrightarrow{a}\mathbf Z_\ell]$ in degrees $0,1$.
It has only $H^1=\mathbf Z_\ell/a$.  The dual character has differential
$u^{-1}-1=-a/u$, so the supported formula has only
\[
 H^{2N}_c(T,\mathcal R(V))=(\mathbf Z_\ell/a)(-N).
\]
Its perfect representative nevertheless occupies degrees $2N-1,2N$.
After derived reduction modulo $\ell$, both of those degrees have
nonzero cohomology, agreeing with the constant mod-$\ell$ coefficient.
The extra degree comes from Tor, not an ordinary reduction of the
integral cohomology group.  Rationalization kills both star and
supported complexes for this nontrivial character.  \end{remark}


The coefficient source matters even for elementary proper images.
The next example shows why the local full-faithfulness statement
cannot be extended to all root-side morphisms.

\begin{proposition}[Proper-image presentations cannot lift all root morphisms]
\label{prop:frag-B-proper-morphism-obstruction}
Proper-image coefficients need not admit arbitrary
root-morphism lifting to generic or analytic coefficient maps,
even for one fixed strict proper map.  In particular the comparison
of star outputs cannot be promoted to a fully faithful identification
of the full realized analytic coefficient category with its root
outputs while retaining all these presentations.
\end{proposition}

\begin{proof}
Let $X=\mathbf A^1_R$, with only the vertical log structure, and set
\[
 X'=\operatorname{Spec}R\amalg\operatorname{Spec}R,
 \qquad h=(t=0)\amalg(t=\pi):X'\longrightarrow X.
\]
Each component is a closed immersion.  Their disjoint union is a
finite, strict proper log morphism; $X'$ is fs saturated log smooth
over the log trait.  Take $Y=X_s=\mathbf A^1_k$.
Its relative inertia rank is zero.  Both components of $Y'$ map
to $0\in Y$.

The two log-lisse perfect inputs $(R_r,0)$ and $(0,R_r)$ have root
outputs
\[
 E_0=i_{0,s,*}R_r=E_\pi
\]
and actual generic outputs
\[
 F_0=i_{0,\bar\eta,*}R_r,\qquad
 F_\pi=i_{\pi,\bar\eta,*}R_r.
\]
Their analytifications on the closed unit-disc tube have disjoint
closed point supports, since $0\ne\pi$ in $\bar K$.  For disjoint
closed immersions $i_a,i_b$, exceptional restriction satisfies
$i_a^!i_{b,*}=0$; adjunction therefore gives
\[
 R\operatorname{Hom}_{T_{\acute et}}(F_0^{\mathrm{an}},
                                      F_\pi^{\mathrm{an}})=0.
\]
The same vanishing holds on the algebraic generic fibre.  On the
root side, the identity of the common nonzero coefficient gives
\[
 \operatorname{Hom}_{Y_{\acute et}}(E_0,E_\pi)\supseteq R_r\ne0.
\]
Hence that identity cannot lift to a generic or analytic morphism
between the two prescribed realizations.  Already at degree zero
full faithfulness and arbitrary morphism lifting are impossible.
The same obstruction holds for the cartesian integral systems;
it persists after their $\mathbf Q_\ell$-realization.  No individual
$R_r$-coefficient is rationalized.

Proper base change still computes the star images in this example;
the loss occurs in passing from the coefficient to its star image.
\end{proof}

\appendix
% Supporting arguments for three_realization_submission.tex.
\section{General charts and coefficient conventions}
\label{sec:frag-chart-coeff}

\begin{assumption}[Ambient six-functor domain]
\label{ass:frag-ambient-six-functor-domain}
All algebraic models, their finite root stacks, and their quotient stacks are
locally of finite type over one fixed affine excellent noetherian scheme
$B$ of finite Krull dimension.  The prime $\ell$ is invertible on $B$, and
finite-type $B$-schemes have uniformly finite $\ell$-cohomological dimension
in each dimension occurring below.  A dualizing object on $B$ and the
constructible enhanced finite and adic six-functor formalisms are fixed once
and for all.  The strictly local face $S$, although used to split the chart,
is obtained by strict localization inside this domain; all statements over
$S$ and $\mathscr S=[S/\Delta]$ mean the corresponding base changes of the
fixed formalism.  In particular, every use of $Rf_*,Rf_!,f^*,f^!$, Verdier
duality, and a Beck--Chevalley mate below takes place in its stated
$D_c^+/D_c^-$ constructible domain.
\end{assumption}

\begin{assumption}[Finite trace domain and ambient duality compatibility]
\label{ass:frag-finite-trace-domain}
For each finite enhanced root map $g_n$ and each coefficient precision $r$,
use the genuine finite Artin-stack operations in the chosen enhanced
formalism.  The coefficient class and its source dual lie in the domain of
\[
Rg_{n,!}E_{n,r}\simeq
\mathbb D_{\mathscr S,r}Rg_{n,*}\mathbb D_{\mathrm{root},n,r}E_{n,r}.
\]
The formalism is required to have smooth purity, proper pushforward, and
nil-invariant coefficient pullback, with their usual composition and
coefficient-change compatibilities.  Lemma~\ref{lem:frag-central-refinement-purity}
deduces the central-root identity $q_{m,n}^!\simeq q_{m,n}^*$ from these
axioms.  It is not an additional root-specific purity conjecture.  This
concerns central faces; ramified maps of normal pairs have different purity formulas.

At each $r$ there is a fixed closed presentable stable category of sheaf
complexes on $\mathscr S$ in which the full-bar colimits and their dual limits
are formed.  Its internal Hom into $\omega_{\mathscr S,r}$ restricts to the
finite Verdier duals in the displayed formula.  No constructible bounded or
one-sided bounded category is assumed to contain arbitrary such limits.
The boundedness and adic recovery of the continuous outputs are proved from
the Iwasawa model.  For the split carriers over a geometric point used in the local trace
theorem, Proposition~\ref{prop:frag-point-finite-domain-verified}
verifies these requirements.  On a general base they remain explicit
additional hypotheses, and no general-base trace theorem is asserted here.
\end{assumption}

\begin{lemma}[Nil transport and central gerbe purity]
\label{lem:frag-central-refinement-purity}
Within the finite six-functor domain, let $i:X_0\hookrightarrow X$ be a
nilpotent closed immersion for which $i^*$ is a coefficient equivalence.
Then $i^!\simeq i^*$ canonically.  For the central refinements, with the
same carrier on both sides, $q_{m,n}^!\simeq q_{m,n}^*$; these identities
are compatible with composition and coefficient change wherever the
specified finite operations have coefficient-change isomorphisms.
\end{lemma}

\begin{proof}
Nil-invariance makes $Ri_*$ the inverse of $i^*$; representable
properness gives $Ri_!=Ri_*$.  Taking right adjoints therefore gives
$i^!=i^*$, compatibly with structural duality.
On reductions a central refinement is locally $BH\to S$ for a finite
etale group $H$.  This map is smooth of stack dimension zero, so
smooth purity gives $q_0^!=q_0^*$.  For reduction immersions $i_m,i_n$,
transitivity gives $i_m^!q^!=q_0^!i_n^!$.  Replacing the two $i^!$ by
$i^*$ and using their inverse equivalences yields $q^!=q^*$.
The carrier map is the identity.  Transitivity of purity and the
uniqueness of adjunction transport provide the composition cells;
coefficient compatibility holds wherever the finite operations have
coefficient-change isomorphisms.
\end{proof}


The preceding six-functor assumptions are standing foundational hypotheses,
rather than a preservation
condition hidden in the definition of a chart morphism.  It is the standard
base range of the Artin-stack six-operation results used here; the additional
continuous Iwasawa calculation is separated from the additional root-site
and geometric comparison theorems.

\begin{definition}[Standard split tame product chart over a strictly local base]
\label{def:frag-standard-product-chart}
Let $R$ be a complete strictly Henselian discrete valuation ring with
fraction field $K$, residue characteristic $p$, and fix $\ell\ne p$.  Let
$M\simeq\mathbf Z^d$ be the free cocharacter (relative characteristic)
lattice, write $M^\vee:=\operatorname{Hom}_{\mathbf Z}(M,\mathbf Z)$ for its
character lattice, and put
\[
  \Gamma_M:=M\otimes_{\mathbf Z}\mathbf Z_\ell(1),
  \qquad
  \Lambda_M^{\mathrm{Iw}}:=\mathbf Z_\ell[[\Gamma_M]]
  =\varprojlim_{U\subset\Gamma_M}\mathbf Z_\ell[\Gamma_M/U],
\]
A \emph{standard split tame product chart} of type $(M,e)$ is a
split log-smooth formal chart over the displayed strictly Henselian base which
is fs, saturated, algebraizable in a
neighbourhood of the chosen relative face, and topologically of finite type
over $\operatorname{Spf}R$, together with the following data.  The selected
tube is required to be separated and compactifiable over the face base:
it admits an open immersion into a proper finite-type chart over that base.
\begin{enumerate}[label=\textup{(\alph*)}]
\item The algebraization consists of a finite-type log-smooth morphism
  $f:X^{\mathrm{alg}}\to\operatorname{Spec}R$ in a neighbourhood of a
  specified locally closed stratum $i:Y\hookrightarrow X_s^{\mathrm{alg}}$,
  an open $j:U\hookrightarrow X_\eta^{\mathrm{alg}}$ whose specialization
  meets precisely that neighbourhood, and a proper log compactification
  $X^{\mathrm{alg}}\hookrightarrow\overline X^{\mathrm{alg}}$ over the face
  base.  Completion along $Y$ is the given formal chart, and the tube of $Y$
  has a fixed product
\[
  ]S[_{\mathrm{an}}
  \simeq \mathbb T_{M,K}^{\mathrm{an}}\times\mathbb D^e,
  \qquad \mathbb T_{M,K}=\operatorname{Spec}K[M^\vee],
\]
  compatible with specialization and with the face projection.  The symbols
  $i,j,U,Y,X^{\mathrm{alg}}$ and their Cech pullbacks always refer to these
  supplied objects; in particular the nearby open is not silently replaced by
  the whole generic fibre.  Here
$\mathbb D^e$ is a relatively open polydisc and the displayed product is an
actual chart isomorphism, not merely a homotopy equivalence.  The chart also
supplies the rank-$e$ smooth tangent bundle $\mathcal E_{\mathrm{sm}}$ on $S$
obtained by restricting the relative tangent bundle of the polydisc factor to
its zero section.  On the split chart it is $\mathcal O_S^e$; its descent
datum is part of the product chart.
\item We use the chosen pro-$\ell$ Kummer subsite: it is generated by the
finite covers
\[
   [\ell^n]:\mathbb T_{M,K}\longrightarrow\mathbb T_{M,K},
   \qquad m\longmapsto \ell^n m,
   \qquad n\geq 0,
\]
where the displayed multiplication is on the cocharacter lattice $M$
(equivalently, characters pull back by $\lambda\mapsto\ell^n\lambda$),
and by their pullbacks to the log chart.  Coefficients are required to be
trivial on the prime-to-$\ell$ tame quotient; the discarded finite
prime-to-$\ell$ part may be retained separately as finite descent data.  The
resulting deck group is $\Gamma_M=M\otimes\mathbf Z_\ell(1)$.  This is an
auxiliary pro-$\ell$ Kummer face site, not the ordinary \'{e}tale site of the
unrigidified quotient stack and not, by notation alone, a geometric Artin
fan.
\item Coefficients are taken in the constructible adic Kummer category of
Definition~\ref{def:frag-adic-coefficients}; in particular, every coefficient
used in a theorem below is represented by a compatible finite-level system
  in the selected cartesian action/module-sheaf model.
  Over a geometric point Proposition~\ref{prop:frag-point-source-iwasawa-admissibility}
  identifies the entire selected class with the classical perfect
  Iwasawa model.
\item The chosen compactification restricts to the fixed toric/polydisc
  compactification of the displayed tube.  Analytic compact support uses
  extension by zero followed by proper direct image.  Its nearby counterpart
  is the costalk construction of \eqref{eq:frag-nearby-star-shriek-definition};
  their comparison is proved on the supported domain specified below.
\end{enumerate}
We denote the pro-$\ell$ site in (b), together with its finite descent
groupoid, by $\mathcal F_M^{\mathrm{rel,ket}}$.
Throughout the paper, $I_t^{\mathrm{base}}$ means the selected pro-$\ell$
quotient $I_{t,\ell}^{\mathrm{base}}$; any retained prime-to-$\ell$ finite
tame factor is incorporated into the finite descent groupoid.
The finite tame descent datum may act on $M$ through a finite group
of lattice automorphisms.  The root coefficient constructions respect
that action, with quotient base specified in
Definition~\ref{def:frag-root-presented-chart}.  Geometric analytic
equivariance additionally uses the actual chart automorphisms and
a compatible compactification.
\end{definition}

\begin{remark}[Actual tubes and the separate log-smooth chart domain]
\label{rem:frag-actual-tube-domain}
The entire-torus product in Definition~\ref{def:frag-standard-product-chart}
is an additional restriction, not a consequence of log smoothness.
For example, in $X=\operatorname{Spec}R[x,y]/(xy-\pi)$, let $Y$ be
$x=y=0$ in the special fibre.  A point specializes to $Y$ precisely when
$|x|<1$ and $|y|<1$.  Since $xy=\pi$, its geometric tube is
\[
 ]Y[_{\bar\eta}=\{\,|\pi|<|x|<1\,\},
\]
an open annulus, not the entire torus in the supplied monomial coordinate.
Adding smooth coordinates vanishing on $Y$ adds open discs.  Also,
completion along $Y$ uses the ideal $(x,y)$; modulo $\pi$ its completed
ring contains $k[[x,y]]/(xy)$, rather than a finite-type $k$-algebra.
Such completions belong to the special formal-scheme domain and need
not be topologically of finite type for the $\pi$-adic topology.

The log-smooth star theorem below works on these actual completions.
It is stated in its own domain, Definition~\ref{def:frag-vertical-log-B-domain},
and does not change the compactification or support hypotheses of the
older product-chart statements.  In particular, a proof for an algebraic
torus cannot be transported to this annulus just by changing its name.
\end{remark}

\begin{definition}[Relative root chart]
\label{def:frag-root-presented-chart}
On a strictly local face base $S$, fix a torsion-free fs sharp monoid
$P$ and a primitive $\rho\in P$ with
$L_{\mathrm{rel}}=P^{\mathrm{gp}}/\mathbf Z\rho$ free.  Put
$M=\operatorname{Hom}(L_{\mathrm{rel}},\mathbf Z)$.
For a nonclosed face use its supplied pair $(P_\tau,\rho_\tau)$.
This root datum is independent of the entire-torus condition of
Definition~\ref{def:frag-standard-product-chart}.
A finite group $\Delta$ acts on $(P,\rho)$ and on $S$, with
\begin{equation}
 |\Delta|\in\mathbf Z_\ell^\times.
 \label{eq:frag-delta-prime-to-ell}
\end{equation}
The quotient base and its Iwasawa algebra are
\begin{equation}
 \mathscr S=[S/\Delta].
 \label{eq:frag-quotient-base}
\end{equation}
On $S$ the algebra is $\Lambda_M^{\mathrm{Iw}}$ with its semilinear
$\Delta$-action; when the base action is trivial one may use
$\Lambda_M^{\mathrm{Iw}}\rtimes\Delta$.

A base-preserving saturated chart map
$\alpha:(P_{\tau'},\rho_{\tau'})\to(P_\tau,\rho_\tau)$ preserves $\rho$
exactly and induces a relative lattice map with torsion-free cokernel.
It must also define a morphism of the supplied log points
\cite[\S1.2]{FujiwaraKatoNakayama}: writing
$\epsilon_Q:\mathcal O_S[Q]\to\mathcal O_S$ for the augmentation,
we require $\epsilon_{P_\tau}\circ\mathcal O_S[\alpha]
=\epsilon_{P_{\tau'}}$.  For the sharp zero-augmentation charts used
here this is precisely $\alpha^{-1}(0)=\{0\}$.
Preserving $\rho$ and saturation alone do not imply this condition.
The map, including its log-point structure, is $\Delta$-equivariant
when descent is present.  Its root maps
and the $2$-cells of the trivial-root sections come from the root
fibre's universal property, hence compose coherently.  Such a map
need not be strict.  Root refinement $P\to\ell^{-n}P$ is a separate
Kummer operation, with torsion cokernel.
Writing $\bar\alpha:L_{\mathrm{rel},\tau'}\to L_{\mathrm{rel},\tau}$,
precomposition gives
\begin{equation}
 \alpha_M:M_\tau\to M_{\tau'},\qquad
 \alpha_{\mathcal K,n}:\mathcal K_{\tau,n}\to\mathcal K_{\tau',n},\qquad
 \alpha_\Lambda:\Lambda_{M_\tau}^{\mathrm{Iw}}
                   \to\Lambda_{M_{\tau'}}^{\mathrm{Iw}}.
 \label{eq:frag-chart-variance}
\end{equation}
For rank-changing maps preservation of the selected perfect
coefficient class must be checked for the operation in question.

For the algebraized log-smooth chart $\mathfrak U$, the fixed carrier is
\begin{equation}
 \mathcal V_{M,e}=
 \mathcal T^{\log}_{\mathfrak U/\mathfrak s_{\mathbf N\rho}}|_S
 \simeq (M\otimes\mathcal O_S)\oplus\mathcal E_{\mathrm{sm}}.
 \label{eq:frag-log-smooth-carrier-bundle}
\end{equation}
It has rank $d+e$.  The bundle is intrinsic to this model, its splitting
uses product coordinates, and its $\Delta$-action is the displayed
diagonal action.  Strict log-etale changes preserve it.  Use
$\mathbf V(\mathcal V)=\operatorname{Spec}_S\operatorname{Sym}(\mathcal V^\vee)$
with compactification $\mathbf P(\mathcal V\oplus\mathcal O_S)$.
Its Thom trace supplies the dimension factor; refinement traces are
constructed in Theorem~\ref{thm:frag-point-C1} over a geometric point.

Write $\mathfrak s_{P/S}$ and $\mathfrak s_{\mathbf N\rho/S}$ for
$S$ with the indicated constant log charts, nonzero characteristic
elements mapping to zero.  The relative root construction uses their
fibre; it does not require $P/\mathbf N\rho$ to be fs or sharp.
\end{definition}



\begin{lemma}[Derived coefficient recovery]
\label{lem:frag-adic-completion}
Put $A=\mathbf Z_\ell$ and $R_s=A/\ell^s$.  Fix an ambient
presentable stable $A$-linear category of module complexes, with its
module categories for $R_s$ and their derived scalar changes.
For a derived-cartesian system
\[
 F_{s+1}\otimes^L_{R_{s+1}}R_s\simeq F_s,
\]
form $\widehat F=R\!\varprojlim_sF_s$ after restriction to that
specified ambient category.  Then the canonical map
\[
 \widehat F\otimes^L_A R_r\xrightarrow{\sim}F_r
\]
is an equivalence for every $r$, and $\widehat F$ is derived
$\ell$-complete.  No Mittag--Leffler hypothesis or exactness of
products of cohomology sheaves is needed.

For the selected adic sheaf categories in this paper, the adic
object is the cartesian system itself.  The displayed ambient
realization recovers its reductions; it is not asserted to be a
bounded locally free ordinary $A$-sheaf on the small etale site.
Integral bounded perfectness at a geometric point is proved in
Lemma~\ref{lem:frag-point-perfect-recovery}.  Rationalization is
performed on the specified integral/adic realization, never
termwise on the torsion coefficients.
\end{lemma}

\begin{proof}
Fix $r$ and discard the finitely many levels $s<r$.  In $D(R_s)$
the multiplication map has the canonical fibre sequence
\[
 R_r[1]\longrightarrow R_s\otimes^L_A R_r
                  \longrightarrow R_r.
\]
Indeed, the middle term is $[R_s\xrightarrow{\ell^r}R_s]$
in degrees $-1,0$; its kernel is $\ell^{s-r}R_s\simeq R_r$ and
its cokernel is $R_r$.  Under $s+1\to s$ the map on that kernel
is multiplication by $\ell$, since
$\ell^{s+1-r}=\ell\,\ell^{s-r}$.  The cokernel map is the identity.
This calculation includes the module structures and the canonical
maps to $R_r$, not just the two cohomology groups.

Tensor this fibre sequence over $R_s$ with $F_s$ and use
cartesianness.  It gives a fibre sequence of towers
\[
 (F_r[1],\ \ell)\longrightarrow
 (F_s\otimes^L_A R_r)_s\longrightarrow(F_r,\ \mathrm{id}).
\]
Compatibility follows by first restricting the ring-level sequence
at $s$ to $R_{s+1}$, and then tensoring with $F_{s+1}$; the prescribed
cartesian cells identify the result with the sequence at $s$.
The left-hand tower is pro-zero: $r$ successive maps multiply by
$\ell^r=0$ on the $R_r$-module $F_r[1]$.  More explicitly, in its
product formula for homotopy limit the shift operator $T$ satisfies
$T^r=0$, so $1-T$ is invertible with inverse
$1+T+\cdots+T^{r-1}$.  Its homotopy limit is zero.

The object $R_r$ is perfect over $A$: it is the cofiber of
$A\xrightarrow{\ell^r}A$.  Tensoring by it therefore commutes with
all limits in the stable ambient category.  Taking homotopy limits
of the preceding fibre sequence proves the asserted canonical
recovery.  The middle tower is not eventually constant; only its
extra Tor fibre is pro-zero.

For any object $K$, the compatible triangles
$K\xrightarrow{\ell^s}K\to K\otimes^L_A R_s$ give
\[
 R\operatorname{Hom}_A(A[1/\ell],K)\longrightarrow K
       \longrightarrow R\!\varprojlim_s(K\otimes^L_A R_s).
\]
Here the first term is the limit of the multiplication-by-$\ell$
tower of $K$.  It vanishes for every $R_s$-module by the same
nilpotent-shift argument.  Internal Hom in its second variable
preserves limits, so it also vanishes for $\widehat F$.
This proves derived completeness in the specified ambient category.
Nothing here exchanges geometric stalks with infinite products or
asserts that ordinary sheaf cohomology of $\widehat F$ is finite.
\end{proof}

\begin{lemma}[Perfect cartesian systems over a geometric point]
\label{lem:frag-point-perfect-recovery}
Let $(F_s)_s$ be a cartesian system in $D(R_s)$ with each $F_s$
perfect of a common Tor-amplitude $[a,b]$.  Then
$R\!\varprojlim_sF_s$ is perfect over $\mathbf Z_\ell$, with the
same Tor-amplitude, and has the given canonical reductions.
Its cohomology is recovered by the ordinary limits of $H^i(F_s)$.
These last cohomology towers are Mittag--Leffler as towers of finite
modules; no analogous general sheaf assertion is made.
\end{lemma}

\begin{proof}
Form $F=R\varprojlim_sF_s$ in $D(A)$, $A=\mathbf Z_\ell$,
using the given enhanced diagram.  Lemma~\ref{lem:frag-adic-completion}
already supplies its canonical reductions.  Since each $H^i(F_s)$
is finite, its tower is Mittag--Leffler.  The module Milnor sequence
therefore gives $H^i(F)=\varprojlim_sH^i(F_s)$ and $F\in D^{[a,b]}(A)$.

We prove perfectness by induction on the common Tor interval, without
choosing compatible chain representatives.  If $F_1=0$, the
minimal-complex test below gives $F_s=0$ for every $s$, and $F=0$.
Otherwise one may decrease $b$ to the actual top degree of $F_1$.
Equivalently, when $H^b(F_1)=0$ the following step uses the empty
basis, $m=0$, and decreases the interval without attaching a cell.
Choose a basis of $H^b(F_1)$.  The triangle $F\xrightarrow{\ell}F\to F_1$ and
$H^{b+1}(F)=0$ lift this basis to a map
$u:A^{\oplus m}[-b]\to F$.  Projecting $u$ to the given diagram
and using scalar adjunction gives a morphism of cartesian systems
$(R_s^{\oplus m}[-b])_s\to(F_s)_s$, with all its compatibility cells.
Let $(G_s)_s$ be its cofiber system.

Each $G_s$ is perfect.  Its reduction to $R_1$ is the cone of a map
which is an isomorphism in degree $b$ and has no other source
cohomology, so it has amplitude $[a,b-1]$.  Over the local Artin ring
$R_s$, a minimal finite free representative has term ranks given
by the cohomology dimensions of its derived residue-field reduction.
Consequently every $G_s$ has Tor-amplitude $[a,b-1]$.  This is a
separate finite-ring test at each $s$, not a simultaneous choice
of representatives.  By induction $G=R\varprojlim_sG_s$ is perfect
in that interval.  Exactness of the stable limit and
$R\varprojlim_sR_s=A$ give a triangle
$A^{\oplus m}[-b]\to F\to G$.  Thus $F$ is perfect in $[a,b]$.
The empty-interval case is zero: a perfect $R_s$-complex with zero
residue-field reduction is zero by the same minimal-complex test.

Canonicity concerns $F$, its maps and its reductions, not the cells
chosen to prove its perfectness.  Indeed scalar reduction
$c:D(A)\to\varprojlim_sD(R_s)$ is left adjoint to
$L(F_\bullet)=R\varprojlim_sF_s$, since
\[
 \operatorname{Map}(cM,F_\bullet)
 =\varprojlim_s\operatorname{Map}_{R_s}(M\otimes_A^LR_s,F_s)
 =\operatorname{Map}_A(M,LF_\bullet).
\]
The preceding recovery lemma identifies the counit $cL F_\bullet\to
F_\bullet$ with an equivalence.  On derived-complete $M$ the unit
is an equivalence as well.  These adjunction maps identify the
\emph{given} coherent systems, all morphisms and their higher cells.
No strictification of an infinite tower is required.
\end{proof}

\begin{definition}[Duality conventions]
\label{def:frag-continuous-duality}
Fix a base dualizing system $\omega_{\mathscr S,r}$.  For
$f:X\to\mathscr S$ set $\omega_{X,r}=f^!\omega_{\mathscr S,r}$,
$\omega_{X/\mathscr S,r}=f^!R_r$ and
$\mathbb D_{X,r}=R\mathcal Hom_{R_r}(-,\omega_{X,r})$.
Use $\mathbb D_{\mathscr S,r}$ for the base dual.
The finite six-functor domain gives
\begin{equation}
 \mathbb D_{Y,r}Rf_!K\simeq Rf_*\mathbb D_{X,r}K,
 \qquad K\in\mathcal D^-_c(X,R_r),
 \label{eq:frag-unbounded-finite-verdier}
\end{equation}
on the stated reflexive domains.  In particular a proper
nonrepresentable classifying-stack map still has distinct star and
shriek images.

For compatible selected systems, $\mathbb D_{X,\mathrm{cont}}$
and $\mathbb D_{\mathscr S,\mathrm{cont}}$ denote the derived limits
of these finite duals.  Perfect biduality applies to the coefficients;
for continuous root outputs, the finite Iwasawa complex proves
boundedness, cartesianness and biduality first.  The levelwise mate then gives
\begin{equation}
 \mathbb D_{\mathscr S,\mathrm{cont}}Rf_!
 \simeq Rf_*\mathbb D_{X,\mathrm{cont}}.
 \label{eq:frag-relative-verdier-mate}
\end{equation}
Coefficient recovery uses Lemmas~\ref{lem:frag-iwasawa-adic-base-change}
and~\ref{lem:frag-adic-completion} in their specified ambient categories.

The coefficient dual has contragredient action:
\begin{equation}
 V_{\mathrm{cont}}^\vee
 =R\operatorname{Hom}^{\mathrm{cts}}_{\mathbf Z_\ell}
       (V_{\mathrm{cont}},\mathbf Z_\ell),
 \qquad (\gamma\phi)(v)=\phi(\gamma^{-1}v).
 \label{eq:frag-continuous-contragredient}
\end{equation}
At a point this is classical continuous $\mathbf Z_\ell$-linear
Hom.  Over a base it denotes the cartesian system of the finite
$R_r$-coefficient duals, with its ambient limit used only for recovery;
it is not ordinary Hom from an unspecified integral etale sheaf.
The $\Delta$-action is likewise semilinear and contragredient.
Rational duals are obtained after inverting $\ell$.
For a perfect base coefficient put
\begin{equation}
 C^{\vee_0}=R\mathcal Hom_{\mathbf Z_\ell}(C,\mathbf Z_\ell).
 \label{eq:frag-base-coefficient-dual}
\end{equation}
Tensor--Hom adjunction and biduality give
\begin{equation}
 \mathbb D_{\mathscr S,\mathrm{cont}}
 (\mathbb D_{\mathscr S,\mathrm{cont}}\mathcal L\otimes C)
 \simeq\mathcal L\otimes C^{\vee_0}.
 \label{eq:frag-base-duality-cancellation}
\end{equation}
Thus coefficient duality and source Verdier duality are distinguished
by their dualizing objects.
\end{definition}

Fix a smooth projective toric compactification
$j_M:\mathbb T_{M,K}\hookrightarrow\overline{\mathbb T}_{M,K}$ and a
compactification of the polydisc factor for which the extension-by-zero
definition gives
\[
 R\Gamma_c(\mathbb D^e,\mathbf Q_\ell)
 \simeq \mathbf Q_\ell(-e)[-2e].
\]
These product calculations fix the normalization of the continuous dual.
The actual local compact-support domain and its analytic duality are
specified in Section~\ref{sec:frag-supported-realization}.  Finite-root
traces use the genuine adjunction counits in the finite-stack domain;
their point implementation is verified in
Proposition~\ref{prop:frag-point-finite-domain-verified}.

Choose a basis $\gamma_1,\ldots,\gamma_d$ of $\Gamma_M$ and put
$T_i=\gamma_i-1\in\Lambda_M^{\mathrm{Iw}}$.  The elements
$T_1,\ldots,T_d$ form a regular sequence.  Consequently the Iwasawa--Koszul
complex
\begin{equation}
 \mathbf K_M^{\mathrm{Iw}}
 :=\bigotimes_{i=1}^{d}
 \bigl[(\Lambda_M^{\mathrm{Iw}})^{-1}
       \xrightarrow{\,T_i\,}(\Lambda_M^{\mathrm{Iw}})^0\bigr]
 \longrightarrow \mathbf Z_\ell
 \label{eq:frag-iwasawa-koszul}
\end{equation}
is a finite free resolution of the augmentation module.  This is the only
Koszul resolution used in the paper.

For comparison, let $\mathbf B_M^{\mathrm{cts}}\to\mathbf Z_\ell$ be the
completed normalized bar resolution of the profinite group $\Gamma_M$.  At a
finite root level and finite coefficient precision we write
$\mathbf B_{M;n,r}$ for the full normalized bar resolution of the finite deck
group $\Gamma_M/\ell^n\Gamma_M$ over $\mathbf Z/\ell^r$.  In general it is
unbounded.  In particular,
\[
 [\,(\mathbf Z/\ell^r)[C_{\ell^n}]
   \xrightarrow{\gamma-1}(\mathbf Z/\ell^r)[C_{\ell^n}]\,]
\]
is not a resolution of the trivial module, and no finite-level
bar--Koszul quasi-isomorphism is asserted.
Indeed, for $n=r=1$ and one generator,
\((\mathbf Z/\ell)[C_\ell]\simeq
(\mathbf Z/\ell)[u]/(u^\ell)\) with $u=\gamma-1$, so
$u^{\ell-1}$ is a nonzero element of the kernel of multiplication by $u$.
This also agrees with the nonvanishing of
$H^q(C_\ell,\mathbf Z/\ell)$ in arbitrarily high degrees.

\begin{lemma}[Iwasawa base change, dual closure, and adic recovery]
\label{lem:frag-iwasawa-adic-base-change}
Put $\Lambda=\Lambda_M^{\mathrm{Iw}}$, $R_r=\mathbf Z/\ell^r$,
$\Lambda_r=\Lambda\otimes^{\mathbf L}_{\mathbf Z_\ell}R_r$, and
$\mathbf K_{M,r}^{\mathrm{Iw}}=
 \mathbf K_M^{\mathrm{Iw}}\otimes^{\mathbf L}_{\mathbf Z_\ell}R_r$.
At a point let $V_{\mathrm{cont}}$ be in the equivalent classes of
Proposition~\ref{prop:frag-point-source-iwasawa-admissibility}.
Over a noetherian split base use instead a selected cartesian
coefficient of Definition~\ref{def:frag-adic-iwasawa-sheaves};
$V_{\mathrm{cont}}\otimes R_r$ then denotes its given $r$th term.
In integral module-sheaf formulas use its ambient derived-limit
realization, whose reductions are those terms by
Lemma~\ref{lem:frag-adic-completion}; no global pseudo-coherent
representative is inferred.  Coefficient duals over a base are
formed levelwise in the adic model.  Set
\[
 H_{*,r}(V):=
 R\!\operatorname{Hom}_{\Lambda_r}
      (\mathbf K_{M,r}^{\mathrm{Iw}},V_{\mathrm{cont}}\otimes R_r).
\]
Then the following hold.
\begin{enumerate}[label=\textup{(\roman*)}]
\item $\mathbf K_{M,r}^{\mathrm{Iw}}$ is a length-$d$ finite free
  $\Lambda_r$-resolution of $R_r$, and finite-projective base change gives
  \[
  R\!\operatorname{Hom}_{\Lambda}
       (\mathbf K_M^{\mathrm{Iw}},V_{\mathrm{cont}})
       \otimes^{\mathbf L}R_r\xrightarrow{\sim}H_{*,r}(V).
  \]
  Hence $(H_{*,r}(V))_r$ is derived cartesian, uniformly bounded, and has
  finite constructible cohomology.
\item The continuous contragredient $V_{\mathrm{cont}}^\vee$ of
  \eqref{eq:frag-continuous-contragredient} remains in the selected
  cartesian category (and is classically Iwasawa-admissible at a point), and
  biduality is an equivalence.  The assertion is equivariant for $\Delta$ and
  for the involution $\gamma\mapsto\gamma^{-1}$ of $\Lambda$.
\item If $N=d+e$ and
  \[
  H_{!,r}(V):=\mathbb D_{\mathscr S,r}H_{*,r}(V^\vee)(-N)[-2N],
  \]
  then $(H_{!,r}(V))_r$ is also derived cartesian, uniformly bounded, and
  recovered from its derived inverse limit by reduction modulo $\ell^r$.
  Over a geometric point the resulting finite-module cohomology
towers are Mittag--Leffler by
Lemma~\ref{lem:frag-point-perfect-recovery}.
\end{enumerate}
These conclusions continue to hold after tensoring with a perfect
constructible object from $\mathscr S$.
\end{lemma}

\begin{proof}
After an \'{e}tale splitting of $M$, the variables $T_1,\ldots,T_d$ are a regular
sequence in $\Lambda_r=R_r[[T_1,\ldots,T_d]]$; hence the reduced Koszul complex
is still a finite free resolution of the augmentation.  Because its terms are
finite projective, derived Hom commutes with reduction without a
$R\!\varprojlim$ argument.  This proves (i), including the amplitude bound.

At each precision the selected action category is closed under
coefficient duality: the contragredient dual of an underlying-perfect
object is underlying-perfect, with reversed Tor interval.  Exact
finite-action realization and derived scalar change preserve its
evaluation and coevaluation.  Taking the cartesian system therefore
proves dual closure and biduality in the independent adic module
model, including semilinear descent.  At a point
Proposition~\ref{prop:frag-point-source-iwasawa-admissibility} gives
classical Iwasawa-admissibility of this dual.  We do not deduce
pseudo-coherence of a global Iwasawa sheaf merely from regularity
of its coefficient ring or from finite-generation at stalks.

The chosen dualizing system is derived cartesian in $r$.  Perfect duality and
the fixed purity factor consequently commute with reduction, proving the
cartesian square for $H_{!,r}$.  Lemma~\ref{lem:frag-adic-completion} now gives
derived recovery in the specified ambient model, without a sheaf
Mittag--Leffler argument.  At a geometric point,
Lemma~\ref{lem:frag-point-perfect-recovery} also proves perfect
integral recovery and the finite-module cohomology assertion.
This proves (iii) in the stated coefficient and duality domain.
\end{proof}


\section{Descent over a noetherian base}
\label{sec:frag-coefficient-descent}

The root tower uses finite quotient actions, whereas its continuous
limit uses an adic Iwasawa module.  We compare these coefficient
categories before discussing geometric direct images.  The finite
quotient may depend on the coefficient precision and on the finite
diagram being represented.

\subsection{Relative finite-action descent on a noetherian base}
\label{sec:frag-relative-action-descent}
This subsection treats the ordinary etale action-sheaf model on a
noetherian scheme $S$.  It is independent of the actual infinite-root
site of Theorem~\ref{thm:frag-actual-root-continuous-topos}.  The distinction between this derived action model and
an independently defined sheaf of pseudocompact Iwasawa modules is kept
until the latter comparison has been established.

\begin{definition}[The relative continuous action-sheaf category]
\label{def:frag-relative-action-sheaves}
Fix $r$ and let $\mathcal A_r=\operatorname{Mod}(S_{\acute et},R_r)$.
Let $\Gamma=\mathbf Z_\ell^d$, $H_n=\ell^n\Gamma$, and $K_n=\Gamma/H_n$.
Write $\mathcal A_{n,r}$ for sheaves in $\mathcal A_r$ with a $K_n$-action,
and $\mathcal A_{\Gamma,r}$ for sheaves with a smooth $\Gamma$-action:
\[
\operatorname*{colim}_n F^{H_n}\xrightarrow{\sim}F
\quad\text{as sheaves in }\mathcal A_r.
\]
Thus local sections are locally fixed by an open subgroup; continuity
is imposed on sheaves, not just on an arbitrarily chosen collection of
stalk representations.  These are Grothendieck abelian categories, with
exact forgetful functors.  Subscripts $c$ denote the full subcategories
whose underlying sheaves on $S$ are constructible.  The derived categories
$D_c^b(\mathcal A_{n,r})$ and $D_c^b(\mathcal A_{\Gamma,r})$ mean the full
subcategories of the ambient derived categories with bounded constructible
cohomology.  They are not defined as a colimit of finite-level categories.
\end{definition}

\begin{lemma}[noetherian action sheaves and a global finite quotient]
\label{lem:frag-noetherian-action-sheaves}
Every object of $\mathcal A_r$ is a union of constructible subobjects,
and its noetherian objects are exactly its constructible objects.
The same assertions hold in $\mathcal A_{n,r}$ and
$\mathcal A_{\Gamma,r}$, with constructibility measured after forgetting
the action.  In particular every object of
$\mathcal A_{\Gamma,r,c}$ has action through a single $K_n$ on all of $S$.
\end{lemma}

\begin{proof}
For a noetherian scheme, constructible $R_r$-sheaves form a strong Serre
subcategory and satisfy the ascending chain condition for arbitrary
subsheaves; every sheaf is a filtered colimit of constructible sheaves
\cite[Sections 59.73--59.74]{StacksConstructibleNoetherian}.
Taking images gives the asserted union.  Conversely, a noetherian object
in such a union equals one of its constructible subobjects.  Here the
hypothesis is that $S$ is noetherian, not merely that the coefficient
stalks are finite.

For a smooth action, $F^{H_n}$ is stable under $\Gamma$ since $H_n$ is
normal.  A constructible subobject $C\subset F^{H_n}$ need not be stable,
but the sum of its translates by the finite group $K_n$ is a
constructible, $\Gamma$-stable subobject of $F$.  Its action factors
through $K_n$.  Such subobjects cover $F$: first use smoothness, then
the constructible-subobject union inside $F^{H_n}$.  Their finite sums
form a directed union.  The analogous finite-group argument applies
to $\mathcal A_{n,r}$.

An equivariant object with constructible underlying sheaf is noetherian,
because equivariant subobjects are underlying subobjects.  Conversely a
noetherian equivariant object must equal one member of the union just
constructed, and is therefore underlying-constructible.  Finally, for
constructible $F$, the increasing chain of subsheaves $F^{H_n}$ is
stationary.  Its union is $F$, so $F^{H_n}=F$ at some level.  This proves
factorization on the entire base; no intersection of infinitely many
stalkwise stabilizers is used.
\end{proof}

\begin{lemma}[Filtered localization and finite diagram descent]
\label{lem:frag-filtered-enhanced-localization}
Suppose exact full inclusions of small abelian categories
$\mathcal B_n$ have filtered union $\mathcal B$, with a common exact
faithful underlying functor detecting cohomology.  Then
\[
 \operatorname*{colim}_nD^b(\mathcal B_n)\simeq D^b(\mathcal B)
\]
in stable infinity-categories with exact functors.  Mapping spectra
are filtered colimits over common stages, and finite simplicial
diagrams, including finitely specified homotopies, descend up to
equivalence to a common stage.  If an object condition is preserved
and detected by the unchanged underlying complex, the same result
holds for the associated full stable subcategories.
\end{lemma}
\begin{proof}
Use $N\operatorname{Ch}^b(\mathcal B_n)$ marked by all
quasi-isomorphisms $W_n$, and its infinity-categorical localization
$N\operatorname{Ch}^b(\mathcal B_n)[W_n^{-1}]$ as $D^b(\mathcal B_n)$.
Inflation here is the exact functor on the actual complexes before
localization.  A bounded complex, a finite string of chain maps,
and a chain homotopy involve finitely many objects and equations.
They all belong to a common stage.  Quasi-isomorphisms are detected
by the underlying functor, so the marked nerve for $\mathcal B$
is the filtered colimit of these marked nerves.

For every infinity-category $\mathcal E$, the universal property is
an equivalence of functor categories, including transformations:
\[
 \operatorname{Fun}(D^b(\mathcal B),\mathcal E)
 \simeq \operatorname{Fun}_{W}(N\operatorname{Ch}^b(\mathcal B),\mathcal E)
 \simeq \varprojlim_n
       \operatorname{Fun}_{W_n}(N\operatorname{Ch}^b(\mathcal B_n),\mathcal E).
\]
The subscript means the full subcategory sending the marked arrows
to equivalences.  Saturating those arrows does not alter this
condition: every such functor already inverts their saturation.
This proves localization commutes with this colimit, with no
extension-closedness assertion for a finite stage in the union.
The functors preserve shifts and cones, so the equivalence is stable.

For a filtered categorical colimit, objects come from finite stages
and mapping spaces are colimits of the mapping spaces at common
stages \cite[Lemma 0.2.1]{RozenblyumFilteredColimits}.  For the finite
diagram assertion one can use the explicit model in \cite[\S0.5]{RozenblyumFilteredColimits}:
a filtered diagram has a quasicategory model whose simplicial
colimit represents the categorical colimit.  A map from a finite
simplicial set factors through a stage, since it has finitely many
nondegenerate simplices and relations.  The same applies to a
specified homotopy (use its product with $\Delta^1$).  Thus this
assertion concerns the whole diagram, not just its vertices or
homotopy-category arrows.  Stability and compatibility with shifts
identify the corresponding mapping spectra as well.

Finally, represent an object of the selected target at a finite
stage.  Its underlying complex has not changed, so that representative
already satisfies the object condition.  Fullness means its mapping
spectra are the ambient ones just compared.  One takes a full
subcategory of the derived category, not complexes all of whose
terms individually satisfy perfectness.  Intermediate resolutions
and Postnikov pieces may therefore have nonperfect terms.
\end{proof}

\begin{lemma}[bounded noetherian derived replacement]
\label{lem:frag-noetherian-derived-replacement}
Let $\mathcal A$ be any of the abelian categories in
Lemma~\ref{lem:frag-noetherian-action-sheaves}, and let $\mathcal B$ be
its subcategory of noetherian objects.  The canonical exact functor is
an equivalence of stable enhanced categories
\[
D^b(\mathcal B)\xrightarrow{\sim}D^b_{\mathcal B}(\mathcal A).
\]
\end{lemma}

\begin{proof}
The noetherian objects form a Serre subcategory.  If $X\twoheadrightarrow Y$
with $Y$ noetherian, the directed union of noetherian subobjects of $X$
has images exhausting $Y$, so one already surjects.  This verifies
\cite[Lemma 13.17.4]{StacksDerivedSerre}.  Good truncation gives
bounded representatives.  The exact functor on derived enhancements
identifies Hom in every shift, hence mapping spectra; bounded
replacement gives essential surjectivity.
\end{proof}

\begin{theorem}[relative finite derived-action descent]
\label{thm:frag-relative-finite-action-descent}
On a noetherian scheme $S$, inflation induces an equivalence
\begin{equation}
 \operatorname*{colim}_n D_c^b(\mathcal A_{n,r})
 \xrightarrow{\sim}D_c^b(\mathcal A_{\Gamma,r}).
 \label{eq:frag-relative-derived-action-colimit}
\end{equation}
This equivalence restricts to the full subcategories with stratified
underlying-$R_r$-perfect complexes.  Its mapping spectra are the
filtered colimits of the finite-level mapping spectra.  Every finite
simplicial diagram with its specified higher cells has a finite-level
presentation.  The construction includes attaching maps across strata
and is natural for pullback between noetherian bases and for compatible
homomorphisms of the split inertia towers.
\end{theorem}

\begin{proof}
Write $\mathcal B_{n,r}$ and $\mathcal B_{\Gamma,r}$ for the noetherian
subcategories.  Lemma~\ref{lem:frag-noetherian-action-sheaves} proves
that each object of $\mathcal B_{\Gamma,r}$ is in some
$\mathcal B_{n,r}$.  A morphism between two such sheaves is equivariant
at their common level if and only if it is $\Gamma$-equivariant.
Consequently the continuous abelian category is the filtered union of
these finite-action abelian categories.  The finite categories are not
asserted to be extension-closed inside that union.

Apply Lemma~\ref{lem:frag-filtered-enhanced-localization} to this
filtered union, using the common underlying sheaf functor to detect
quasi-isomorphisms.  Lemma~\ref{lem:frag-noetherian-derived-replacement}
then identifies each localization with the corresponding full
constructible subcategory of the ambient action derived category.
This gives \eqref{eq:frag-relative-derived-action-colimit}, its mapping
spectra and finite higher-diagram descent.  The argument uses sheaf
complexes on the entire base, not bounded representatives chosen only
at stalks.  Inflation need not be fully faithful after deriving at a
fixed stage.

The underlying complex is unchanged by inflation.  Stratified
perfectness can therefore be imposed after the equivalence.  The
intermediate noetherian terms need not be $R_r$-perfect, so this
restriction does not discard torsion Postnikov pieces.  Because all
complexes and roofs above are sheaves on the entire base, the argument
already includes cross-stratum attaching morphisms.  Finally pullback
is exact, preserves constructible sheaves under the stated noetherian
base changes, and commutes with inflation and quasi-isomorphisms.
Its canonical localization squares give naturality; no particular
choice of a bounded representative is made into a functor.
\end{proof}

\begin{lemma}[Derived rectification of finite semilinear descent]
\label{lem:frag-finite-equivariant-rectification}
For the Grothendieck abelian action-sheaf categories $\mathcal A$ used
here, finite semilinear descent satisfies
$D(\mathcal A^\Delta)\simeq D(\mathcal A)^{h\Delta}$.
On the selected bounded perfect subcategories, finite-action descent
can consequently be performed in $\mathcal A^\Delta$ before deriving.
\end{lemma}
\begin{proof}
Forgetful evaluation has the exact left adjoint
$M\mapsto\bigoplus_{\delta\in\Delta}\delta^*M$.
On derived categories evaluation is conservative and preserves
colimits: it is exact and coproducts and filtered colimits are
computed on the underlying complexes.  Its monad is the same finite
sum with multiplication given by the group composition maps.
The derived monadic comparison therefore identifies
$D(\mathcal A^\Delta)$ with modules for this monad on $D(\mathcal A)$.
Evaluation from the homotopy fixed-point category has the same
left adjoint and monad, is conservative, and computes colimits on
the underlying objects.  Its monadic comparison gives the same
category.  This proves rectification; exact averaging alone would
not prove it.

For the finite root levels, work first with semilinear sheaves.
A finite union or sum of the $\Delta$-translates of a constructible
subobject is still constructible; $\Delta$ normalizes $H_n$.
The noetherian finite-action replacement and finite Postnikov
arguments therefore apply inside this abelian descent category.
They descend objects, attaching maps and finite diagrams to a
common root level.  The just proved rectification identifies this
with the displayed homotopy fixed-point notation.  The equivalence
preserves the underlying complex, so restricts to the selected
perfect class and commutes with its scalar reductions.  No exchange
of a filtered colimit with the infinite nerve $B\Delta$ is used.
\end{proof}

\begin{corollary}[finite descent and adic action systems]
\label{cor:frag-relative-adic-action-descent}
Theorem~\ref{thm:frag-relative-finite-action-descent} also holds for
finite semilinear descent by a group $\Delta$ acting on $S$ and
normalizing the subgroups $H_n$.  On the stratified underlying-perfect
subcategory its equivalences commute with derived coefficient
reduction.  In particular they induce an equivalence on the uniformly
bounded constructible cartesian adic systems:
\[
 \varprojlim_r\left(
   \operatorname*{colim}_n D_{c,\mathrm{perf}}^b(\mathcal A_{n,r})
                 \right)
 \simeq
 \varprojlim_r D_{c,\mathrm{perf}}^b(\mathcal A_{\Gamma,r}),
\]
with the same restrictions on the two limit categories.  For two such
systems represented at finite levels at each precision, the mapping
spectrum is
\[
 \varprojlim_r\operatorname*{colim}_{n\ge n_0(r)}
 \operatorname{Map}_{D(\mathcal A_{n,r})}(E_{n,r},F_{n,r}).
\]
No finite integral action level, and no reversed order of these limits,
is asserted.
\end{corollary}

\begin{proof}
Lemma~\ref{lem:frag-finite-equivariant-rectification} identifies the
resulting derived descent category with homotopy-coherent finite
descent.  Work in the abelian category of descent sheaves from the start.  Saturating a constructible subobject by the
finitely many $\Delta$-translates preserves constructibility.  The
$H_n$ are normalized by $\Delta$, so the invariant-subobject and
finite-saturation arguments still apply.  The resulting abelian
categories satisfy the same noetherian lifting property.  Apply the
proof of Theorem~\ref{thm:frag-relative-finite-action-descent} there.
This does not commute a colimit with the infinite nerve $B\Delta$;
it proves the equivariant assertion before deriving.

Inflation commutes with derived extension of coefficient scalars: both
retain the same underlying complex and action, and the canonical
extension-of-scalars comparison follows from tensor associativity.
On the stratified underlying-perfect class the reductions remain
bounded and constructible; no such claim is made for an arbitrary
bounded complex over the nonregular ring $R_r$.  The levelwise
comparison squares therefore form an equivalence of coefficient
systems of categories.  Taking their categorical limits proves the
adic statement.  The conditions on cohomological amplitude,
constructibility, and the prescribed adic recovery class are the
same after forgetting the action, so the equivalence restricts to them.
Mapping spectra in a limit category are limits of mapping spectra;
apply the fixed-$r$ equivalence to obtain the displayed formula.
Finite presentations and reduction cells may be chosen successively
with an increasing bound $\nu(r)$, using descent of each finite diagram.
This proves compatibility without positing one bound for the entire
adic system.
\end{proof}

\begin{proposition}[the two finite coefficient topologies]
\label{prop:frag-iwasawa-ideal-cofinality}
Put $\Lambda=\mathbf Z_\ell[[T_1,\ldots,T_d]]$, $I=(T_1,\ldots,T_d)$,
and identify $1+T_i$ with the chosen generators of $\Gamma$.  The ideals
\[
 J_{r,n}=(\ell^r,(1+T_1)^{\ell^n}-1,\ldots,
                      (1+T_d)^{\ell^n}-1)
 \quad\text{and}\quad Q_{r,a}=(\ell^r,I^a)
\]
define cofinal systems, already at each fixed $r$.  Consequently smooth
$\Gamma$-action sheaves of $R_r$-modules are exactly sheaves of
$\Lambda/\ell^r$-modules which are locally $I$-power torsion.
This is an exact equivalence of abelian categories, and hence of their
derived categories and bounded constructible-cohomology subcategories.
It concerns the derived category of torsion modules; comparison with
the unrestricted or completed module derived category is not asserted.
\end{proposition}

\begin{proof}
In $\Lambda/J_{r,n}=R_r[K_n]$, reduction modulo $\ell$ identifies the
augmentation algebra with
$\mathbf F_\ell[T_1,\ldots,T_d]/(T_1^{\ell^n},\ldots,T_d^{\ell^n})$.
Thus for $L=d(\ell^n-1)+1$ one has $I^L\subset\ell R_r[K_n]$,
and hence $I^{rL}=0$.  This proves $Q_{r,rL}\subset J_{r,n}$.
Conversely, for $1\le k<a$ the identity
\[
 k\binom{\ell^n}{k}=\ell^n\binom{\ell^n-1}{k-1}
\]
gives $v_\ell\binom{\ell^n}{k}\ge n-v_\ell(k)$.
Choose $n$ with $\ell^n\ge a$ and
$n\ge r+\lceil\log_\ell(\max\{1,a\})\rceil$.
All coefficients in degrees $1,\ldots,a-1$ are divisible by $\ell^r$;
the higher terms lie in $I^a$.  Therefore $J_{r,n}\subset Q_{r,a}$.
The case $d=0$ is the trivial augmentation case.

A section killed by $J_{r,n}$ has a finite $K_n$-action, and the first
inclusion makes it $I$-power torsion.  A section killed by $I^a$ has
an action through $K_n$ by the second inclusion.  These constructions
are compatible on overlapping annihilator bounds and with restriction
of sheaves.  They give inverse exact functors, without choosing a
uniform bound for all sections of an arbitrary sheaf.  The global
bound for constructible sheaves is supplied separately by
Lemma~\ref{lem:frag-noetherian-action-sheaves}.  Deriving an exact
abelian equivalence proves the last assertion.
\end{proof}

\begin{remark}[Why the completion model must be specified]
\label{rem:frag-small-etale-completion-test}
A cartesian adic system is not in general a locally free ordinary
$\mathbf Z_\ell$-sheaf on the small etale site.  For example let
$\eta=\Spec k((t))$, with $k$ algebraically closed and
$\operatorname{char}k\ne\ell$.  Choose a generator of the tame
$\mathbf Z_\ell$-quotient and let it act on $V=\mathbf Z_\ell$ by
$u=1+\ell^2$.  The reductions $V_r$ define nonzero rank-one lisse
$R_r$-sheaves $\mathcal L_r$ on $\eta$, with cartesian reduction maps.
Every finite extension has inertia image of finite index in this
quotient.  Its integral invariant module is zero since
$u^m-1\ne0$ in $\mathbf Z_\ell$ for every positive integer $m$.
Limits of sections commute with invariants, so the ordinary inverse
limit sheaf $\varprojlim_r\mathcal L_r$ is zero on the small etale site.
In particular $H^0(R\!\varprojlim_r\mathcal L_r)=0$ there.  This does
not assert that the derived inverse limit itself is zero; it rules
out interpreting it as the expected degree-zero locally free sheaf.
The example also occurs on the generic stratum of a strictly henselian
trait.  It does not obstruct an adic-system, pro-sheaf, or pro-etale
realization with a correctly specified completion functor.
\end{remark}

\begin{lemma}[Sections and torsion on the etale basis]
\label{lem:frag-etale-sections-torsion}
Let $S$ be noetherian, $A$ a noetherian ring, and
$\mathcal M=\operatorname{Mod}(S_{\acute et},\underline A)$.
On the basis of quasi-compact etale $U\to S$ the following hold:
sections commute with filtered colimits; evaluation $J\mapsto J(U)$
has an exact left adjoint; and, for injective $J\in\mathcal M$ and a
finitely generated ideal $I=(t_1,\ldots,t_d)$, the augmented stable
Cech map
\[
 \Gamma_IJ\longrightarrow
 \bigotimes_{i=1}^d[\underline A\to\underline A[1/t_i]]\otimes_AJ
\]
is a quasi-isomorphism.  Here $\Gamma_I$ means locally $I$-power
torsion, and the Cech complex starts in degree zero.
\end{lemma}
\begin{proof}
A quasi-compact etale scheme over $S$ is noetherian.  Its covering
families have finite subcovers, and the fibre products in their
matching conditions are again quasi-compact.  The finite-cover
sheaf condition is an equalizer of finite products.  Filtered
colimits of modules commute with these equalizers, so the presheaf
colimit restricted to this basis is already a sheaf.  Its value
at $U$ is therefore the colimit of the values at $U$.

For the etale localization $j_U:U_{\acute et}\to S_{\acute et}$,
put $L_U(M)=j_{U,!}\underline M_U$.  The constant-sheaf adjunction
and the localization adjunction give
\[
 \operatorname{Hom}_{\mathcal M}(L_U M,J)
       =\operatorname{Hom}_A(M,J(U)).
\]
At a geometric point $\bar s$ its stalk is
$\bigoplus_{\bar u\in U_{\bar s}}M$; this is exact in $M$.
Enough geometric points detect exactness, so $L_U$ is exact.
The left adjoint is extension along an etale localization, not a
structural proper or analytic shriek.  Applying this adjunction to
an injective $J$ shows that $J(U)$ is an injective $A$-module.

The annihilator of $I^a$ is a finite kernel, since $I^a$ is finitely
generated.  The first part gives
$(\Gamma_IJ)(U)=\Gamma_I(J(U))$.
Localization by $t_i$ is a filtered colimit of multiplication maps;
hence $J[1/t_i](U)=J(U)[1/t_i]$, also for multiple localizations.
The displayed sheaf Cech complex therefore has, on every basis
object, exactly the usual module Cech complex of $J(U)$.
For a noetherian ring its only cohomology on an injective module is
$\Gamma_I(J(U))$ in degree zero
\cite[Lemma 47.10.1]{StacksNoetherianTorsion}.
The augmented cone is consequently exact on every basis object.
Taking the filtered colimit over etale neighbourhoods makes its
geometric stalks acyclic, so the sheaf cone is acyclic.  This last
step does not require sections to be exact on arbitrary sheaves.
\end{proof}

\begin{proposition}[derived torsion embedding over the base]
\label{prop:frag-sheaf-derived-torsion}
Let $S$ be noetherian, $A=R_r[[T_1,\ldots,T_d]]$, and $I=(T_1,\ldots,T_d)$.
Let $\mathcal M=\operatorname{Mod}(S_{\acute et},\underline A)$ and let
$\mathcal T\subset\mathcal M$ be the sheaves locally killed by powers
of $I$.  Inclusion induces an equivalence
\[
D^+(\mathcal T)\xrightarrow{\sim}D^+_{I\text{-tor}}(\mathcal M),
\]
and hence is fully faithful on bounded constructible objects, with their
mapping spectra.  The target is a full subcategory of the unrestricted
module-sheaf derived category, not the derived category of a separately
chosen finite quotient ring.
\end{proposition}

\begin{proof}
Write $\Gamma_I F=\operatorname*{colim}_a\underline{\operatorname{Hom}}_A
(A/I^a,F)$.  It is the right adjoint to the exact inclusion of torsion
sheaves.  We verify its derived behavior on sheaves, rather than invoking
the corresponding module statement without a sheaf comparison.

Lemma~\ref{lem:frag-etale-sections-torsion} supplies the sheaf-level
calculation on injectives.  More explicitly, $\Gamma_I$ sends
injectives of $\mathcal M$ to injectives of $\mathcal T$, since its
left adjoint (inclusion) is exact.  On a bounded-below injective
resolution the finite-length flat Cech complex from that lemma
computes the underlying object of $R\Gamma_I$.  Its augmentation
is the inclusion counit, not just an identification of its cohomology.
Bounded-below totalization and the finite Cech length suffice here;
no unbounded torsion-completion equivalence is needed.

For a torsion sheaf all the displayed localizations vanish.  More
generally, for a bounded-below complex with torsion cohomology they
are acyclic, by flatness.  Thus the counit
$R\Gamma_I K\to K$ is an equivalence on precisely this part of the
argument.  For $K$ in $D^+(\mathcal T)$, apply this assertion to its included
complex.  The triangle identity identifies the included unit with
the inverse counit; exact faithful inclusion detects its being an
equivalence.  Thus both unit and counit of the derived adjunction
are equivalences on the stated subcategories, proving the claim.  Its canonical enhanced functor identifies
Hom in every shift, hence mapping spectra.  No unbounded sheaf
injective-completion assertion is needed here.
\end{proof}

\begin{definition}[Independent adic Iwasawa module-sheaf model]
\label{def:frag-adic-iwasawa-sheaves}
At precision $r$ let $\mathcal M_{r}^{\mathrm{sel}}$ be the full
subcategory of $D(\operatorname{Mod}(S_{\acute et},\underline{\Lambda_r}))$
with bounded $I$-power-torsion cohomology, constructible underlying
$R_r$-cohomology, and stratified underlying-$R_r$-perfectness.
Here $\Lambda_r=R_r[[T_1,\ldots,T_d]]$; modules and derived morphisms
are defined in the unrestricted module-sheaf category before selecting
this subcategory.  Define
\[
\mathcal M_{\mathrm{Iw,ad}}^{\mathrm{sel}}(S)
 =\varprojlim_{r,\,-\otimes^L_{R_{r+1}}R_r}\mathcal M_r^{\mathrm{sel}},
\]
restricted to a common cohomological and underlying-$R_r$ Tor
interval and a finite constructible stratification for the system.
These are all the selection conditions; no further finite-Tor
condition over an integral global Iwasawa sheaf is implicit.  Its mapping spectrum, denoted
$\operatorname{Map}_{\mathrm{Iw,ad}}$, is the mapping spectrum in this
limit category.  This is an independent algebraic definition: it uses
module sheaves over completed coefficient rings, with no group action
or root-site category in its construction.

This is the explicit relative completed coefficient model used in the
proved comparison below.  It does not define a global topologically
projective resolution theory on the small etale site.  The stronger
identification with any such separately proposed pseudocompact sheaf
category is a separate comparison problem and is not assumed implicitly.
\end{definition}

\begin{theorem}[Adic Iwasawa realization of continuous actions]
\label{thm:frag-F-adic-iwasawa}
On the noetherian split bases of
Theorem~\ref{thm:frag-relative-finite-action-descent}, including the
specified finite semilinear descent, there is a natural equivalence
\[
 \varprojlim_r D^b_{c,\mathrm{perf}}(\mathcal A_{\Gamma,r})
 \xrightarrow{\sim}\mathcal M_{\mathrm{Iw,ad}}^{\mathrm{sel}}(S).
\]
The same uniform conditions are imposed on both sides.  This equivalence,
together with relative finite-action descent, proves the coefficient equivalence for this
explicit target model.  At fixed precision it identifies relative
continuous invariants with
\[
R\underline{\operatorname{Hom}}_{\Lambda_r}(R_r,E_r),
\]
computed by the length-$d$ augmentation Koszul resolution.  This
invariants complex commutes with derived coefficient reduction on the
selected class.  Integral mapping spectra are the coefficient inverse
limits of the unrestricted $\Lambda_r$-module-sheaf mapping spectra.
\end{theorem}

\begin{proof}
Proposition~\ref{prop:frag-iwasawa-ideal-cofinality} identifies the smooth
action abelian category with $\mathcal T$.  Proposition~\ref{prop:frag-sheaf-derived-torsion}
identifies its bounded derived objects and their mapping spectra with
the full torsion-cohomology subcategory of unrestricted $\Lambda_r$
module sheaves.  Restrict to the underlying-perfect constructible
objects to obtain the equivalence at each $r$.  The functors preserve
the underlying $R_r$-complex and are compatible with scalar change;
on this selected class reduction remains bounded.  The canonical
squares for reduction consequently give an equivalence of the limit
categories, including their stated uniform restrictions.  Finite
semilinear descent is handled in the same abelian sheaf categories;
all the constructions, including torsion and scalar change, respect it.
This is a proof of the comparison with an independently defined target,
not its definition by transport.

Trivial relative action is restriction along $\Lambda_r\to R_r$.
Full faithfulness of the torsion embedding identifies its right
adjoint with the displayed internal derived Hom: test the adjunction
against arbitrary base complexes with trivial relative action in the
bounded-below domain.  The sequence $T_1,\ldots,T_d$ is regular in
$\Lambda_r$, even when $r>1$, so its finite free Koszul complex resolves
$R_r$.  Applying internal Hom constructs the sheaf-level comparison
and proves coefficient reduction and boundedness for these invariants.
Coordinates are used for the computation only; the right adjoint is
intrinsic.  Mapping spectra in the adic category are limits by its
universal property.  This last assertion is categorical and does not
interchange sheaf stalks with inverse limits.
\end{proof}

\begin{proposition}[The point source and Iwasawa perfectness]
\label{prop:frag-point-source-iwasawa-admissibility}
Let $A=\mathbf Z_\ell$, $\Gamma=\mathbf Z_\ell^d$ and
$\Lambda=A[[\Gamma]]$.  The entire category
$\mathcal S_{\mathrm{ad}}(\Gamma)$ of
Definition~\ref{def:frag-local-adic-source-target} is equivalent to
\[
 \{M\in\operatorname{Perf}(\Lambda):
                   M|_A\in\operatorname{Perf}(A)\}.
\]
The equivalence identifies the specified reductions with
$M\otimes_\Lambda^L\Lambda_r=M\otimes_A^LR_r$ and their mapping
spectra.  These objects have classical pseudocompact representatives
and are derived complete for $(\ell,T_1,\ldots,T_d)$.
Underlying-$A$-perfectness in the displayed class is essential:
it makes every reduction have finite, $I$-power-torsion cohomology.
For $d>0$, the free module $\Lambda$ alone is not in this class.
In particular no extra Iwasawa finite-Tor or pseudo-coherence
condition is needed in the definition of the point source.
\end{proposition}
\begin{proof}
\emph{Recovery with the action and the specified arrows.}
Write $\Lambda_r=\Lambda\otimes_A^LR_r$ and use the stable
localizations of module complexes.  Ideal cofinality and the
derived torsion embedding
(Propositions~\ref{prop:frag-iwasawa-ideal-cofinality}
and~\ref{prop:frag-sheaf-derived-torsion}, over a point) identify
bounded finite-cohomology continuous coefficients with the
corresponding full subcategory of $D(\Lambda_r)$.
These enhanced inclusions commute with derived coefficient change,
as in Theorem~\ref{thm:frag-F-adic-iwasawa}.  Thus the input,
including its specified scalar-change cells, gives an object
$V_\bullet$ of $\mathcal L=\varprojlim_rD(\Lambda_r)$.
Here the limit uses derived scalar extension, not restriction.

There is an adjunction
\[
 c:D(\Lambda)\rightleftarrows\mathcal L:L,\qquad
 cM=(M\otimes_\Lambda^L\Lambda_r)_r,\qquad
 L(V_\bullet)=R\varprojlim_r\operatorname{Res}_\Lambda V_r.
\]
Indeed the mapping spectrum into a cartesian section is
\[
 \operatorname{Map}_{\mathcal L}(cM,V_\bullet)
 =\varprojlim_r\operatorname{Map}_{\Lambda_r}
       (M\otimes_\Lambda^L\Lambda_r,V_r)
 =\operatorname{Map}_{\Lambda}(M,L V_\bullet).
\]
Set $M=L V_\bullet$.  Its $\Lambda$-action is therefore present
before any finiteness argument.  The counit has components
$\rho_r:M\otimes_\Lambda^L\Lambda_r\to V_r$, adjoint to the limit
projections.  The forgetful functor $D(\Lambda)\to D(A)$ preserves
limits, and scalar extension here is tensor by the central
$A$-perfect object $R_r$.  After forgetting, $\rho_r$ is exactly
the recovery map of Lemma~\ref{lem:frag-adic-completion} applied
to the given cartesian system.  That lemma applies directly in
$D(\Lambda)$ as well.  The forgetful functors are conservative,
so $\rho_r$ is an equivalence of $\Lambda_r$-objects.  Moreover
\[
 \rho_r=\theta_r\circ
   (\rho_{r+1}\otimes_{R_{r+1}}^LR_r)
\]
under scalar associativity, with the homotopy supplied by the
counit in $\mathcal L$; $\theta_r$ is the original reduction.
The same counit supplies the higher cells.  This proves recovery
of the given action diagram, not merely of its underlying values.
Lemma~\ref{lem:frag-point-perfect-recovery} now makes $M|_A$
perfect, and the recovery lemma makes $M$ derived $\ell$-complete.

\emph{Perfectness over the Iwasawa ring.}
Each $H^i(M)$ is finite over $A$.  Its finite $A$-generating set
also generates it over $\Lambda$, since $A\subset\Lambda$ and
the action has just been constructed.  A basis of $\Gamma$ identifies
$\Lambda$ with the regular local ring $A[[T_1,\ldots,T_d]]$ of
dimension $d+1$.  The finite-global-dimension theorem for regular
local rings makes this bounded finite-cohomology complex perfect.
A bounded finite free $\Lambda$-representative gives its classical
pseudocompact model and derived maximal-ideal completeness.
Its reductions have uniform $\Lambda_r$-Tor bounds.

\emph{Return to continuous and finite actions.}
Conversely, let $M$ belong to the displayed category.  The reductions
$N_r=M\otimes_A^LR_r$ are underlying-$R_r$-perfect with one Tor
interval, and their cohomology modules are finite sets.  On each
nonzero cohomology module the image of the local ring $\Lambda_r$
in its endomorphism ring is a finite local quotient.  Its maximal
ideal is nilpotent; thus the cohomology is $I$-power torsion.
This does not assert that $N_r$ itself is a complex over one finite
quotient merely because its cohomology has finite action.
Instead Proposition~\ref{prop:frag-sheaf-derived-torsion} identifies
$N_r$, through the counit $R\Gamma_I N_r\to N_r$, with a bounded
object of the derived torsion category.  Ideal cofinality identifies
that abelian category with discrete continuous $\Gamma$-modules.
Proposition~\ref{prop:frag-point-finite-action-descent} then descends
the whole derived object: successive Postnikov attaching maps, not
only its cohomology modules, descend to common finite quotients.
Full faithfulness of the enhanced torsion inclusion transports the
specified scalar arrows and cells of $N_\bullet$ back to the source.
Their boundedness and Tor bounds follow from the unchanged
underlying complexes.  This produces the required source system.

\emph{Maps and independence of coordinates.}
Every $M$ under consideration is derived $\ell$-complete, since
its underlying $A$-perfect complex is complete and forgetting is
conservative.  Thus both the unit of $c\dashv L$ on these objects
and its counit on the selected systems are equivalences.  In particular
\[
 \operatorname{Map}_\Lambda(M,N)
 =\varprojlim_r\operatorname{Map}_{\Lambda_r}
       (M\otimes_\Lambda^L\Lambda_r,N_r).
\]
This follows from completeness in the second variable and scalar
adjunction; no compactness of $M$ is required for this identity.
The finite-precision torsion equivalences identify it with the source
mapping spectrum.  All functors used in recovery are intrinsic to
$A[[\Gamma]]$.  A basis is used only to prove regularity; changing it
acts by an automorphism of this same ring.  Supplied semilinear actions
therefore preserve the recovery maps and their cells.  The statement
asserts no global pseudocompact sheaf comparison over a general base.
\end{proof}

\begin{proposition}[Agreement with classical pseudocompact Hom at a point]
\label{prop:frag-point-pc-Hom}
Over a geometric point, let $V,W$ be classical pseudocompact
$\Lambda$-complexes in the stated admissible class, in particular
$V$ is perfect over $\Lambda$ and $W$ is derived $\ell$-complete.
Their image in Definition~\ref{def:frag-adic-iwasawa-sheaves} satisfies
\[
R\operatorname{Hom}^{\mathrm{cts}}_\Lambda(V,W)
 \simeq R\!\varprojlim_r
 R\operatorname{Hom}_{\Lambda_r}(V\otimes^LR_r,W\otimes^LR_r).
\]
\end{proposition}

\begin{proof}
Choose a bounded finite-projective $\Lambda$-resolution of $V$.
Finite projective modules are topologically projective and every
$\Lambda$-linear map from them to a pseudocompact module is continuous.
Thus the same finite Hom complex computes classical continuous and
algebraic derived Hom.  Its reduction computes the right-hand term
at each $r$.  Finite projectivity commutes with the derived inverse
limit recovering $W$, proving the formula.  This identifies the new
relative model with the classical meaning on points; it does not
infer a global topologically projective sheaf resolution from that
point calculation.
\end{proof}

\begin{remark}[Proved adic model and the stronger pseudocompact interpretation]
\label{rem:frag-F-iwasawa-remainder}
The adic model on a noetherian split chart retains finite-action
descent, attaching maps and finite semilinear descent.  Its derived Hom is computed in unrestricted $\Lambda_r$
module sheaves before coefficient completion.  Over a point this agrees
with the classical pseudocompact calculation on the admissible class.

A global pseudocompact interpretation would identify this model with
some separately defined global topologically-projective pseudocompact
sheaf category.  That category and its realization are not supplied by
taking an ordinary inverse limit of small-etale sheaves.  It remains
outside the proved comparison, and statements explicitly using that
interpretation retain the separate global pseudocompact comparison.  The actual-root topos is identified separately in the next section.
\end{remark}


\section{The relative root site over a base}
\label{sec:frag-root-realization}

For the geometric root statements, a \emph{split} base includes a
chosen compatible identification of the relative finite etale
inertia tower with the constant groups $K_n=(\mathbf Z/\ell^n)^d$.
On the strictly local base of Definition~\ref{def:frag-root-presented-chart}
this follows by choosing compatible primitive roots of unity, since
$\ell$ is invertible.  The semilinear descent action records its
action on these choices.  On another noetherian base such a compatible
splitting must be supplied; it is not inferred for an arbitrary
finite etale group scheme.  The preceding abstract action-sheaf
results use a constant profinite group on any noetherian base.
The finite-quotient lemma below itself needs only a finite etale
group scheme and does not assume it constant.


We now identify the geometric source of continuous inertia.  The
relative root fibre is first constructed as an actual fpqc stack.
Finite-presentation descent then compares its representable etale site
with finite root levels; the resulting topos identifies the structural
star functor, not merely the values of selected cohomology groups.

\subsection{The infinite-root central fibre and its relative
\texorpdfstring{pro-$\ell$}{pro-l} category}
\label{frag:geometric-artin-face}

Talpo--Vistoli introduced the pullback of an infinite root stack along the
lifted closed point as its central fibre; it is a pro-algebraic stack and need
not itself be an infinite root stack
\cite{TalpoVistoliInfiniteRoot}.  The same construction is used in
Scherotzke--Sibilla--Talpo
\cite{ScherotzkeSibillaTalpoMcKay}.  The object below is its chartwise relative
form followed by the prescribed finite tame descent.  Its construction is
prior art.  What we isolate on its intrinsic full pro-root site is a
relative-inertia-selected perfect constructible category and the removal of
the base inertia direction.  It is not
the closed stratum of the ordinary Artin fan, which retains a diagonalizable
stabilizer in that direction.
Whenever a face $\tau$ is fixed in this subsection, we write
$P=P_\tau$ and $\rho=\rho_\tau$ (and likewise for the associated relative
lattice).  Thus all root-stack constructions below are understood facewise;
the suppressed face index may be restored along the supplied augmented
log-point maps of Lemma~\ref{lem:frag-root-face-functoriality}.
This convention constructs no specialization or incidence map between
strata.

For a fine saturated monoid $Q$ and an integer $a$ invertible on $S$, write
$\sqrt[a]{\mathfrak s_{Q/S}}$ for the $a$th root stack.  In the chart above
it has the standard presentation
\[
 \sqrt[a]{\mathfrak s_{Q/S}}
 =\left[
   \bigl(\operatorname{Spec}S[\tfrac1aQ]
        \times_{\operatorname{Spec}S[Q]}S\bigr)
   \mathbin{/}
   D\!\left((\tfrac1aQ)^{\mathrm{gp}}/Q^{\mathrm{gp}}\right)
  \right],
\]
where $D(A)=\underline{\operatorname{Hom}}(A,\mathbf G_m)$ and the map
$\operatorname{Spec}S\to\operatorname{Spec}S[Q]$ is the log-point
augmentation.  Thus every finite level is a genuine finite-type geometric
Artin stack; when $Q^{\mathrm{gp}}$ is free and $S$ is strictly local its
 reduced geometric fibre is $B\mu_a^{\operatorname{rk}Q}$.
For example, for $Q=\mathbf N$ over a field $k$ one obtains
\[
 \sqrt[a]{\mathfrak s_{\mathbf N/k}}
 \simeq[\operatorname{Spec}k[x]/(x^a)\,/\,\mu_a],
 \qquad
 (\sqrt[a]{\mathfrak s_{\mathbf N/k}})_{\mathrm{red}}
 \simeq B\mu_a.
\]
This elementary example is why all classifying-stack statements below are
made for the reduction or for the nil-invariant representable \'{e}tale topos,
never for the unreduced finite root stack itself.
For a possibly nonreduced base $S$, the notation $(-)_{\mathrm{red}}$ below
means the \emph{relative monomial reduction}: quotient by the nilpotent ideal
generated by positive fractional monomials while leaving $S$ unchanged.  Its
 classifying stack at order $\ell^n$ is therefore $B_S\mathcal K_n$.

Let $\mathbf{St}_{\mathrm{fpqc}}(S)$ denote the infinity-topos of
fpqc sheaves of spaces on $S$-schemes (in a fixed universe).
The root stacks are its $1$-truncated objects, equivalently stacks
in groupoids.  Limits of these remain $1$-truncated and compute
their $2$-categorical limits.  In the following display every inverse limit is formed in
$\mathbf{St}_{\mathrm{fpqc}}(S)$, not merely in the pro-category:
\[
 \sqrt[\infty]{\mathfrak s_{Q/S}}
 :=\varprojlim_{(a,p)=1}\sqrt[a]{\mathfrak s_{Q/S}},
 \qquad
 \sqrt[\ell^\infty]{\mathfrak s_{Q/S}}
 :=\varprojlim_{n\geq0}\sqrt[\ell^n]{\mathfrak s_{Q/S}}.
\]
The first is the prime-to-$p$ infinite root stack and the second is its
actual pro-$\ell$ limit stack.  We reserve
$\{\sqrt[\ell^n]{\mathfrak s_{Q/S}}\}_n$ for the associated pro-object.
The distinction matters: the representable finitely presented etale site of
the actual limit will be compared with the filtered finite-level site below,
and that comparison is a theorem rather than a definition.  The map
$\mathbf N\rho\to P$ induces maps at every root
level.  The zero (trivial-root) section, induced at each finite level by the
trivial line bundle with its zero section,
\[
 e_{\rho,S}:S\longrightarrow
 \sqrt[\infty]{\mathfrak s_{\mathbf N\rho/S}}
\]
is compatible with the finite-level sections $e_{\rho,S,n}$.  Define first the
split relative infinite-root central fibre over $S$ and then its descended
quotient over $\mathscr S=[S/\Delta]$:
\begin{equation}
 \widetilde{\mathcal A}_M^{\mathrm{rel}}
 :=
 e_{\rho,S}\times^{(2)}_{\sqrt[\infty]{\mathfrak s_{\mathbf N\rho/S}}}
 \sqrt[\infty]{\mathfrak s_{P/S}},
 \qquad
 \boxed{\mathcal A_M^{\mathrm{rel}}
 :=[\widetilde{\mathcal A}_M^{\mathrm{rel}}/\Delta]
 \longrightarrow\mathscr S.}
 \label{eq:frag-geometric-artin-face}
\end{equation}
For an arbitrary root order $a$ prime to $p$, reserve bracketed notation
\begin{equation}
 \widetilde{\mathcal A}_{M,[a]}^{\mathrm{rel}}
 :=e_{\rho,S,[a]}
   \times^{(2)}_{\sqrt[a]{\mathfrak s_{\mathbf N\rho/S}}}
   \sqrt[a]{\mathfrak s_{P/S}}.
 \label{eq:frag-arbitrary-root-order}
\end{equation}
Unbracketed $n$ below always denotes the exponent in the pro-$\ell$ root
order $\ell^n$; coefficient precision is denoted by $r$.  Thus the three
indices $a,n,r$ are never interchangeable.
Its pro-$\ell$ refinement is obtained by replacing every infinite root stack
in~\eqref{eq:frag-geometric-artin-face} with its $\ell^\infty$-part; denote
this restricted geometric object by
\begin{equation}
 \widetilde{\mathcal A}_{M,\ell}^{\mathrm{rel}}
 :=
 e_{\rho,S}\times^{(2)}_{\sqrt[\ell^\infty]{\mathfrak s_{\mathbf N\rho/S}}}
 \sqrt[\ell^\infty]{\mathfrak s_{P/S}},
 \qquad
 \mathcal A_{M,\ell}^{\mathrm{rel}}
 :=[\widetilde{\mathcal A}_{M,\ell}^{\mathrm{rel}}/\Delta]
 \longrightarrow\mathscr S.
 \label{eq:frag-proell-geometric-artin-face}
\end{equation}
The finite-level relative faces are
\[
 \widetilde{\mathcal A}_{M,n}^{\mathrm{rel}}
 :=
 e_{\rho,S,n}\times^{(2)}_{\sqrt[\ell^n]{\mathfrak s_{\mathbf N\rho/S}}}
 \sqrt[\ell^n]{\mathfrak s_{P/S}},
 \qquad
\mathcal A_{M,n}^{\mathrm{rel}}
 :=[\widetilde{\mathcal A}_{M,n}^{\mathrm{rel}}/\Delta]
 \longrightarrow\mathscr S.
\]

The bare central fibre records relative characteristic inertia but not the
ordinary smooth directions of the product chart.  It is therefore the correct
coefficient space for the root comparison, but it is not by itself the third geometric
object for compact supports.  Let
\[
 v:\mathbf V_{M,e}:=\mathbf V(\mathcal V_{M,e})\longrightarrow S,
 \qquad N:=d+e,
\]
and define the \emph{smooth-enhanced root faces}
\begin{equation}
 \widetilde{\mathcal R}_{M,e,n}^{\mathrm{rel}}
 :=\widetilde{\mathcal A}_{M,n}^{\mathrm{rel}}
       \times_S\mathbf V_{M,e},
 \qquad
 \mathcal R_{M,e,n}^{\mathrm{rel}}
 :=[\widetilde{\mathcal R}_{M,e,n}^{\mathrm{rel}}/\Delta],
 \qquad
 g_{M,e,n}:\mathcal R_{M,e,n}^{\mathrm{rel}}\to\mathscr S.
 \label{eq:frag-smooth-enhanced-root-face}
\end{equation}
In addition to the inverse system of these finite levels, define its actual
fpqc limit stack and structural map by
\begin{equation}
 \widetilde{\mathcal R}_{M,e,\ell}^{\mathrm{rel}}
 :=\widetilde{\mathcal A}_{M,\ell}^{\mathrm{rel}}
       \times_S\mathbf V_{M,e},
 \qquad
 \mathcal R_{M,e,\ell}^{\mathrm{rel}}
 :=[\widetilde{\mathcal R}_{M,e,\ell}^{\mathrm{rel}}/\Delta],
 \qquad
 g_{M,e,\ell}^{\mathrm{act}}:
 \mathcal R_{M,e,\ell}^{\mathrm{rel}}\to\mathscr S .
 \label{eq:frag-actual-smooth-enhanced-root-face}
\end{equation}
The notation
$\{\mathcal R_{M,e,n}^{\mathrm{rel}}\}_n$ is reserved for its finite-level
pro-object.  Coefficients on the actual limit are pulled back from the
central-root factor using Corollary~\ref{cor:frag-A0-adic}.  Homotopy invariance gives
$Rv_*v^*\simeq\mathrm{id}$ on these coefficients, whereas smooth purity gives
$Rv_!v^*\simeq(-N)[-2N]$.  Thus the carrier leaves the $*$-realization
unchanged and geometrizes, rather than inserts formally, the full
logarithmic--smooth compact-support shift.  In the stress-test case
$d=0,e=1$ and $P=\mathbf N\rho$, the object is
$\mathbf A^1_S\to S$; its constant-coefficient shriek image is correctly
$\mathbf Z/\ell^r(-1)[-2]$.

\begin{proposition}[Intrinsic tangent carrier on the fixed log model]
\label{prop:frag-intrinsic-tangent-carrier}
The vector bundle $\mathcal V_{M,e}$ and its zero section are independent of
the chosen splitting in
\eqref{eq:frag-log-smooth-carrier-bundle}.  Two strict log-etale product charts
on the same log model induce canonically identified carriers, projective-bundle
compactifications, Thom orientations, and finite-level $Rg_*$/$Rg_!$
operations.  Under the product decomposition, the determinant line is
\[
 \det(\mathcal V_{M,e})
 \simeq \det(M\otimes\mathcal O_S)\otimes\det(\mathcal E_{\mathrm{sm}}).
\]
\end{proposition}

\begin{proof}
For a log-smooth morphism the logarithmic cotangent complex is a locally free
sheaf in degree zero; its dual commutes with strict log-etale pullback.  The
relative characteristic exact sequence removes the base summand
$\mathbf Z\rho$ and identifies the remaining logarithmic tangent directions
with $M\otimes\mathcal O_S$.  The ordinary smooth tangent directions form the
quotient $\mathcal E_{\mathrm{sm}}$.  A product coordinate splits this exact
sequence but does not change its middle term.  Vector-bundle formation, the
projective completion $\mathbf P(\mathcal V_{M,e}\oplus\mathcal O_S)$, Thom
purity, and the six operations are functorial under the resulting bundle
isomorphism.  Taking determinants gives the displayed identity.
\end{proof}

\begin{remark}[Canonicity and locality of the carrier]
\label{rem:frag-carrier-locality}
Proposition~\ref{prop:frag-intrinsic-tangent-carrier} makes the carrier
intrinsic to the fixed log-smooth model and independent of a product
splitting.  It is not asserted to be intrinsic to the bare infinite-root
central fibre, nor is invariance under a logarithmic modification proved.
Comparison between genuinely different log models still requires an
isomorphism or a future modification-descent theorem.  Within the fixed
six-functor formalism, $Rg_!$ is independent of the chosen compactification,
but this standard independence does not imply log-model independence.
Theorem~\ref{thm:frag-node-complete-local-datum} uses this fixed-model, local,
canonicity.
\end{remark}

\begin{lemma}[Relative $2$-fibres are the finite-level inverse system]
\label{lem:frag-relative-fibre-limit}
The compatible finite-level root maps induce canonical equivalences of
groupoids, functorial in the test scheme over $S$ before descent,
\[
 \widetilde{\mathcal A}_{M,\ell}^{\mathrm{rel}}
 \ \simeq\ \varprojlim_{n\geq 0}
       \widetilde{\mathcal A}_{M,n}^{\mathrm{rel}},
\qquad
 \widetilde{\mathcal A}_{M}^{\mathrm{rel}}
 \ \simeq\ \varprojlim_{(a,p)=1}
                  \widetilde{\mathcal A}_{M,[a]}^{\mathrm{rel}}.
\]
In the first display, $n$ is the exponent in the root order $\ell^n$; in the
second, $a$ runs through all integers prime to $p$.  Taking the quotient
by $\Delta$ gives the corresponding inverse systems over $\mathscr S$.
Here the limits are limits of $2$-groupoids (equivalently, of the
corresponding fibred groupoids).  In particular, the relative $2$-fibre is not
being defined by first choosing the auxiliary Kummer site.
\end{lemma}

\begin{proof}
For a test scheme $S'\to S$, an object of the left-hand side is a compatible
system of roots on the $P$-chart over $S'$, together with a compatible
$2$-isomorphism between its restriction along $\mathbf N\rho\to P$ and the
trivial root section.  Giving this datum at the infinite root level is, by
definition of the inverse-limit root stack, the same as giving it at every
finite level with compatibility under divisibility of the root index.  The
same assertion holds for morphisms: a morphism and its $2$-isomorphism are
determined levelwise, and the compatibility equations are exactly those in
the inverse-limit groupoid.  Thus the two sides have equivalent mapping
groupoids for every $S'$, and the displayed equivalence follows from the
$2$-Yoneda lemma.

All transition maps are $\Delta$-equivariant.  Since the prescribed splitting
groupoid is finite and its descent is effective, taking the corresponding
descent groupoid commutes with this levelwise inverse limit.  This proves the
statement after the displayed quotient by $\Delta$ as well.
\end{proof}

\begin{lemma}[Torsor transitions in the relative quotient atlas]
\label{lem:frag-relative-atlas-torsors}
In the split chart, put $L=P^{\mathrm{gp}}/\mathbf Z\rho$ and
$H_n=\operatorname{Hom}(L,\mu_{\ell^n})$.  The relative finite root
fibre has a compatible presentation $[Z_n/H_n]$ with $Z_n$ affine.
For an affine test object $T$ carrying compatible maps to these fibres,
let $P_n=T\times_{[Z_n/H_n]}Z_n$.  Then $P_n$ is an $H_n$-torsor,
and $P_m\to P_n$ is a torsor under $\ker(H_m\to H_n)$ for $m\geq n$.
In particular these test-torsor transitions are finite faithfully flat.
This statement does not assert flatness of $Z_m\to Z_n$.
\end{lemma}

\begin{proof}
Work over an affine open $S=\operatorname{Spec}B$ of the split base.
For root order $a=\ell^n$, the absolute root atlas and the rigidified
relative atlas have coordinate rings
\[
 B[a^{-1}P]/(X^p:p\in P\setminus\{0\}),\qquad
 B_n=B[a^{-1}P]/(X^p:p\ne0,\ X^{\rho/a}),
\]
respectively: the second quotient is pullback along the zero section
of the base-root coordinate.  The base log point sends every positive
monomial to zero.  Choose nonzero generators $p_1,\ldots,p_t$ of $P$.
The ideal $J_n=(X^{p_1/a},\ldots,X^{p_t/a})$ in $B_n$ satisfies
$J_n^{t(a-1)+1}=0$, since each generator has $a$th power zero.
Sharpness gives $B_n/J_n=B$, even when $B$ itself is nonreduced.
This proves the required relative nilpotence with an explicit bound;
no uniform bound as $n$ varies is used.

The residual action is by the finite etale group $H_n$ below.
Its nerve has coordinate schemes
$U_{n,q}=\operatorname{Spec}B_n\times_S H_n^q$, affine and qcqs.
For $m\ge n$ the monoid map sends $X^{p/\ell^n}$ to
$(X^{p/\ell^m})^{\ell^{m-n}}$; together with $H_m\to H_n$ it
makes all nerve transitions affine.  These formulas commute with
localization on $S$, so the presentations and their properties
descend from affine opens.  They are the geometric input to the
finite-presentation limit theorem, not consequences of that theorem.

Write the absolute root chart as $[V_n/G_{P,n}]$ and the base root
chart as $[W_n/G_{\rho,n}]$, where
$G_{P,n}=\operatorname{Hom}(P^{\mathrm{gp}},\mu_{\ell^n})$ and
$G_{\rho,n}=\mu_{\ell^n}$.  Primitivity of $\rho$ and freeness of
$L$ imply that $G_{P,n}\to G_{\rho,n}$ is a finite faithfully flat
surjection with kernel $H_n$.  Pulling back along the trivial base-root
object is equivalent to choosing a trivialization of its associated
$G_{\rho,n}$-torsor.  The remaining structure group is $H_n$, and the
coordinate scheme is $Z_n=V_n\times_{W_n}S$, where $S\to W_n$ is
the zero-root section.  Thus the relative two-fibre is $[Z_n/H_n]$.
This argument uses the groupifications and an affine fibre product;
it does not require the fs monoid $P$ itself to be free.

Refinement gives an equivariant map $Z_m\to Z_n$ for the surjection
$H_m\to H_n$.  In the quotient-stack description a map from $T$
is a torsor with an equivariant map to $Z_n$.  Compatibility of the
test objects therefore gives the specified isomorphism
$P_m\times^{H_m}H_n\simeq P_n$.  After a faithfully flat cover
trivializing $P_m$, its map to $P_n$ is exactly $H_m\to H_n$.
The kernel is finite etale of rank $\ell^{(m-n)\operatorname{rk}L}$,
since $\ell$ is invertible.  The corresponding coordinate algebra is
locally free of this positive rank; flatness, faithful flatness and the
torsor identity descend.  This proves the assertion without any
flatness claim about the monoid coordinate maps.  The construction
uses the supplied refinement two-cells, so it respects composition.
For finite $\Delta$-descent the same assertion is checked after the
fpqc cover $S\to[S/\Delta]$; torsors and finite faithfully flat maps
descend, also when no $\Delta$-equivariant lattice splitting is chosen.
\end{proof}

\begin{proposition}[Effective atlas of the actual infinite-root fibre]
\label{prop:frag-actual-root-effective-atlas}
Let $U_{n,\bullet}\to
\widetilde{\mathcal A}_{M,n}^{\mathrm{rel}}$ be the compatible affine
action-groupoid presentations constructed from the two finite root charts,
and put
\[
 U_{\infty,q}:=\varprojlim_n U_{n,q}.
\]
Then $U_{\infty,0}$ is affine over $S$, the map
\[
 U_{\infty,0}\longrightarrow
 \widetilde{\mathcal A}_{M,\ell}^{\mathrm{rel}}
\]
is an effective fpqc epimorphism, and its Cech nerve is canonically
$U_{\infty,\bullet}$.  Consequently
\begin{equation}
 \widetilde{\mathcal A}_{M,\ell}^{\mathrm{rel}}
 \simeq |U_{\infty,\bullet}|
 \quad\text{in }\mathbf{St}_{\mathrm{fpqc}}(S),
 \label{eq:frag-actual-root-effective-presentation}
\end{equation}
and this is simultaneously the actual inverse-limit stack and the stack
presented by the limiting root groupoid.  The same statement holds after
effective finite $\Delta$-descent.
\end{proposition}

\begin{proof}
Test after an arbitrary affine map
$T\to\widetilde{\mathcal A}_{M,\ell}^{\mathrm{rel}}$.  Its projections give
compatible maps $T\to\widetilde{\mathcal A}_{M,n}^{\mathrm{rel}}$, and
\[
 P_n:=T\times_{\widetilde{\mathcal A}_{M,n}^{\mathrm{rel}}}U_{n,0}
\]
is the finite diagonalizable torsor of
Lemma~\ref{lem:frag-relative-atlas-torsors}, trivializing the level-$n$
relative root data.  For $m\geq n$, the transition $P_m\to P_n$ is the torsor under
the kernel of the finite diagonalizable quotient
$\mathcal H_m\to\mathcal H_n$ (with the base-root factor rigidified);
in particular it is finite faithfully flat.  This remains so after
\emph{any} base change.  In particular, for a geometric point
$t:\operatorname{Spec}\Omega\to T$, the maps
$\mathcal O(P_n\times_Tt)\to\mathcal O(P_m\times_Tt)$ are injective.
This is faithful-flat injectivity after base change, not an assertion
that tensoring an arbitrary injective ring map remains injective.  The inverse limit
$P_\infty=\varprojlim_nP_n$ is
$\operatorname{Spec}_T(\varinjlim_n\mathcal O_{P_n})$.  Each $P_n\to T$ is
finite locally free and faithfully flat.  Its filtered-colimit algebra is
flat.  Tensoring with $\Omega$ commutes with the filtered colimit;
the preceding fibrewise injectivity makes its nonzero unit survive.
Thus every geometric fibre is nonempty.  We use the elementary
criterion that a flat algebra $B$ over a ring $A$ is faithfully flat
if and only if $B/\mathfrak mB\ne0$ for every maximal ideal
$\mathfrak m$ of $A$.  Indeed, for $0\ne x\in M$ the cyclic
submodule $Ax=A/I$ embeds in $M$.  Choose $\mathfrak m\supset I$;
$B/IB$ surjects onto the nonzero $B/\mathfrak mB$, and flatness
injects $B/IB$ into $M\otimes_AB$.  Thus tensor detects nonzero
modules; the converse follows by testing residue fields.
Nonzero geometric fibres imply this residue-field condition by
faithfulness of field extension.  Apply the criterion to the
filtered-colimit algebra above.
Hence $P_\infty\to T$ is affine faithfully flat and quasi-compact.  Since
this holds after every affine test map, $U_{\infty,0}\to
\widetilde{\mathcal A}_{M,\ell}^{\mathrm{rel}}$ is fpqc.

Work in the sheaves-of-spaces infinity-topos specified above; its
$1$-truncated root objects compute the stack limits.  Limits commute
with limits in this ambient category.  Hence
\[
 \underbrace{U_{\infty,0}\times_{X_\infty}\cdots
                 \times_{X_\infty}U_{\infty,0}}_{q+1}
 \simeq\varprojlim_n
 \underbrace{U_{n,0}\times_{X_n}\cdots\times_{X_n}U_{n,0}}_{q+1}
 =U_{\infty,q}.
\]
Here $X_n$ denotes the relative root fibre at level $n$.
The equivalences preserve the projections and diagonals defining
all face and degeneracy maps.  They therefore identify the whole
Cech nerve, not just its objects degree by degree.  The representable fpqc map is locally surjective on objects, hence
epimorphic on $\pi_0$ in this sheaf infinity-topos.  Such a map is
an effective epimorphism: by the effective-epimorphism/Cech-descent
property of an infinity-topos, its Cech realization recovers the target.
On the other hand finite limits commute with inverse limits in the same
infinity-topos; applying this to the defining relative $2$-fibre identifies
that target with the actual stack
$\widetilde{\mathcal A}_{M,\ell}^{\mathrm{rel}}$ of
\eqref{eq:frag-geometric-artin-face}.  This proves
\eqref{eq:frag-actual-root-effective-presentation} and supplies fpqc effectivity for this presentation.  The comparison with
the representable etale site is proved in
Theorem~\ref{thm:frag-actual-root-fp-descent}; its continuous-action
interpretation and derived structural realization are proved separately in
Theorem~\ref{thm:frag-actual-root-continuous-topos} and
Corollary~\ref{cor:frag-actual-root-derived-realization}.
The atlas construction is equivariant for $\Delta$.  The quotient is
formed from its action groupoid and fpqc descent; it is not a finite limit
merely because $\Delta$ is finite.  Its etale coefficient comparison is
given by finite semilinear descent and the root-topos comparison below.
\end{proof}

\begin{definition}[representable etale objects over an fpqc stack]
\label{def:frag-actual-repet-sources}
For an fpqc stack $X$, an object of $\operatorname{RepEt}_{\mathrm{fp}}(X)$
is a morphism $Y\to X$ of fpqc stacks such that, for every scheme
$T\to X$, the pullback $Y\times_XT$ is an algebraic space etale and
of finite presentation over $T$.  The source $Y$ need not be algebraic.
Morphisms are morphisms over $X$, with the canonical identifications
in the slice; representability makes this an ordinary category up to
its canonical equivalence.  Coverings are jointly surjective etale
families.  At finite algebraic root levels this convention gives the
usual representable etale objects; at the actual limit it permits
their base changes, which may be non-algebraic fpqc stacks.
\end{definition}

\begin{lemma}[The scheme-source test at infinite inertia]
\label{lem:frag-repet-scheme-source-test}
Over an algebraically closed field of characteristic different from
$\ell$, let $G=\varprojlim_n K_n$ be an infinite proconstant group
scheme.  There is no nonempty algebraic space $Y$ admitting a
representable etale finitely presented morphism $Y\to BG$.
Thus restricting the source in Definition~\ref{def:frag-actual-repet-sources}
to schemes or algebraic spaces does not give a covering site for $BG$.
\end{lemma}

\begin{proof}
Pull back by the trivial torsor $\Spec k\to BG$.  The resulting
$P=Y\times_{BG}\Spec k$ is an etale finitely presented algebraic
space over $k$, hence a finite discrete set, and $P\to Y$ is a
$G$-torsor.  Nonemptiness of $Y$ implies nonemptiness of $P$.
The torsor action is free because $Y$ is an algebraic space; after
taking geometric points this would be a free action of the infinite
group $G(k)$ on a nonempty finite set.  This is impossible.  Sources
such as $BH\to BG$ for an open subgroup $H\subset G$ are instead
allowed by the stated fpqc-stack convention.
\end{proof}

\begin{lemma}[Representable etale objects on a finite quotient]
\label{lem:frag-finite-quotient-etale-descent}
Let a finite etale group scheme $G$ act on an affine scheme $U$.
For the fpqc quotient stack $X=[U/G]$, pullback to $U$ identifies
$\operatorname{RepEt}_{\mathrm{fp}}(X)$ with $G$-equivariant etale
finitely presented algebraic spaces $V\to U$.  Its inverse sends
$V$ to $[V/G]\to[U/G]$.  This morphism is representable even when
the action on $V$ has stabilizers.  Morphisms, their equalities,
and covering families are detected on the atlas $U$.
\end{lemma}
\begin{proof}
The action is an isomorphism between the two pullbacks of $V$ to
$G\times U$, with its unit and cocycle identities on $U$ and
$G\times G\times U$.  Form $[V/G]$ as the stack of $G$-torsors
with equivariant maps to $V$.  A test map $T\to[U/G]$ is a torsor
$P\to T$ with an equivariant map $P\to U$.  Directly from these
objects and arrows,
\[
 T\times_{[U/G]}[V/G]\simeq (P\times_U V)/G.
\]
The diagonal action here is free because $P$ is a torsor, regardless
of stabilizers on $V$.  Its quotient is the algebraic space obtained
by etale descent along the finite etale cover $P\to T$.
Explicitly, an etale scheme atlas of $P\times_UV$ and its translates
give an etale equivalence relation in schemes; the sheaf quotient
is that descended algebraic space.  Its pullback to $P$ is
$P\times_U V$, etale and finitely presented over $P$.  These
properties descend along this finite etale cover.  This proves
representability and the required finiteness over every $T$.
In particular its pullback to $U$ is exactly $V$.

Conversely, pulling back a representable object on $X$ gives this
action datum.  The stack descent equivalence reconstructs the object:
locally a lift consists of a map to $V$, with precisely the prescribed
action isomorphisms.  A map of such objects is a $G$-equivariant map
of the algebraic spaces; its compatibility is one equation on
$G\times U$.  Relative representability makes the mapping groupoids
discrete, so there are no additional independent higher $2$-cell
choices to descend.  A family covers if and only if its pullback
to $U$ is jointly surjective.  Composition and pullback of covers
are the corresponding operations on etale algebraic spaces.
\end{proof}

The following result concerns the specified relative pro-$\ell$ root
fibre over $S$.  Talpo--Vistoli's
\cite[Theorem 6.22, arXiv v2 numbering]{TalpoVistoliInfiniteRoot}
already gives the absolute Kummer-etale/infinite-root topos equivalence;
\cite[\S\S5--6]{CarchediEtAlLogHomotopy} compare the actual root stack,
its pro-system and its representable etale site after profinite shape.
Here we must additionally track the chosen base-root neutralization,
restriction to the $\ell$-power tower, and the structural coefficient
functor over $S$.  The next argument descends actual finite-presentation
objects; the later coefficient theorem handles precision $r$ and its
reduction arrows.  We use the effective chart atlas above and
descend its degree-$0,1,2$ etale data.  The conclusion below is a
site equivalence on this chart; its identification over $S$ and its
continuous derived structural functor are proved afterwards.

\begin{theorem}[finite-presentation descent on the actual root site]
\label{thm:frag-actual-root-fp-descent}
Let $X_n=\widetilde{\mathcal A}_{M,n}^{\mathrm{rel}}$ and
$X_\infty=\widetilde{\mathcal A}_{M,\ell}^{\mathrm{rel}}$ on the fixed
split chart.  Use their compatible affine action-groupoid presentations
$U_{n,\bullet}$ and the effective limit presentation of
Proposition~\ref{prop:frag-actual-root-effective-atlas}.  Explicitly,
over each affine open of $S$, $U_{n,q}$ is affine (hence qcqs),
its root transitions are affine,
and $U_{n,\bullet}$ is the nerve of a finite etale root-group action;
$\ell$ is invertible on the base.  These properties follow from the
explicit coordinate formulas in Lemma~\ref{lem:frag-relative-atlas-torsors}
and are the hypotheses used below.
Then pullback induces an equivalence of sites
\[
 \operatorname*{colim}_n\operatorname{RepEt}_{\mathrm{fp}}(X_n)
 \xrightarrow{\sim}\operatorname{RepEt}_{\mathrm{fp}}(X_\infty),
\]
with the source convention of Definition~\ref{def:frag-actual-repet-sources}
and the finite-cover topology on the left.  It is natural over $S$,
under the specified strict base changes, and under finite $\Delta$-descent.
In particular it proves the actual-site interpretation of the filtered
site.  Identification of this topos with continuous action sheaves and
its derived structural functor is identified in
Corollary~\ref{cor:frag-actual-root-derived-realization}.
\end{theorem}

\begin{proof}
Work first over an affine open of $S$.  There all $U_{n,q}$ are
affine and their transitions are affine.  The resulting site
comparisons commute with restriction to affine overlaps and hence
descend to $S$ by the same etale descent of objects and maps.  These
conditions come from the explicit monoid-algebra and finite-group
presentations.  The proof uses the algebraic-space limit theorem for
finite-presentation objects and morphisms, and descent of etaleness
and surjectivity after increasing the index
\cite{StacksSpaceLimitDescent}.  It does not use a limit theorem for
stacks with nonrepresentable root transition maps.

The descent input in addition to the algebraic-space limit theorem
is the following.  Relative representability makes an etale object
over $X_\infty$ a $0$-truncated object in the fpqc slice topos,
although its total stack may be $1$-truncated.  On the effective
atlas it is therefore determined by an algebraic space $V$ over
$U_{\infty,0}$, an isomorphism between its two pullbacks to
$U_{\infty,1}$, and its unit and cocycle equations on degrees zero
and two.  Higher simplices give the iterated composites of that
isomorphism; they are not further freely chosen descent cells.
Morphisms are maps over degree zero satisfying the degree-one
equivariance equation.

Conversely these data glue as a sheaf of sets in the fpqc slice,
by effectivity of the atlas.  To verify relative representability,
test on a scheme $T\to X_\infty$.  Its pulled-back atlas $T'\to T$
is fpqc, and the data give an etale finitely presented algebraic
space over $T'$ with descent datum.  Effective fpqc descent of
etale algebraic spaces yields an etale finitely presented space
over $T$.  Thus the glued slice object has exactly the source type
specified in Definition~\ref{def:frag-actual-repet-sources}.
This establishes the degree-$0,1,2$ description on which the
following finite-stage descent operates.  At a finite stage the
quotient construction is verified separately below.

Let $Y_\infty\to X_\infty$ be an object of the right-hand site.
Its pullback $V_\infty\to U_{\infty,0}$ is an etale finitely presented
algebraic space, with the descent isomorphism between its two
pullbacks to $U_{\infty,1}$.  First descend $V_\infty$ to an etale
finitely presented $V_n\to U_{n,0}$.  The two pullbacks of $V_n$
to $U_{n,1}$ are finitely presented; their isomorphism and its inverse
descend to a common later stage by the morphism part of the same
limit theorem.  The inverse identities and the unit identity descend
as equalities of morphisms.  The cocycle identity descends on
$U_{n,2}$ after one further refinement.  This is a finite amount of
ordinary groupoid descent data, not an interchange with the entire
infinite Cech totalization of derived coefficients.

At this stage the unit and degree-two cocycle are exactly a
$G_n$-action on $V_n$ over $U_{n,0}$, where $G_n$ is the finite
etale group of the root quotient chart.  Lemma~\ref{lem:frag-finite-quotient-etale-descent}
constructs $Y_n=[V_n/G_n]\to[U_{n,0}/G_n]=X_n$ and proves it
representable etale of finite presentation.  In particular the
quotient's representability is not supplied by the algebraic-space
limit theorem.  Pulling it
back to $X_\infty$ gives exactly the descent datum of $Y_\infty$.
Both pullbacks carry the same degree-zero object and degree-one
isomorphism with the degree-two cocycle.  Descent for objects and
arrows of the fpqc stacks along the effective atlas identifies $Y_n\times_{X_n}X_\infty$ with $Y_\infty$.  This proves
essential surjectivity without constructing a new etale realization
by definition.

A morphism between two such objects pulls back to a morphism of
finitely presented algebraic spaces over $U_{\infty,0}$, with its
compatibility equation over $U_{\infty,1}$.  Descend the morphism
and then this equation.  Equality of two morphisms is likewise
detected after increasing the stage.  This proves precisely the
filtered Hom formula and full faithfulness.

For a finite family of arrows, pull back to the affine atlas of
the target object.  Joint surjectivity is surjectivity of a single
finite coproduct of finitely presented etale morphisms.  If it holds
at the limit, the surjectivity part of algebraic-space limit descent
makes it hold at a later stage.  The converse follows by base change.
An arbitrary covering of a finitely presented object has a finite
subcover: its pullback has open images covering a quasi-compact
algebraic space.  The quotient lemma detects these covers on the atlas; their
pullbacks and finite composites are again detected there.  Thus
the finite-cover topology generates the full induced etale topology.
The only limit inputs were descent of finite-presentation algebraic
spaces on degrees $0,1,2$, followed by the finite quotient construction.

All comparisons are induced by pullback; the noncanonical finite
stage choices do not enter the resulting equivalence.  Finite
$\Delta$-descent can be included in the finite action-groupoid
presentations before the argument.  Their degrees $0,1,2$ still
have the required quasi-compactness and affine transitions, giving
the same proof and the asserted naturality.
\end{proof}

\begin{definition}[Pro-root site and double-index constructible category]
\label{def:frag-restricted-root-site}
For a root level $n\geq0$, let
\[
 \mathcal C_n:=
 \operatorname{RepEt}_{\mathrm{fp}}
 (\widetilde{\mathcal A}_{M,n}^{\mathrm{rel}})
\]
be its full representable, finitely presented \'{e}tale site.  If $m\geq n$,
root pullback gives $q_{m,n}^*:\mathcal C_n\to\mathcal C_m$.  Define
$\mathcal C_\infty$ as follows.  An object is a pair $(n,Y_n)$; a morphism is
\[
 \operatorname{Hom}_{\mathcal C_\infty}((n,Y_n),(n',Y'_{n'}))
 :=\varinjlim_{m\geq n,n'}
 \operatorname{Hom}_{\mathcal C_m}
 (q_{m,n}^*Y_n,q_{m,n'}^*Y'_{n'}).
\]
A finite family is covering if, after one common refinement $m$, it is jointly
surjective in $\mathcal C_m$.

The action of $\Delta$ on the root tower induces an action on its sheaf
category.  We use equivariant descent in that category:
\begin{equation}
\operatorname{Sh}(\mathcal C_\infty^\Delta)
:=\operatorname{Sh}(\mathcal C_\infty)^{h\Delta}
=\operatorname*{lim}_{B\Delta}\operatorname{Sh}(\mathcal C_\infty).
\label{eq:frag-delta-site-definition}
\end{equation}
Here the diagram consists of the inverse-image equivalences of the action;
its objects are sheaves with coherent equivariance data.  The notation on
the left abbreviates the action-groupoid sheaf category, not a claim that
taking equivariant representable objects is already a comparison of sites.
The analogous finite-level notation is used for $\mathcal C_n$.
The nerve $B\Delta$ has arbitrarily high simplicial degrees even though
$\Delta$ is finite.  No commutation with a filtered colimit is inferred from
calling that nerve a finite diagram.  Its coefficient descent is Corollary~\ref{cor:frag-relative-adic-action-descent}.

Let
$\mathscr X_{\mathrm{proot}}:=\operatorname{Sh}(\mathcal C_\infty^\Delta)$.
The intrinsic notion of constructibility used here is the coherent-topos one.
A finite stratification is a finite decomposition of the terminal topos by
coherent locally closed subtopoi; a complex is constructible if its cohomology
objects are locally constant with finite stalks on those subtopoi.  This
definition uses neither a Kummer presentation nor a finite root level.  Write
$\mathcal D^b_{c,\mathrm{pf}}$ for the full subcategory whose restrictions to
all strata are perfect over $\mathbf Z/\ell^r$.  The site and this perfect
constructible category are intrinsic.  The following Iwasawa-admissible
subcategory is then selected by the relative inertia supplied by $(P,\rho)$:
\begin{equation}
 \mathcal D^b_{c,\mathrm{Iw}}
 (\mathcal A_{M,\ell}^{\mathrm{rel}}/\mathscr S,\mathbf Z_\ell)
 :=\varprojlim_r
   \mathcal D^b_{c,\mathrm{pf}}(\mathscr X_{\mathrm{proot}},
                   \mathbf Z/\ell^r)_{\mathrm{unif,Iw}},
 \label{eq:frag-intrinsic-root-category}
\end{equation}
where the inverse limit is formed in stable $\infty$-categories.  The
subscript means: common finite cohomological and underlying-$R_r$ Tor
intervals, a finite stratification on which all cohomology sheaves have finite
stalks, the adic reduction equivalences, and
the underlying-$R_r$ perfectness at each precision in the displayed
category.  The Iwasawa interpretation is
Definition~\ref{def:frag-adic-iwasawa-sheaves}; at points its
additional classical finite-Tor properties follow from
Proposition~\ref{prop:frag-point-source-iwasawa-admissibility}.  The root
index does not enter this definition.  These conditions specify the selected perfect adic system, rather
than a locally free ordinary small-etale integral sheaf, and do not
impose torsion-free cohomology;
the complex $[\mathbf Z_\ell\xrightarrow{\ell}\mathbf Z_\ell]$ at $M=0$
belongs to the category.  Its finite-level derived presentation follows from
Theorem~\ref{thm:frag-relative-finite-action-descent}; its actual-root
interpretation is Corollary~\ref{cor:frag-actual-root-derived-realization}.
\end{definition}

\begin{lemma}[Continuous mapping complexes on a pro-root stratum]
\label{lem:frag-continuous-mapping-complexes}
Let $\mathscr T$ be a coherent stratum on which the relative inertia tower is
the quotient tower $\mathcal K_m$ of $\mathcal K_\ell$, and let $F_n,G_n$ be
perfect locally constant $\mathbf Z/\ell^r$-complexes at one finite level.  For
$m\ge n$ write $F_m,G_m$ for their inflations.  Then the natural map is an
equivalence of mapping spectra
\begin{equation}
 \operatorname*{hocolim}_{m\ge n}
 R\!\operatorname{Hom}_{B\mathcal K_m}(F_m,G_m)
 \xrightarrow{\sim}
 R\!\operatorname{Hom}_{B\mathcal K_\ell}^{\mathrm{cts}}(F,G).
 \label{eq:frag-continuous-mapping-complexes}
\end{equation}
After a split geometric pullback, both sides are represented by
\[
 C^*_{\mathrm{cts}}\!\left(
   \Gamma_M,R\!\operatorname{Hom}(F,G)\right)
 =\varinjlim_m
 C^*\!\left(\Gamma_M/\ell^m\Gamma_M,
             R\!\operatorname{Hom}(F,G)\right).
\]
No finite-level Ext-amplitude is asserted.  In particular, the equality allows
the periodic high-degree classes at finite levels to die under inflation.  The
statement is compatible with finite descent groups and finite recollement
diagrams.
\end{lemma}

\begin{proof}
At finite level the mapping object on the classifying stratum is the full
normalized inhomogeneous cochain complex on $\mathcal K_m$ with coefficients
in $R\!\operatorname{Hom}(F,G)$.  On the pro-stratum it is the analogous
continuous cochain complex.  A continuous cochain from
$\mathcal K_\ell^q$ to a finite discrete coefficient factors through some
$\mathcal K_m^q$; hence the displayed equality is degreewise before taking
cohomology.  Since $F$ and $G$ are perfect, their internal derived Hom has
finite amplitude, so each total degree contains only finitely many
bar/internal bidegrees.  Exact filtered colimits of finite modules therefore
commute with this totalization and with cohomology, proving
\eqref{eq:frag-continuous-mapping-complexes}.  This argument deliberately uses
the entire bar complex; boundedness of the objects does not truncate group
cohomology.  Finite recollement uses only finite limits and colimits.  For
$\Delta$, \eqref{eq:frag-delta-prime-to-ell} supplies the averaging idempotent
$|\Delta|^{-1}\sum_{\delta\in\Delta}\delta$; hence invariants are exact and
$H^{>0}(\Delta,-)=0$ on the torsion coefficient categories used here.  The
same calculation therefore commutes with the prescribed equivariant descent
without introducing another infinite group-cohomological tail.
\end{proof}

\begin{proposition}[Actual site descent and selected derived comparison]
\label{lem:frag-pro-root-site-descent}
The finite-covering rule of Definition~\ref{def:frag-restricted-root-site}
defines the filtered finite-level site.  By
Theorem~\ref{thm:frag-actual-root-fp-descent}, with the source convention
of Definition~\ref{def:frag-actual-repet-sources}, its comparison with
the actual root site is
\begin{equation}
\operatorname{RepEt}_{\mathrm{fp}}
(\widetilde{\mathcal A}_{M,\ell}^{\mathrm{rel}})
\simeq\mathcal C_\infty.
\label{eq:frag-pro-site-descent}
\end{equation}
After effective descent for the finite action groupoid, the corresponding
statement is
\begin{equation}
\operatorname{Sh}_{\acute et}^{\mathrm{rep}}
(\mathcal A_{M,\ell}^{\mathrm{rel}})
\simeq\operatorname{Sh}(\mathcal C_\infty)^ {h\Delta}
=: \operatorname{Sh}(\mathcal C_\infty^\Delta).
\label{eq:frag-pro-site-delta-descent}
\end{equation}
On the fixed noetherian split chart, the continuous-action topos
identification of Theorem~\ref{thm:frag-actual-root-continuous-topos} and
Theorem~\ref{thm:frag-relative-finite-action-descent} then gives
\begin{equation}
\mathcal D^b_{c,\mathrm{pf}}(\mathscr X_{\mathrm{proot}},R_r)
\simeq\operatorname*{colim}_n
\mathcal D^b_{c,\mathrm{pf}}
(\widetilde{\mathcal A}_{M,n}^{\mathrm{rel}},R_r)^{h\Delta}.
\label{eq:frag-pro-root-derived-descent}
\end{equation}
These comparisons concern the selected action/adic coefficient model;
the separate global pseudocompact interpretation retains the separate global pseudocompact comparison.
None is an unrestricted interchange of a
derived category with an inverse limit.
\end{proposition}

\begin{proof}
Identity coverings, base change of a finite covering, and composition of
finite coverings can each be represented at one common finite root level;
this proves the pretopology assertion.  The first site equivalence is
Theorem~\ref{thm:frag-actual-root-fp-descent}; finite etale descent for the quotient by $\Delta$ gives the second.
Theorem~\ref{thm:frag-actual-root-continuous-topos} identifies the
actual topos and structural pullback.  Relative derived-action descent
then gives the coefficient equivalence; finite semilinear descent is
performed in its action-sheaf category, without commuting a filtered
colimit through an arbitrary homotopy fixed-point limit.
On affine presentations the
candidate transition maps are
\begin{equation}
U_{m,q}\longrightarrow U_{n,q},\qquad m\ge n.
\label{eq:frag-compatible-affine-root-groupoids}
\end{equation}
Affine finite-presentation continuity supplies part of this comparison; see
\cite[Tags 07SK, 07SL, and 07SN]{StacksProject}.  It does not say that the
non-finitely-presented atlas $U_{\infty,0}$ is a hypercover in the full
representable etale site.  The fpqc effectivity statement and the required
etale cohomological descent statement must therefore remain separate.
For the independent adic Iwasawa coefficient model, the inverse
comparison and attaching mapping spectra are now proved by
Theorem~\ref{thm:frag-F-adic-iwasawa} and relative action descent.
The needed topos identification is the separate generating-site argument
of Theorem~\ref{thm:frag-actual-root-continuous-topos}.
\end{proof}

The actual pro-$\ell$ stack in the preceding display is the fpqc inverse limit
of the finite root fibres of order $\ell^n$; its associated inverse system is
a pro-DM, hence pro-geometric, object.  The actual limit is not asserted to be
a finite-type algebraic stack.  Proposition~
\ref{prop:frag-actual-root-effective-atlas} is the bridge between that actual
limit and the finite-level presentation; the transition maps are induced by
divisibility of the root indices.
Before finite descent by $\Delta$, the relative root fibre over the chosen
split chart is generally a nilpotent thickening of a classifying stack, not the
classifying stack itself.  Precisely,
\begin{equation}
 (\widetilde{\mathcal A}_{M,n}^{\mathrm{rel}})_{\mathrm{red}}
 \simeq B_S\mathcal K_n,
 \qquad
 \operatorname{Sh}_{\acute et}^{\mathrm{rep}}
   (\widetilde{\mathcal A}_{M,n}^{\mathrm{rel}})
 \simeq
 \operatorname{Sh}_{\acute et}^{\mathrm{rep}}(B_S\mathcal K_n)
 \quad(\Delta=1),
 \label{eq:frag-root-reduction-BK}
\end{equation}
with $\mathcal K_n$ as in Lemma~\ref{lem:frag-root-levels}.  The second
equivalence is
nil-invariance of the representable \'{e}tale topos.  After descent, the same
statement is retained as a $\Delta$-equivariant equivalence of topoi.
The quotient by $\Delta$ is shorthand for the stack obtained by effective
descent along the prescribed finite splitting groupoid.  In particular it
does not discard the $\Delta$-action or replace equivariant objects by
invariants.  The central fibre is canonical relative to the chosen chart,
section, and descent groupoid: it uses the actual relative log point and does
not require the quotient
$P/\mathbf N\rho$ to be fs.  It is the infinite-root enhancement of the
relative central face of the local Artin fan.  It should not be confused with
the ordinary quotient stack
$[\operatorname{Spec}_S\mathcal O_S[P]/D_S(P^{\mathrm{gp}})]$.
Before descent there is a canonical characteristic map over $S$
\[
 \widetilde{\mathcal A}_M^{\mathrm{rel}}
 \longrightarrow \sqrt[\infty]{\mathfrak s_{P/S}}
 \longrightarrow \mathcal A_{P/S},
 \qquad
 \mathcal A_{P/S}:=
 [\operatorname{Spec}_S\mathcal O_S[P]/D_S(P^{\mathrm{gp}})],
\]
which exhibits the split relative root face over the usual relative
local Artin cone.  These maps are $\Delta$-equivariant.  After descent
the corresponding arrows are over $\mathscr S=[S/\Delta]$:
\[
 \mathcal A_M^{\mathrm{rel}}
 \longrightarrow[\sqrt[\infty]{\mathfrak s_{P/S}}/\Delta]
 \longrightarrow[\mathcal A_{P/S}/\Delta].
\]
No map from the descended source to the unquotiented targets over $S$
is asserted.  The same convention applies to the pro-$\ell$ versions,
with every root index replaced by $\ell^\infty$.
The absolute expression with a residue
field $k$ is used only after taking a geometric fibre in
Proposition~\ref{prop:frag-artin-obstruction}.

\begin{lemma}[Functoriality of relative root faces]
\label{lem:frag-root-face-functoriality}
Let
$\alpha:(P_{\tau'},\rho_{\tau'})\to(P_\tau,\rho_\tau)$ be one of the
base-preserving saturated chart morphisms in
Definition~\ref{def:frag-root-presented-chart}, so that
$\alpha(\rho_{\tau'})=\rho_\tau$, the induced relative map is saturated,
and $\alpha$ respects the log-point augmentations (equivalently
$\alpha^{-1}(0)=\{0\}$ on these sharp charts).
For every root order $a$ invertible on the base, functoriality of the finite root-stack
construction and the trivial-root sections induces the corresponding map of
relative $a$th-root fibres.  These maps are compatible with divisibility in
$a$, with the finite descent groupoid $\Delta$, and with composition of chart
maps of the supplied log points.  Consequently, in a selected
codimension-two diagram whose two composite log-point morphisms agree,
the relative root composites agree by the canonical $2$-isomorphism,
before passing to the pro-$\ell$ limit.  Equality of the relative
lattice maps alone is not the hypothesis of this assertion.
This is functoriality of relative root charts, not an incidence theorem:
it supplies no cross-stratum specialization, link or Gysin operation.
\end{lemma}

\begin{proof}
Put $P'=P_{\tau'}$ and $P=P_\tau$.  At finite level the algebra map
$\mathcal O_S[\tfrac1aP']\to\mathcal O_S[\tfrac1aP]$
sends $X^{p'/a}$ to $X^{\alpha(p')/a}$.
The augmentation identity is exactly what allows this map to pass
to the base changes along the zero sections.  Thus it induces the
map between the presentations (for
$Q=P_{\tau'}$ and $Q=P_\tau$)
\[
 \left[\bigl(\operatorname{Spec}(S[\tfrac1aQ])
       \times_{\operatorname{Spec}(S[Q])}\operatorname{Spec}S\bigr)
       /D((\tfrac1aQ)^{\mathrm{gp}}/Q^{\mathrm{gp}})\right]
\]
for the two log points, in the direction from the $P$-root stack
to the $P'$-root stack.  The map is equivariant for the induced
homomorphism of the displayed diagonalizable groups.  Without the
augmentation identity, the projection $\mathbf N^2\to\mathbf N$,
$(a,b)\mapsto a$, would send a positive monomial to $1$ and would
not descend through the zero-section relation setting it to $0$.
This projection is therefore not an allowed log-point map here.
Exact preservation of $\rho$ makes the trivial-root
sections commute with this map, so the universal property of the $2$-fibre
gives the relative map.  The same presentation shows compatibility with
 refinement of $a$ and with $\Delta$.  For a composite of monoid maps the
induced maps of monoid algebras and diagonalizable groups compose strictly;
the universal property supplies the canonical associativity $2$-cell on the
fibres.  Augmentation compatibility is preserved by composition.
If one also wants to use a composite as a selected chart map for
coefficient operations, its other selection conditions must be
checked; no closure of finite-Tor coefficient operations under an
arbitrary rank change is inferred.  Passing to the inverse system
proves the assertion.
\end{proof}

\begin{lemma}[Finite root reductions and the relative Kummer group]
\label{lem:frag-root-levels}
For every root exponent $n\geq0$, the reduction of the split relative
geometric fibre is $B_S\mathcal K_n$ before $\Delta$-descent, where the finite
diagonalizable group scheme is
\[
 \mathcal K_n
 :=\underline{\operatorname{Hom}}
       (P^{\mathrm{gp}}/\mathbf Z\rho,\mu_{\ell^n}).
\]
The primitivity of $\rho$ and the torsion-freeness of
$L_{\mathrm{rel}}=P^{\mathrm{gp}}/\mathbf Z\rho$ give an exact sequence
\[
 1\longrightarrow \mathcal K_n\longrightarrow
 \underline{\operatorname{Hom}}(P^{\mathrm{gp}},\mu_{\ell^n})
 \longrightarrow\mu_{\ell^n}\longrightarrow1.
\]
Consequently
\[
 \mathcal K_\ell:=\varprojlim_n\mathcal K_n
 \simeq \operatorname{Hom}(L_{\mathrm{rel}},\mathbf Z_\ell(1))
 \quad\text{as a pro-\'{e}tale group sheaf over }S.
\]
For a strict geometric point $\bar s\to S$ and the chosen split-lattice
identification,
\[
 \mathcal K_{\ell,\bar s}
 :=\varprojlim_n\mathcal K_n(\bar s)
 \simeq M\otimes\mathbf Z_\ell(1)=\Gamma_M.
\]
After descent, the inertia-selected coefficient category carries the same continuous
$\Gamma_M$-action together with the prescribed $\Delta$-descent groupoid; no
assertion is made that the absolute stabilizer of the central fibre is just
$\Gamma_M$.
\end{lemma}

\begin{proof}
At level $\ell^n$, the standard presentation of the root stack of the log
point with chart $Q$ has an affine atlas whose ideal generated by the positive
fractional monomials is nilpotent.  Because $Q$ is sharp, its reduction is the
zero section.  Hence the reduced root stack is the classifying stack of
\(D(Q^{\mathrm{gp}}/\ell^nQ^{\mathrm{gp}})\), although the unreduced root
stack is not.  Taking the $2$-fibre over the base-root section restricts
characters along
\(\mathbf Z\rho\hookrightarrow P^{\mathrm{gp}}\).  The stabilizer of the
reduction is therefore
\(D(P^{\mathrm{gp}}/(\mathbf Z\rho+\ell^nP^{\mathrm{gp}}))\), equivalently
the displayed $\mathcal K_n$.  Nil-invariance gives
\eqref{eq:frag-root-reduction-BK}.  Since $\rho$ is primitive and the quotient
is torsion-free, restriction onto the base factor is surjective.  Passing to
the inverse limit and using
\(M=\operatorname{Hom}(L_{\mathrm{rel}},\mathbf Z)\) gives the final
identification.  The finite group $\Delta$ acts on every term and hence on
the resulting equivariant descent groupoid.
\end{proof}

\begin{remark}[Torsion and wild boundaries]
The torsion-free hypothesis on $L_{\mathrm{rel}}$ is substantive.  If the
relative quotient has torsion, the finite kernels $\mathcal K_n$ acquire an additional
finite diagonalizable factor, and the auxiliary site must be enlarged by the
corresponding finite tame characters.  If $\ell=p$, the groups
$\mu_{\ell^n}$ are not \'{e}tale and the present Kummer--\'{e}tale argument
does not apply.  Neither case is included in the selected coefficient theorem.
\end{remark}

\begin{definition}[Auxiliary relative Kummer coefficient category]
\label{def:frag-relative-kummer-category}
At root exponent $n$, let
$\mathcal F_{M,n}^{\mathrm{rel,ket}}$ be the Kummer--\'{e}tale chart site over
$S$ generated by the saturated Kummer--\'{e}tale extension
$P\to\frac1{\ell^n}P$ together with a trivialization of its restriction to
$\mathbf N\rho$, its strict-\'{e}tale base changes, and their descent data.
This definition uses the logarithmic chart and not the root central fibre.
The Cech action groupoid of its standard atlas has deck group scheme
$\mathcal K_n$.  Put
\begin{equation}
 \mathcal D^b_{c,\mathrm{ad}}
 (\mathcal F_M^{\mathrm{rel,ket}}/\mathscr S,\mathbf Z_\ell)^{\mathrm{Iw}}
 :=\varprojlim_r\left(
 \left(\varinjlim_n
 \mathcal D^b_c(\mathcal F_{M,n}^{\mathrm{rel,ket}},
                  \mathbf Z/\ell^r)\right)^\Delta
 \right)_{\mathrm{unif,Iw}}.
 \label{eq:frag-relative-kummer-double-category}
\end{equation}
Here ``$\mathrm{unif,Iw}$'' selects derived-cartesian coefficient
systems with a common finite cohomological interval, a common
underlying-$R_r$ Tor interval, and a finite constructible
stratification on which the terms at every precision are
underlying-$R_r$-perfect with finite cohomology stalks.
These are exactly the adic-system conditions in
\eqref{eq:frag-intrinsic-root-category}.  Reductions and their
coherence maps are taken in the ambient unbounded module-sheaf
categories; the displayed notation does not assert that reduction
preserves every bounded complex before this selection.
The superscript $\Delta$ means coherent semilinear descent along
the given action groupoid, not ordinary invariants.

For clarity the Iwasawa notation in this comparison means
\[
 \operatorname{Perf}^{\mathrm{cts}}_{\mathbf Z_\ell}
 (\mathscr S,\Lambda_{M,\mathscr S}^{\mathrm{Iw}})
 :=\bigl(\mathcal M_{\mathrm{Iw,ad}}^{\mathrm{sel}}(S)\bigr)^{h\Delta},
\]
with the semilinear action and the module-sheaf category of
Definition~\ref{def:frag-adic-iwasawa-sheaves}.
In that category the precision-$r$ coefficients are modules over
$\underline{\Lambda_{M,r}^{\mathrm{Iw}}}$ with bounded inertia-ideal
torsion cohomology; the Tor selection is over $R_r$, not over an
integral global Iwasawa sheaf.  The symbol
$\operatorname{Perf}^{\mathrm{cts}}_{\mathbf Z_\ell}$ here denotes
this selected adic system category.  It does not impose global
$\mathbf Z_\ell$-perfectness or completed-inertia Tor-amplitude on
an independently recovered sheaf.  At a geometric point the stronger
classical properties follow from
Proposition~\ref{prop:frag-point-source-iwasawa-admissibility};
that result is not used to assert global sheaf recovery.
Finite $\ell$-torsion cohomology remains allowed.
The adjective ``auxiliary'' is deliberate: this site is constructed without
using the root central fibre, but it is generated by the selected Kummer
charts, so its identification with the action groupoid is a descent
calculation rather than a novelty claim.  The actual-root comparison proves the passage from the specified
relative pro-root site to this category, including its structural
functor; it is not an assertion of a new absolute Kummer-topos equivalence.
\end{definition}

\begin{lemma}[Finite-level root--Kummer action-groupoid comparison]
\label{lem:frag-finite-root-kummer-comparison}
For every root exponent $n$, there are equivalences of topoi over $S$
\begin{equation}
 \operatorname{Sh}_{\acute et}^{\mathrm{rep}}
 (\widetilde{\mathcal A}_{M,n}^{\mathrm{rel}})
 \simeq
 \operatorname{Sh}_{\acute et}^{\mathrm{rep}}(B_S\mathcal K_n)
 \simeq
 \operatorname{Sh}(\mathcal F_{M,n}^{\mathrm{rel,ket}}).
 \label{eq:frag-finite-root-kummer-topoi}
\end{equation}
They are compatible with root refinement, strict base change, the
$\Delta$-action, and the augmentation to $S$.
\end{lemma}

\begin{proof}
The first equivalence is nil-invariance applied to
Lemma~\ref{lem:frag-root-levels}.  For the second, the standard relative
Kummer atlas has Cech nerve
$\mathcal K_n^{\times_S\bullet}$: trivializing the base root removes the
quotient character along $\mathbf Z\rho$, leaving exactly the kernel
$\mathcal K_n$.  Descent for this finite \'{e}tale action groupoid identifies
sheaves on the Kummer chart site with sheaves on $B_S\mathcal K_n$.
The monoid maps for $n\leq m$ give the quotient homomorphisms
$\mathcal K_m\to\mathcal K_n$ on both nerves, proving refinement
compatibility.  Every construction is functorial in strict base change and
$\Delta$-equivariant.
\end{proof}

\begin{lemma}[constructible etale spaces and finite action objects]
\label{lem:frag-constructible-etale-action-site}
Let $S$ be noetherian.  Etale algebraic spaces of finite presentation
over $S$ correspond to constructible sheaves of sets on $S_{\acute et}$.
For a finite constant group $K_n$ (the fixed split inertia group
in the root applications), representable etale finitely presented
morphisms to $B_SK_n$ correspond to such sheaves with a $K_n$-action.
Joint surjectivity corresponds to a jointly epimorphic family of sheaves.
These correspondences commute with inflation and with base pullback.
\end{lemma}

\begin{proof}
The etale-space construction represents a sheaf of sets by an etale
algebraic space, under the usual universe convention
\cite{StacksEtaleSpaceRealization}.  Here is the needed finiteness check.
A constructible sheaf $F$ is noetherian and is generated by finitely
many sections on quasi-compact etale schemes.  Thus there is an
epimorphism $h_U\to F$ with $U$ a finite coproduct of such schemes.
The relation $h_U\times_F h_U$ is represented by the open equality
locus in $U\times_SU$: equality of two sections is etale-local, and
subterminal sheaves over an etale scheme are its open subsets.
This open is quasi-compact since $U\times_SU$ is noetherian.  The
quotient etale algebraic space is consequently quasi-compact and
quasi-separated over $S$, hence of finite presentation.

Conversely, choose a finite quasi-compact etale scheme atlas $U$ of
an etale finitely presented space $V/S$.  Its relation admits a
finite quasi-compact etale scheme cover $W$.  Both $U$ and $W$ are
etale of finite presentation over $S$.  Their representable sheaves
are constructible, and their coequalizer is $h_V$, which is
constructible by closure under finite colimits of sheaves of sets
\cite[Section 59.73]{StacksConstructibleNoetherian}.  This also proves
full faithfulness, since the etale-space construction and restriction
to the small site are inverse on these objects.

Pullback along the finite etale torsor $S\to B_SK_n$ identifies a
representable etale object with an etale space $V/S$ and a $K_n$-action.
Conversely the quotient $[V/K_n]\to B_SK_n$ is representable etale
and finitely presented.  The description respects morphisms and
joint surjectivity by base change along that torsor.  Both inflation
and base pullback are the literal pullbacks of these descent data.
\end{proof}

\begin{proposition}[a generating site for continuous action sheaves]
\label{prop:frag-continuous-action-generators}
Let $\mathscr E=\operatorname{Sh}(S_{\acute et})$ and let
$\mathscr E_\Gamma$ be the category of sheaves of sets with a smooth
$\Gamma$-action, defined by $\operatorname*{colim}_n F^{H_n}=F$ as
in Definition~\ref{def:frag-relative-action-sheaves}.
This is a Grothendieck topos.  Its full subcategory
$\mathscr E_{\Gamma,c}$ of underlying-constructible objects with finite
stalks on a finite stratification, equipped
with the finite jointly epimorphic topology, is a generating site, and
\[
\operatorname*{colim}_n\mathscr E_c^{K_n}
 \simeq\mathscr E_{\Gamma,c},
\qquad
\operatorname{Sh}(\mathscr E_{\Gamma,c})\simeq\mathscr E_\Gamma.
\]
The target topos is independently defined by smooth sheaf actions.
\end{proposition}

\begin{proof}
The constructible-subobject and ascending-chain results used in
Lemma~\ref{lem:frag-noetherian-action-sheaves} also hold for sheaves
of sets.  Every smooth action sheaf is the union of its $H_n$-fixed
subsheaves.  In each such subsheaf, saturate a constructible subobject
by its finitely many $K_n$-translates.  Finite unions remain
constructible, so these finite-action objects cover every smooth
object.  If the underlying sheaf is constructible, the fixed-subobject
chain stabilizes globally, proving factorization through one $K_n$.
Equivariant morphisms are detected at a common level.  This proves
the first displayed equivalence without deriving anything.

For an arbitrary abstract $\Gamma$-action sheaf $A$, set
$\operatorname{Sm}(A)=\operatorname*{colim}_n A^{H_n}$, a union of
subsheaves.  Normality of $H_n$ makes this a smooth invariant subsheaf.
Every equivariant map from a smooth object to $A$ lands locally in
these fixed subsheaves and factors uniquely through $\operatorname{Sm}(A)$.
Thus $\operatorname{Sm}$ is right adjoint to inclusion of smooth actions,
and arbitrary limits are
\[
 \lim_i^{\mathrm{sm}}F_i=\operatorname{Sm}(\lim_i^{\mathrm{abs}}F_i).
\]
For $F_i=\Gamma/H_i$, the tuple of identity cosets is excluded when
its stabilizer $\bigcap_iH_i$ is not open.  Infinite products are
therefore not asserted to be the underlying products.

Finite limits and colimits of smooth actions are computed on underlying
sheaves; finite limits are smooth because finitely many local
stabilizers have open intersection.  Arbitrary colimits are likewise
smooth, since their sections are locally generated by sections of the
inputs.  Equivalence relations are effective and colimits are
universal, as for underlying sheaves.  The finite-action constructible
objects just constructed form a small generating family (after fixing
a universe).  These facts give the topos assertion by the usual
Grothendieck topos criterion.

The full generating subcategory is closed under finite limits.  A
jointly epimorphic family onto a constructible object has a finite
subfamily: finite unions of its images form an ascending directed
family of subobjects of a noetherian underlying sheaf.  The topology
induced by epimorphic families is therefore the stated finite topology.
The generating-site comparison lemma now identifies its sheaves with
$\mathscr E_\Gamma$.  Equivalently, an object is recovered from all
maps from these generators and their jointly epimorphic relations;
this is the canonical presentation of an object of a topos by a full
generating site.  It uses objects and their relations on the entire
base, not equality of the fibres of two proposed topoi.
\end{proof}

\begin{theorem}[The relative root topos]
\label{thm:frag-actual-root-continuous-topos}
On the fixed noetherian split tame chart, with the fpqc-stack source
convention of Definition~\ref{def:frag-actual-repet-sources}, there is
an equivalence of topoi over $S$
\[
\operatorname{Sh}\bigl(\operatorname{RepEt}_{\mathrm{fp}}(X_\infty)\bigr)
 \simeq \mathscr E_\Gamma,
\qquad X_\infty=\widetilde{\mathcal A}_{M,\ell}^{\mathrm{rel}}.
\]
The structural inverse image from $S_{\acute et}$ is identified with
trivial relative action.  The equivalence is compatible with root
refinement, the specified strict base changes, and finite semilinear
$\Delta$-descent.
\end{theorem}

\begin{proof}
At each finite level, etale finite-presentation objects lift uniquely
across the nilpotent monomial reduction.  This can be checked on the
finite affine root atlas: nil-invariance of etale spaces lifts the
object and every descent isomorphism uniquely, so also lifts its
unit and cocycle equations.  Consequently the finite site is the
site of representable etale fp objects over $B_SK_n$.  By
Lemma~\ref{lem:frag-constructible-etale-action-site}, it is the site
$\mathscr E_c^{K_n}$, and the identifications commute with root
pullback and covering families.

Theorem~\ref{thm:frag-actual-root-fp-descent} identifies the actual
site with the filtered colimit of these finite sites.  Proposition~
\ref{prop:frag-continuous-action-generators} identifies that colimit
with $\mathscr E_{\Gamma,c}$ and then its sheaves with
$\mathscr E_\Gamma$.  This constructs the topos equivalence, including
its topology.  The infinite fpqc atlas itself has not been declared
an etale hypercover.

For an etale finitely presented $U/S$, its structural pullback to a
finite root stack corresponds under nil reduction and the torsor
atlas to $h_U$ with trivial $K_n$-action.  The same holds in the
filtered site and hence in the actual topos.  Such $h_U$ generate the
base topos; both inverse-image functors preserve colimits and finite
limits.  Thus their canonical agreement on these generators extends
to the structural inverse image on all sheaves, proving the statement
over $S$ rather than only an abstract equivalence of topoi.

Every step is induced by pullback or by its canonical descent
comparison.  It therefore commutes with base change and with maps
of the given root groupoids.  Finite $\Delta$-descent is applied to
this equivalence with its supplied finite descent diagram, giving the asserted
comparison over $\mathscr S=[S/\Delta]$.
\end{proof}

\begin{corollary}[actual derived root star and adic coefficients]
\label{cor:frag-actual-root-derived-realization}
On the chart of Theorem~\ref{thm:frag-actual-root-continuous-topos}, let
$g:X_\infty\to\mathscr S$ be the structural morphism and let $E_r$
be a selected bounded coefficient.  Then the actual right direct
image is naturally
\[
Rg_*^{\mathrm{act}}E_r
 \simeq R\underline{\operatorname{Hom}}_{\Lambda_r}(R_r,E_r)
 \simeq C^*_{\mathrm{cts}}(\Gamma,E_r),
\]
with finite semilinear descent understood.  The actual root coefficient
category is equivalent to the adic Iwasawa module-sheaf category of
Definition~\ref{def:frag-adic-iwasawa-sheaves}, with the same selected
uniform restrictions.  The star outputs are uniformly bounded and
derived cartesian when the input system is, and their adic value is
formed by taking the coefficient inverse limit after this recovery.
The assertion concerns this actual site and coefficient model;
it does not construct an exceptional direct image on $X_\infty$.
\end{corollary}

\begin{proof}
The preceding theorem gives an exact equivalence of coefficient
abelian categories, identifying structural pullback with trivial
action.  Its right adjoint is the invariant-sheaf functor, since a
map from a trivial-action object has image in the invariant subsheaf.
The exact equivalence identifies injective resolutions, so the
actual derived right adjoint is derived continuous invariants.
Theorem~\ref{thm:frag-F-adic-iwasawa} identifies this with the
displayed internal Hom and its finite Koszul calculation.  The
comparison morphism has thus been constructed before it is checked
by that calculation.

Constructibility is respected by
Lemma~\ref{lem:frag-coherent-constructible-root-descent}.  More explicitly,
subterminal objects in the
smooth-action topos are exactly invariant open subobjects of the
terminal sheaf, hence the corresponding base opens.  On a stratum,
a locally constant finite object pulls back to a locally constant
finite base sheaf.  Conversely a finite locally constant base sheaf
with finite $K_n$-action becomes constant after a finite etale
trivialization on the base and the regular $K_n$-cover in the action
topos.  This identifies the constructible strata and the selected
coefficient subcategories; finite Tor-amplitude is tested on their
local complexes.  Relative derived-action descent handles their
attaching maps, so this is not just an assertion about cohomology
stalks.  Applying the exact equivalence to complexes and then to
cartesian coefficient systems yields the adic coefficient comparison.

The augmentation Koszul complex has length $d$ and finite free terms.
On the selected coefficients it commutes with derived reduction and
increases amplitude by at most $d$, uniformly in $r$.  This proves
recovery for the actual star output already identified above.  The
statement about its adic limit follows in the category of coefficient
systems.  On these bounded-cohomology coefficients, derived sheaf
realization is compatible with hyperdescent through the equivalence
of topoi; no unbounded hypercompletion theorem for arbitrary root
coefficients is asserted.
\end{proof}

\begin{lemma}[Constructible strata and their finite root presentations]
\label{lem:frag-coherent-constructible-root-descent}
In the actual-root topos just identified, intrinsic finite coherent
stratifications and finite locally constant coefficients agree with
the underlying-constructible continuous-action model.  Selected
bounded derived coefficients and their attaching maps descend to
finite root levels at each coefficient precision.
\end{lemma}
\begin{proof}
A subterminal object in $\mathscr E_\Gamma$ is a subobject of the
terminal base sheaf, with its unique action.  Hence opens, and then
locally closed subtopoi obtained by complements and restrictions,
are exactly the base ones.  On the noetherian base these opens are
quasi-compact.  A finite coherent stratification thus occurs already
at root level zero, before any coefficient is descended.
On each stratum, forgetting the action carries a finite locally
constant object to a finite locally constant base sheaf.  Conversely,
a finite locally constant base sheaf with smooth action factors
through a finite quotient by the noetherian fixed-subobject argument.
An etale trivialization of that base sheaf, together with the regular
finite-group torsor, trivializes it in the action topos.
Theorem~\ref{thm:frag-relative-finite-action-descent}, applied to the
whole constructible sheaf category, then descends the bounded objects
and attaching morphisms, not just the separate stalk values.
Finite semilinear descent is handled by
Lemma~\ref{lem:frag-finite-equivariant-rectification}.  Underlying
perfection is unchanged by inflation.  This proves the selected
constructible statement; it makes no compactness claim for infinite
stalks or infinitely many strata.
\end{proof}



\begin{proposition}[the relative log-fibre coefficient interface]
\label{prop:frag-relative-log-root-interface}
Let $Y$ be a noetherian split log stratum with characteristic chart
$\mathbf N\rho\to P$, with $\rho$ primitive and
$L=P^{\mathrm{gp}}/\mathbf Z\rho$ free.  Fix a geometric logarithmic
point over the base special point, including compatible prime-to-$p$
roots of the base parameter.  Write
$\varepsilon_{Y,\mathrm{rel}}:\widetilde Y_{\mathrm{ket}}\to Y_{\acute et}$
for the projection from the resulting relative logarithmic fibre.
On coefficients trivial on the non-$\ell$ relative tame factor, its
selected bounded/adic coefficient category agrees with that of the
actual pro-$\ell$ root fibre, and
\[
 R\varepsilon_{Y,\mathrm{rel},*}E
 \simeq C^*_{\mathrm{cts}}(\Gamma,E),\qquad
 \Gamma=\operatorname{Hom}(L,\mathbf Z_\ell(1)).
\]
Here the adic assertion uses the uniform restrictions of Theorem~\ref{thm:frag-F-adic-iwasawa}.
This is not an equivalence between the full prime-to-$p$ logarithmic
topos and the full pro-$\ell$ root topos.
\end{proposition}

\begin{proof}
On an etale split test object of $Y$, Kummer covers are generated by
adjoining prime-to-$p$ roots of the chart and then taking strict etale
covers.  The finite chart construction and its character pairing give
absolute inertia $\operatorname{Hom}(P^{\mathrm{gp}},\widehat{\mathbf Z}^{(p')}(1))$.
Fixing the roots of the base parameter takes the kernel of its map to
$\operatorname{Hom}(\mathbf Z\rho,\widehat{\mathbf Z}^{(p')}(1))$.
Primitivity makes this map surjective and the kernel is
\[
 G_{\mathrm{rel}}=\operatorname{Hom}(L,\widehat{\mathbf Z}^{(p')}(1))
                  =\Gamma\times H,
\]
where all finite quotients of $H$ have order prime to $\ell$.
There is consequently no extra finite component factor in this split
relative fibre.  A different chosen logarithmic lift is compared by
its torsor of lifts, not by declaring a canonical splitting of absolute
inertia.

The Kummer chart criterion says that every finitely presented Kummer
etale object is etale locally obtained from such a finite chart cover;
its strict part is ordinary etale descent.  Hence the same generating-site
argument used in Theorem~\ref{thm:frag-actual-root-continuous-topos}
identifies this relative log topos over $Y$ with smooth
$G_{\mathrm{rel}}$-action sheaves.  On the pro-$\ell$ selected subcategory,
inflation identifies the coefficient abelian category with smooth
$\Gamma$-actions.  Averaging over finite quotients proves that smooth
$H$-invariants are exact, as in Lemma~\ref{lem:frag-wild-inertia-criterion}.
Thus deriving invariants gives
$R\Gamma(G_{\mathrm{rel}},E)=R\Gamma(\Gamma,E)$ on this subcategory.
The exact equivalence, relative derived-action descent, and Theorem~\ref{thm:frag-F-adic-iwasawa}
identify the selected derived coefficients and their adic systems with
those of the actual root fibre.  The structural inverse image on both
sides is trivial relative action, so this comparison identifies the
stated structural right adjoints as well as their sources.

The chart construction commutes with strict base pullback and with
its finite Kummer refinements.  Its character pairing is equivariant
for the specified finite semilinear descent data.  This supplies these
compatibilities on a split stratum; it asserts no comparison for
incidence maps between different characteristic strata.
\end{proof}


\begin{corollary}[Intrinsic adic coefficients]
\label{cor:frag-A0-adic}
On the fixed noetherian split chart, with representable etale fpqc-stack
sources as in Definition~\ref{def:frag-actual-repet-sources}, the selected
adic coefficient category on the actual infinite-root central fibre is
equivalent to the independently defined category
$\mathcal M_{\mathrm{Iw,ad}}^{\mathrm{sel}}$ of
Definition~\ref{def:frag-adic-iwasawa-sheaves}, with finite semilinear
descent understood on both sides.  This equivalence identifies mapping
spectra, attaching maps and coefficient reduction.  Structural star is
continuous invariants, with derived adic recovery on the selected uniform
class.  This target does not require a global pseudocompact interpretation.
\end{corollary}

\begin{proof}
Apply Theorem~\ref{thm:frag-actual-root-continuous-topos} to identify the
actual coefficient topos, then Theorem~\ref{thm:frag-F-adic-iwasawa} to
identify its independent adic module-sheaf target.  The structural and
recovery assertions are Corollary~\ref{cor:frag-actual-root-derived-realization}.
This identifies the explicitly defined adic module-sheaf target.  No geometric operation
transport is inferred from this coefficient equivalence.
\end{proof}


\begin{proposition}[Actual root star image and its coefficient scope]
\label{prop:frag-actual-root-star-direct-image}
Work on the fixed noetherian split chart with the selected adic
module-sheaf coefficient category of
Definition~\ref{def:frag-adic-iwasawa-sheaves}.
Then the actual structural direct image is computed by
\begin{equation}
Rg_{M,e,\ell,*}^{\mathrm{act}}E_r
\simeq C^*_{\mathrm{cts}}(\mathcal K_\ell,E_r).
\label{eq:frag-actual-root-star-equals-continuous-bar}
\end{equation}
The right side is an internal sheaf-valued continuous cochain object.
For an admissible adic coefficient its derived-cartesian recovery gives
$Rg_*^{\mathrm{act}}E\simeq\mathsf G_{*,\mathrm{ad}}^{\mathrm{cts}}(E)$.
\end{proposition}

\begin{proof}
Theorem~\ref{thm:frag-actual-root-continuous-topos} identifies structural pullback with the
functor assigning trivial relative inertia action.  Its right adjoint is
continuous invariants; deriving that adjunction in the specified coefficient
realization gives the displayed formula.  The finite-projective Iwasawa
calculation then proves boundedness and adic recovery for the selected
coefficients.  Homotopy invariance removes the affine carrier for the star
image.  This argument requires a comparison over the base and its derived
coefficient realization.  An abstract equivalence of selected categories, or
an effective fpqc atlas by itself, does not identify this right adjoint.
\end{proof}

\begin{remark}[Full site, selected coefficients]
The coefficient comparison uses the full representable \'{e}tale site of every finite central
fibre.  What is selected on that intrinsic site is the coefficient subcategory with
continuous relative pro-$\ell$ inertia, not a chartwise-defined target site.
The ordinary lisse--\'{e}tale cohomology of the unrigidified Artin fan still
contains the residual classifying-stack direction exposed by
Proposition~\ref{prop:frag-artin-obstruction}.
\end{remark}


\section{The continuous dualizing module}
\label{sec:frag-iwasawa}

Continuous cohomology has a finite Iwasawa resolution.  Finite root
stacks do not: their full bars remain unbounded.  We first describe
the continuous algebra and its dualizing bimodule, and then show that
geometric finite-stack traces converge to this dual partner.

\subsection{Continuous bar--Iwasawa models and the dualizing module}



\begin{lemma}[Equivariant bar--Iwasawa algebra and the top orientation]
\label{lem:frag-equivariant-bar-iwasawa}
Put
\[
 W_M:=\operatorname{Hom}_{\mathrm{cts}}(\Gamma_M,\mathbf Z_\ell)
      \simeq M^\vee\otimes\mathbf Z_\ell(-1).
\]
The basis-free object
\[
 C^*_{\mathrm{cts}}(\Gamma_M,\mathbf Z_\ell)
 :=R\!\operatorname{Hom}^{\mathrm{cts}}_{\Lambda_M^{\mathrm{Iw}}}
      (\mathbf Z_\ell,\mathbf Z_\ell)
\]
carries its canonical homotopy-coherent cup product and, for every finite
group $\Delta$ acting on $M$, its canonical $\Delta$-equivariant structure.
The completed bar model, with its shuffle product, and every basis Koszul
model represent this same multiplicative derived object.  In particular,
\begin{equation}
 H^q C^*_{\mathrm{cts}}(\Gamma_M,\mathbf Z_\ell)
   \simeq \bigwedge^q W_M,
 \qquad
 \det W_M\simeq\mathfrak o_M^\vee(-d),
 \qquad
 \mathfrak o_M:=\det(M\otimes\mathbf Z_\ell),
 \label{eq:frag-equivariant-top-orientation}
\end{equation}
and all three identifications are $\Delta$-equivariant.  Thus the determinant
and Tate factors used below descend without choosing a basis.  What is
canonical is the derived cochain algebra; no strictly functorial Koszul basis
is asserted.
\end{lemma}

\begin{proof}
The augmentation ideal $I\subset\Lambda_M^{\mathrm{Iw}}$ satisfies
$I/I^2\simeq\Gamma_M\otimes_{\mathbf Z_\ell}\mathbf Z_\ell$, functorially in
$M$, and the associated graded algebra is the symmetric algebra on $I/I^2$.
The regular-sequence calculation therefore identifies the Yoneda algebra of
the augmentation module with $\bigwedge^\bullet W_M$.  Continuous group
cochains give the intrinsic multiplicative model.  The shuffle map on the
completed bar resolution represents its cup product, while comparison of
projective resolutions transports that product to any chosen Koszul model;
the result is canonical in the chosen enhancement through comparison
with the intrinsic continuous cochain algebra.  No contractibility
claim for an unspecified space of model choices is needed.  Every construction commutes with lattice automorphisms, hence with
$\Delta$.  Taking the top exterior power and using
$\det(M^\vee)=\det(M)^\vee$ gives the last two identities in
\eqref{eq:frag-equivariant-top-orientation}.
\end{proof}

For an Iwasawa-admissible coefficient define
\begin{align}
 \mathsf K^{\mathrm{ad}}_{M,*}(V_{\mathrm{cont}})
  &:=R\!\operatorname{Hom}^{\mathrm{cts}}_{\Lambda_M^{\mathrm{Iw}}}
       (\mathbf K_M^{\mathrm{Iw}},V_{\mathrm{cont}}),
  \label{eq:frag-koszul-star-ad}\\
 \mathsf K^{\mathrm{ad}}_{M,!}(V_{\mathrm{cont}})
  &:=\bigl(
       \mathsf K^{\mathrm{ad}}_{M,*}(V_{\mathrm{cont}}^\vee)
       \bigr)^{\vee_0}
       \otimes^{\mathbf L}_{\mathbf Z_\ell}\mathbf Z_\ell(-d-e)[-2(d+e)].
  \label{eq:frag-koszul-shriek-ad}
\end{align}
Their rationalizations are denoted by $\mathsf K_{M,\bullet}(V)$.
The factor $(-d-e)[-2(d+e)]$ is the inverse of the relative dualizing factor
$(d+e)[2(d+e)]$ inside the product-chart identity
\[
 \mathbb D_X(f^*\mathcal L\otimes\widetilde V)
 \simeq f^*\mathbb D_{\mathscr S}(\mathcal L)\otimes
 \widetilde{V^\vee}\otimes\omega_{X/\mathscr S}.
\]
It is not an additional orientation correction.
On the third realization it is the compact-support class of the actual
rank-$(d+e)$ affine carrier \eqref{eq:frag-smooth-enhanced-root-face}; it is no
longer attached to the bare central fibre by normalization.
If $V_{\mathrm{cont}}$ is trivial in the characteristic-torus direction,
Iwasawa--Koszul self-duality gives
\begin{equation}
 \mathsf K^{\mathrm{ad}}_{M,!}(V_{\mathrm{cont}})
 \simeq
 \mathsf K^{\mathrm{ad}}_{M,*}(V_{\mathrm{cont}})\otimes
 \det(M\otimes\mathbf Z_\ell)[-d]
 \otimes\mathbf Z_\ell(-e)[-2e],
 \label{eq:frag-orientation-factor-ad}
\end{equation}
and hence, after tensoring with $\mathbf Q_\ell$,
\begin{equation}
 \mathsf K_{M,!}(V)
 \simeq
 \mathsf K_{M,*}(V)\otimes
 \det(M\otimes\mathbf Q_\ell)[-d]
 \otimes\mathbf Q_\ell(-e)[-2e].
 \label{eq:frag-orientation-factor}
\end{equation}
Thus the determinant line is an object-level adic factor, not a $K_0$ artifact.

\begin{proposition}[The base-linear continuous dualizing bimodule]
\label{prop:frag-continuous-dualizing-bimodule}
Put
\begin{equation}
 A^{*,\mathrm{ad}}_M:=
 C^*_{\mathrm{cts}}(\Gamma_M,\mathbf Z_\ell),
 \qquad
 A^{!,\mathrm{ad}}_{M,e}:=
 (A^{*,\mathrm{ad}}_M)^{\vee_0}
 \otimes\mathbf Z_\ell(-d-e)[-2(d+e)].
 \label{eq:frag-star-shriek-algebras}
\end{equation}
Their rationalizations are denoted
$A_M^*:=A_M^{*,\mathrm{ad}}[1/\ell]$ and
$A_{M,e}^!:=A_{M,e}^{!,\mathrm{ad}}[1/\ell]$; over $\mathscr S$ these are
understood with their specified lattice and finite descent data.

The cup product makes $A^{*,\mathrm{ad}}_M$ an $E_\infty$-algebra.
Dualizing its left and right regular actions makes
$A^{!,\mathrm{ad}}_{M,e}$ an
$(A^{*,\mathrm{ad}}_M)^e$-module, where
\[
 (A^{*,\mathrm{ad}}_M)^e:=
 A^{*,\mathrm{ad}}_M
 \widehat\otimes^{\mathbf L}_{\mathbf Z_\ell}
 (A^{*,\mathrm{ad}}_M)^{\mathrm{op}}.
\]
The evaluation map
\begin{equation}
 A^{!,\mathrm{ad}}_{M,e}
 \widehat\otimes_{\mathbf Z_\ell}
 A^{*,\mathrm{ad}}_M
 \longrightarrow \mathbf Z_\ell(-d-e)[-2(d+e)]
 \label{eq:frag-balanced-dualizing-evaluation}
\end{equation}
is balanced for the two cup actions.  Because the Iwasawa--Koszul model is a
bounded finite free $\mathbf Z_\ell$-complex, it is dualizable over the
coefficient base.  Put $L=\mathbf Z_\ell(-d-e)[-2(d+e)]$.
The displayed evaluation and the coevaluation
$L\to A^{*,\mathrm{ad}}_M\widehat\otimes A^{!,\mathrm{ad}}_{M,e}$
are the ordinary base-duality maps tensored with $L$.
After removing this invertible twist, the pair
$(A^{*,\mathrm{ad}}_M,(A^{*,\mathrm{ad}}_M)^{\vee_0})$ obeys the
usual triangle identities.  Balancing and dinaturality with respect
to regular multiplication are asserted; no $A^e$-linear coevaluation
from the base unit with an augmentation action is asserted.  No assertion
that $A^{*,\mathrm{ad}}_M$ is perfect as a module over its enveloping algebra
is used.  Thus
$A^{!,\mathrm{ad}}_{M,e}$, rather than a bare determinant line, is the local
candidate compact-support dualizing bimodule, identified geometrically
in Corollary~\ref{cor:frag-C2-local-closed} on its geometric domain.
For every coefficient $V_{\mathrm{cont}}$ in the selected source,
$\mathsf K^{\mathrm{ad}}_{M,*}(V_{\mathrm{cont}})$ is naturally an
$A^{*,\mathrm{ad}}_M$-module and
$\mathsf K^{\mathrm{ad}}_{M,!}(V_{\mathrm{cont}})$ has the dual action; all
geometric comparison maps preserve these actions under the
multiplicative star and supported comparisons below.  Formula
\eqref{eq:frag-orientation-factor-ad} is the constant-character calculation
of this bimodule, not a replacement for it.
In particular, this proposition does not assert an $A^e$-linear equivalence
$A^{*,\mathrm{ad}}_M\simeq A^{!,\mathrm{ad}}_{M,e}$ and does not equip
$A^{*,\mathrm{ad}}_M$ with a Frobenius or Calabi--Yau structure.  The proven
object is the base-linear dualizing $A^e$-module together with its balanced
evaluation and base-linear coevaluation.
\end{proposition}

\begin{proof}
Continuous invariants are the derived right adjoint to the strong
symmetric-monoidal trivial-action functor.  Their canonical lax
symmetric-monoidal structure sends the unit to the stated
$E_\infty$-algebra and defines its regular actions, before a lattice
basis is chosen.  The completed normalized bar with its
Alexander--Whitney diagonal models the binary cup product; that
diagonal alone is not being used to supply all $E_\infty$ cells.  The finite free Iwasawa--Koszul replacement proves
that the resulting augmented cochain object is perfect over
\(\mathbf Z_\ell\).  It is therefore dualizable in the enhanced coefficient
category.  Dualizing the two regular actions gives the displayed bimodule;
evaluation is balanced by associativity of cup product, and the ordinary
dualizability coevaluation has source
\(\mathbf Z_\ell(-d-e)[-2(d+e)]\) and target
\(A^{*,\mathrm{ad}}_M\widehat\otimes
A^{!,\mathrm{ad}}_{M,e}\).  After removing the factor $L$, these are the ordinary triangle
identities for $A$ and $A^{\vee_0}$.  Under
$A\otimes A^{\vee_0}\simeq R\operatorname{Hom}_{\mathbf Z_\ell}(A,A)$,
coevaluation corresponds to the identity.  Its dinaturality for
left and right multiplication, and evaluation's balancing
$\operatorname{ev}(\phi a,x)=\operatorname{ev}(\phi,ax)$
(with the graded conventions), are the required action identities.
They do not make the base unit into an unspecified $A^e$-module.  Every construction commutes with
strict-etale descent and lattice automorphisms; hence the bimodule and its
trace descend over \(\mathscr S\) with the \(\Delta\)-action.  Applying the
same bar action to a coefficient module proves the final assertion.
\end{proof}

For the dual Kummer local system
\(V_M^{\mathrm{Kum}}:=M^\vee\otimes\mathbf Q_\ell(-1)\) and
\(\mathfrak o_M:=\det(M\otimes\mathbf Q_\ell)\), the determinant bookkeeping is
\[
 \det V_M^{\mathrm{Kum}}=\mathfrak o_M^\vee(-d),
 \qquad
 (\det V_M^{\mathrm{Kum}})^\vee=\mathfrak o_M(d).
\]
This is the convention used when the orientation factor is rewritten in
terms of the torus orientation line.

\begin{remark}[Rank-one constant-coefficient check]
For $d=1$, $e=0$, and the constant coefficient $V=\mathbf Q_\ell$, the
two-term model gives
\[
 \mathsf K_{M,*}(\mathbf Q_\ell)
 \simeq \mathbf Q_\ell\oplus\mathbf Q_\ell(-1)[-1],
 \qquad
 \mathsf K_{M,!}(\mathbf Q_\ell)
 \simeq \mathbf Q_\ell[-1]\oplus\mathbf Q_\ell(-1)[-2].
\]
The second identity is the normalized continuous Verdier/Borel--Moore mate,
including the rank-one orientation shift.
By Theorem~\ref{thm:frag-theorem-B-local-closed} and
Corollary~\ref{cor:frag-C2-local-closed}, the same two identities hold for the
geometric complexes $\mathsf G_*$ and $\mathsf G_!$ on a chart satisfying the
stated hypotheses.
\end{remark}


\section{Geometric comparison tools}
\label{sec:frag-analytic-nearby}

The root calculation becomes a realization theorem only after an
actual generic coefficient has been constructed.  We specify the
specialization sites, compare ordinary and tame nearby cohomology,
and construct the local Kummer source.  The final mapping comparison
establishes full faithfulness on that source.

\begin{definition}[The actual analytic specialization sites]
\label{def:frag-analytic-specialization-sites}
For the geometric tube $T$ of a special formal scheme with reduction $Y$,
Berkovich's construction uses morphisms of topoi
\[
 T_{\acute et}\xleftarrow{\ \mu\ }T_{q\acute et}
                  \xrightarrow{\ \nu\ }Y_{\acute et},\qquad
 R\Theta(F)=R\nu_*\mu^*F.
\]
Here $q\acute et$ denotes the quasi-etale site.  Formal etale lifting
followed by generic fibre is quasi-etale, and does not in general define
a morphism from the ordinary analytic etale topos to $Y_{\acute et}$.
This is the construction of \cite[\S2, Proposition 2.1 and the
definition of $\Theta_K$ following it]{BerkovichFormalII}: formal
etale lifting defines $\nu$, and the inclusion of ordinary etale
objects in the quasi-etale site defines $\mu$.

In all analytic star formulas in this paper, $R\operatorname{sp}_*F$
is shorthand for $R\Theta(F)$ with this factorization.  The analytic
coefficient $F$ itself is on $T_{\acute et}$.  Pullback of a general
sheaf $L$ on $Y$ gives $\nu^*L$ on $T_{q\acute et}$; it is not thereby
an algebraically constructible generic coefficient, nor has a lift to
$T_{\acute et}$ been specified.  For finite-torsor descent of the formal
star functor one works on $T_{q\acute et}$ with the actual coefficient
$\mu^*F$ and applies $R\nu_*$.  This requires no interchange of $\mu^*$
with an arbitrary infinite totalization.
\end{definition}



Fix a geometric generic point $\bar\eta$ and a completed algebraic closure
$\widehat{\overline K}$.  Let
$\operatorname{sp}_{\bar\eta}:]S[_{\bar\eta}^{\mathrm{an}}/\Delta\to\mathscr S$
be geometric specialization.  The analytic functors use the full geometric
etale site, not arithmetic $K$-analytic cohomology.  Base tame inertia is an
external category action; an equivariant object requires a cartesian section.
For the integral Kummer local system $E_\bullet=(E_r)_r$, set
\[
\mathsf A_{*,r}(E_r):=R\operatorname{sp}_{\bar\eta,*}E_r.
\]
For the supported comparisons in this paper the target is a geometric
point.  Analytic compact support is the actual functor
\begin{equation}
\mathsf A_{!,r}(E_r):=R\Gamma_c(T,E_r).
\label{eq:frag-analytic-sp-shriek-definition}
\end{equation}
When computed through a compactification, extension by zero is followed
by the proper analytic structural direct image.  This formula defines
neither $\nu_!$ nor an exceptional functor for relative specialization.
Define $\mathsf A_{\star,\mathrm{ad}}=R\!\varprojlim_r\mathsf A_{\star,r}$
only with its stated coefficient system; its bounded constructible adic
recovery is a conclusion of the relevant comparison, not automatic.
Write $\mathsf A_\star=\mathsf A_{\star,\mathrm{ad}}[1/\ell]$ afterwards.

Let $\overline{\mathfrak U}$ be the fixed proper log compactification,
$\bar j_\eta:\mathfrak U_{\bar\eta}\hookrightarrow
\overline{\mathfrak U}_{\bar\eta}$ its generic open immersion,
$\bar i_Y:Y\hookrightarrow\overline{\mathfrak U}_s$ the selected locally
closed face, and $f_Y:Y\to\mathscr S$ its structural map.
The face need not be proper: being locally closed in a proper model does
not imply properness.  Require $f_Y$ to be in the compactifiable domain when
using $Rf_{Y,!}$, and otherwise leave the supported functor undefined.
For shorthand put $\Psi_r=R\Psi^t_{\overline{\mathfrak U}}$ at precision $r$.
The two geometrically typed candidates are
\begin{equation}
\begin{aligned}
\mathsf N_{*,r}(E_r)&:=
Rf_{Y,*}\bar i_Y^*\Psi_r(R\bar j_{\eta,*}E_r),\\
\mathsf N_{!,r}^{\mathrm{cost}}(E_r)&:=
Rf_{Y,!}\bar i_Y^!\Psi_r(R\bar j_{\eta,!}E_r).
\end{aligned}
\label{eq:frag-nearby-star-shriek-definition}
\end{equation}
We use $\mathsf N_{!,r}$ for the second, costalk construction.  The superscript
may be restored when comparing conventions.  Its identification with analytic
compact support is proved below for geometric closed faces.  These formulas do not assert that nearby cycles
commute with duality or that restriction to a selected face preserves a
compact-support comparison.  If $f_Y$ is proper, only the outer $!$ may be
replaced by $*$.


\begin{definition}[Kummer-basic source and relative tame inertia]
\label{def:frag-kummer-basic-source}
Let $\mathcal L$ be a perfect constructible adic object on $\mathscr S$, and
let $V_{\mathrm{cont}}$ be an Iwasawa-admissible coefficient.  The notation
$\mathcal L\boxtimes V$ denotes its pullback to a realization chart tensored
with the Kummer local system defined by $V$.  On the root/action side
this is defined intrinsically.  On the analytic/nearby side it is only
notation for an actually supplied common source realization; it is not
a construction for every constructible $\mathcal L$.  In particular,
$\nu^*\mathcal L$ is a quasi-etale coefficient and need not come from
an algebraic generic object, by
Proposition~\ref{prop:frag-constructible-source-obstruction}.
The local theorem supplies that realization in its stated domain.
For a general statement, a \emph{Kummer-basic common source category}
means a specified enhanced category with the actual realization
functors and coefficient maps.  Closure under finite cones and
retracts uses maps in that source category; tensoring by a perfect
base object requires its compatible actual realization as well.
It is not automatically the full stable root subcategory generated
by the displayed object values.  In Theorem~\ref{thm:frag-theorem-B-local-closed},
Corollary~\ref{cor:frag-B-local-essential-image} proves the required
full mapping and cone closure.  On the full proper-image source,
Proposition~\ref{prop:frag-B-proper-morphism-obstruction} prevents
that automatic identification.  This qualification is part of the
source typing, not an assumed theorem lifting arbitrary root maps.

On the tame nearby side, the chosen base-root section singles out the
relative characteristic kernel of the finite deck extension in
Lemma~\ref{lem:frag-root-levels}; on a split geometric fibre its limit
is $\Gamma_M$.  We write $R\Psi^t_{\mathrm{rel}}$ for restriction to
this relative inertia, not for invariants under base tame inertia.
Nearby cycles retain their own geometric tame action.  No global
weak base-inertia action on the root coefficient category, or
comparison with such an action, is asserted here.  Additional
equivariance requires a supplied action diagram on the chart, its
compactification and the coefficients.
\end{definition}


\begin{proposition}[actual finite Galois descent and its comparison map]
\label{prop:frag-actual-finite-galois-descent}
Let $\pi:U\to S$ be a morphism of the actual coefficient topoi under
consideration, and let $q:V\to U$ be an actual finite etale $K$-torsor,
with $K$ finite.  For a bounded-below coefficient $\mathcal E$, writing
$p=\pi q$, there is a canonical equivalence
\[
R\pi_*\mathcal E\xrightarrow{\sim}
 \operatorname{Tot}C^*(K,Rp_*q^*\mathcal E).
\]
It is natural for maps of these torsors and their quotient groups.
On the unit coefficient it is multiplicative in the enhanced sense;
coefficient comparisons are module comparisons over it.  If
$q^*\mathcal E=p^*E$ is a specified equivariant trivialization, the
adjunction unit further supplies a natural map
\[
 C^*(K,E)\longrightarrow
 \operatorname{Tot}C^*(K,Rp_*p^*E)
 \simeq R\pi_*\mathcal E.
\]
The second map is not asserted to be an equivalence at finite level.
\end{proposition}

\begin{proof}
The augmented Cech diagram of the finite etale cover resolves
$\mathcal E$.  This may be checked after the exact conservative
pullback $q^*$, where the cover has a section and the augmented
diagram is split.  Its totalization is therefore $\mathcal E$ in
the bounded-below derived category.  Applying the right adjoint
$R\pi_*$ commutes with this limit.  The torsor identifications
$V\times_U\cdots\times_UV=V\times K^a$ identify the resulting
cosimplicial terms and all face and degeneracy maps with the full
group-cochain diagram on $Rp_*q^*\mathcal E$.  This proves the first
formula without identifying $Rp_*q^*\mathcal E$ with its degree-zero
cohomology.

The map in the second formula is induced by the equivariant unit
$E\to Rp_*p^*E$.  All maps come from the actual cover and the
pullback--pushforward adjunction.  Consequently a refinement of
torsors gives a commuting diagram, including its composition cells.
Pullback is symmetric monoidal and its right adjoint has its
canonical lax symmetric-monoidal structure; the unit and Cech maps
respect these structures.  The same augmented descent argument in
commutative algebra objects and their modules gives the claimed
multiplicative/module comparisons.  Thus the enhancement belongs to
the constructed map, not to a chosen cohomology-ring isomorphism.
There is no assertion here that an unbounded finite group bar
commutes with derived coefficient reduction.
\end{proof}

\begin{lemma}[algebraic torus cohomology and root weights]
\label{lem:frag-algebraic-torus-weights}
Let $S$ be noetherian with $\ell$ invertible and use a split Kummer
tower.  Let $\pi:T\times_S\mathbf A_S^e\to S$, where
$T=\mathbb T_{M,S}$.  If $E$ is a constructible $R_r$-sheaf on $S$,
then
\[
R^j\pi_*\pi^*E
 \simeq E\otimes_{R_r}\bigwedge^j(M^\vee\otimes R_r)(-j),
 \qquad 0\le j\le d,
\]
with zero in the other degrees.  The identification is natural in
$E$ and base pullback.  On these sheaves, Kummer refinement by
$[\ell^{m-n}]$ acts as $\ell^{j(m-n)}$ in degree $j$; deck
translations act trivially on the exterior factor.  This is a
statement about relative cohomology sheaves, not a trivialization of
the entire derived deck action.
\end{lemma}

\begin{proof}
The affine-space factor has ordinary cohomology $E$ by homotopy
invariance.  The same fact and the one-dimensional torus calculation
can be obtained uniformly from $\mathbf P^1_S$.  If
$p:\mathbf P^1_S\to S$ and $i$ is one of its standard sections,
smooth purity and $pi=\mathrm{id}$ give
$i^!p^*E=E(-1)[-2]$.  The projective bundle formula gives
$Rp_*p^*E=E\oplus E(-1)[-2]$.  The localization triangle for the
complement of infinity cancels the second summand and gives the
assertion for $\mathbf A^1_S$.  Removing both zero and infinity
instead gives a map from two copies of $E(-1)[-2]$ to this second
summand; both section classes have degree one.  Its cofiber has
cohomology $E$ in degree zero and $E(-1)$ in degree one.
The relative degree-one generator is the Kummer class of the
coordinate $t$.  These calculations use the actual proper
compactification and its purity maps, so apply to coefficients
pulled from $S$ and commute with base change.

Iterating the one-dimensional calculation and the projection formula
gives the displayed exterior cohomology sheaves.  No cohomology
sheaf of the coefficient is assumed flat: the exterior factors are
finite free, and the projective-bundle and localization computations
are carried out on the coefficient itself.  The Kummer sequence sends
$t^a$ to $a$ times the class of $t$, proving the stated degree-one
weight.  Cup products of the coordinate classes give its $j$th
exterior power.  Translation $t\mapsto\zeta t$ changes the Kummer
class only by a class from the base, which vanishes in relative
$R^1\pi_*$.  It therefore acts trivially on the exterior cohomology
factor.  This last observation does not erase the higher derived
coherence of the deck action retained in Proposition~
\ref{prop:frag-actual-finite-galois-descent}.
\end{proof}

\begin{theorem}[an actual algebraic torus star comparison]
\label{thm:frag-algebraic-torus-star}
Use the algebraic product of Lemma~\ref{lem:frag-algebraic-torus-weights}.
Let $\mathcal E$ be the Kummer local system associated to a finite
$K_{n_0}$-action on a constructible base coefficient $E$.  The actual
geometric torsors $q_n:U_n\to U$ supply a canonical comparison
\[
 C^*_{\mathrm{cts}}(\Gamma,E)
 \xrightarrow{\sim} R\pi_*\mathcal E.
\]
It extends to bounded derived Kummer coefficients by finite derived
action descent and is compatible with their sheaf attaching maps.
On the unit it is a comparison of commutative algebra objects; on
coefficients it is a module comparison.  On selected cartesian adic
systems the outputs have bounded derived-cartesian recovery.
This theorem is for the algebraic product $T\times\mathbf A^e$;
it does not identify an analytic polydisc specialization with that
algebraic projection.
\end{theorem}

\begin{proof}
For $n\ge n_0$, torsor descent gives the equivariant trivialization
$q_n^*\mathcal E=p_n^*E_n$.  Proposition~
\ref{prop:frag-actual-finite-galois-descent} constructs compatible maps
\[
C^*(K_n,E_n)\longrightarrow
F_n:=\operatorname{Tot}C^*(K_n,Rp_{n,*}p_n^*E_n)
 \simeq R\pi_*\mathcal E.
\]
Take their filtered colimit.  The left side is continuous cochains;
the right sides are a compatible diagram of equivalences with the
fixed actual direct image.  It remains to show that the resulting
map is an equivalence, rather than merely that its endpoints have
the same expected cohomology.

Use the quotient Postnikov tower $\tau_{\ge j}$ of
$Rp_{n,*}p_n^*E_n$, with arrows from $j$ to $j+1$.  It has length
$d+1$ by Lemma~\ref{lem:frag-algebraic-torus-weights}.  After derived
invariants the fibres of its successive arrows are
\[
 C^*(K_n,E_n\otimes\bigwedge^j(M^\vee\otimes R_r)(-j))[-j].
\]
The transition on the positive-$j$ piece is inflation followed by
multiplication by $\ell^{j(m-n)}$, and is zero after sufficiently
large refinement at fixed $r$.  The degree-zero comparison induced
by the unit is the identity on $E_n$.  Exactness of filtered colimits
and the finite length of this filtration therefore prove that the
constructed map is an equivalence.  This argument uses neither
formality of the finite derived deck complex nor interchange with
an arbitrary infinite totalization.

The associated-sheaf functors at finite level are exact on the
coefficient abelian categories and compatible with inflation.  The
finite derived descent of Theorem~\ref{thm:frag-relative-finite-action-descent}
therefore extends the construction to the stated derived source,
including its global sheaf maps.  Equivalence on bounded objects
follows by finite truncation triangles in that larger constructible
category, without requiring their cohomology objects to be perfect.
The multiplicative/module structure is inherited from the actual
unit maps and descent diagrams of Proposition~\ref{prop:frag-actual-finite-galois-descent}.  Filtered colimits preserve
these commutative algebra and module structures.  Finally the
continuous augmentation Koszul calculation identifies the already
compared actual output with a finite construction on $E_r$, proving
uniform boundedness and derived coefficient recovery on the selected
adic class.  Rationalization is performed only afterwards.
\end{proof}


\begin{proposition}[ordinary nearby comparison away from the generic boundary]
\label{prop:frag-ordinary-nearby-boundary-disjoint}
Let $X/R$ be the chosen algebraization, $Y\subset X_s$ the selected
locally closed subscheme, and $j:U\hookrightarrow X_\eta$ the supplied
open.  Assume
\[
Y\cap\overline{X_\eta\setminus U}=\varnothing
\quad\text{inside }X.
\]
Let $\mathcal E$ be a bounded constructible $R_r$-complex on $U$.
For ordinary geometric nearby cycles, retaining the full inertia
action, there is a canonical comparison equivalence
\[
 i_Y^*R\Psi_X(Rj_*\mathcal E)
 \xrightarrow{\sim}
 R\operatorname{sp}_{Y,\bar\eta,*}
       (\mathcal E^{\mathrm{an}}|_{]Y[_{\bar\eta}}).
\]
After the structural direct image from $Y$ to the face base this
compares the analytic star functor with the \emph{ordinary} nearby
star functor.  It is natural on algebraized finite Kummer diagrams
for which the displayed boundary condition holds.  It does not by
itself replace ordinary nearby cycles by tame nearby cycles, and it
asserts no compact-support or incidence comparison.
\end{proposition}

\begin{proof}
Remove the closure of $X_\eta\setminus U$ from $X$, then choose an
open neighbourhood on which $Y$ is closed.  This leaves the formal
completion along $Y$ unchanged, and the entire generic fibre of this
neighbourhood lies in $U$.  On this neighbourhood the coefficient
$Rj_*\mathcal E$ restricts to $\mathcal E$ on the generic fibre.
This explains the use of the boundary condition; the original
$U$ has not silently been replaced by the whole generic fibre of $X$.

Use the standard comparison morphism from algebraic nearby cycles
restricted to $Y$ to the formal nearby cycles of the completion
$\widehat X_{/Y}$.  It is induced by restricting algebraic etale
neighbourhoods to their completions and then to their analytic
generic fibres.  Those restriction maps give compatible maps of
coefficient complexes and hence a derived comparison morphism.
Berkovich's comparison theorem for completion along a subscheme
\cite[Theorem 3.1]{BerkovichFormalII} identifies its cohomology
sheaves for prime-to-residue-characteristic constructible torsion
coefficients.  Here the algebraization is finite type over the complete valuation
ring $R$, the localized $Y$ is closed in its special fibre, and
$\widehat X_{/Y}$ is a special formal completion, as in that theorem.
The coefficient is algebraic constructible $R_r$-torsion on the
localized generic fibre, so its torsion order is invertible in $k$.
The theorem identifies the canonical map from restricted algebraic
nearby cycles to the formal nearby functor on this completion.
Proposition 1.3 identifies its generic fibre with the tube.  Geometric base change
to the chosen completed algebraic closure gives the displayed
geometric version.  Finite truncations give bounded complexes.

The formal nearby functor is the actual geometric specialization
direct image on the generic fibre of $\widehat X_{/Y}$, and that
fibre is the tube $]Y[_{\bar\eta}$.  This identifies the target in
the display.  Naturality is the naturality of the restriction maps
used to construct the comparison, so applies to the stated finite
Kummer diagrams.  Applying $Rf_{Y,*}$ proves the facewise assertion.
The proposition uses this published ordinary comparison in its
specified domain.  Proposition~\ref{prop:frag-completion-comparison-monoidal}
identifies this construction with a tensor-compatible specialization
mate and supplies its algebra/module enhancement; this is not inferred
solely from the isomorphisms on cohomology sheaves.
\end{proof}

\begin{proposition}[the tensor structure of completion comparison]
\label{prop:frag-completion-comparison-monoidal}
The ordinary comparison of
Proposition~\ref{prop:frag-ordinary-nearby-boundary-disjoint} is a
morphism of lax symmetric-monoidal derived functors with their
coefficient-module actions.  It is natural for the algebraized finite
Kummer diagrams in that proposition and for coefficient change wherever
the derived comparison is defined.  This statement supplies the
multiplicative clause of Proposition~\ref{prop:frag-ordinary-nearby-boundary-disjoint} in its stated geometric domain.
\end{proposition}

\begin{proof}
\emph{Topoi and the inverse-image transformation.}
Localize as in Proposition~\ref{prop:frag-ordinary-nearby-boundary-disjoint},
so that $Y$ is closed and the generic boundary is absent.  Put $A=R_r$.
All categories in this proof are the stable localizations
$D_A(\mathcal Z)$ of complexes of $A$-module sheaves on the indicated
topos; functors are derived unless they are exact inverse images.
The directions of the geometric morphisms are
\[
\begin{array}{ccc}
 (X_{\bar\eta})_{\acute et}&\xrightarrow{a}&X_{\acute et}
                           \xleftarrow{i}Y_{\acute et},\\
 T_{\acute et}&\xrightarrow{c}&(X_{\bar\eta})_{\acute et},\\
 T_{\acute et}&\xleftarrow{\mu}&T_{q\acute et}
                           \xrightarrow{\nu}Y_{\acute et}.
\end{array}
\]
Here $a$ is geometric generic inclusion, $c$ is analytification
restricted to the tube, and $(\mu,\nu)$ is exactly
Definition~\ref{def:frag-analytic-specialization-sites}.
Thus the two inverse-image functors below both have domain
$D_A(X_{\acute et})$ and target $D_A(T_{q\acute et})$:
\[
 \sigma:\nu^*i^*\longrightarrow\mu^*c^*a^*.
\]
We specify its monoidal structure before taking a mate.
For an etale $V\to X$, complete $V$ along $V_Y$.  Formal etale
lifting identifies its generic fibre over $T$ with the quasi-etale
object defining $\nu^*h_{V_Y}$.  The canonical map from that generic
fibre to $V_{\bar\eta}^{\mathrm{an}}$ defines $\sigma_{h_V}$.
For a finite family $V_1,\ldots,V_t$, completion and generic fibre
identify the lift of $\prod_XV_j$ with the product of the lifts.
The restriction map for this product is the product of the
restriction maps, and the map for $V=X$ is the unit map.
For a permutation or a substitution of finite families, both
composites are this same restriction of etale objects.  This gives
the finite-arity symmetry, unit and substitution identities, not
merely a binary cup-product identity.

These constructions respect etale coverings and morphisms, so they
extend by sheafification to inverse images on sheaves of sets.
On free module generators they give
\[
 \sigma_{A[h_{V_1\times_X\cdots\times_XV_t}]}
           =\bigotimes_{j=1}^t\sigma_{A[h_{V_j}]}.
\]
This is a natural transformation of exact inverse-image functors
on module sheaves.  Applying it degreewise to complexes and localizing
at quasi-isomorphisms gives a specified transformation in the chosen
derived enhancement: both functors preserve quasi-isomorphisms, and
the naturality squares descend with them.
To compute its tensor structure, resolve a complex of presheaves by
a semi-free complex, with cells sums of shifts of free representables,
and then sheafify.  Exact sheafification gives a resolution of the
original sheaf complex.  Its cell filtration makes it K-flat: free
representable sheaves are flat, tensoring an acyclic complex with
each cell is acyclic, and filtered colimits of module sheaves are exact.
The inverse images of these resolutions have the same property,
since inverse image preserves free modules, sums and these filtrations.
Thus the sheaf-level finite-arity identities above compute the derived
tensor identities on these resolutions.  They are compatible with
quasi-isomorphisms and hence descend, together with symmetry and
substitution, to the localization.  This defines the symmetric-monoidal
transformation $\sigma$; it need not be an equivalence.

\emph{The mate, including its multiplication.}
Write $B=\mu^*c^*$, $Q=i^*Ra_*$ and $P=R\nu_*B$, functors from
$D_A((X_{\bar\eta})_{\acute et})$ to $D_A(Y_{\acute et})$ for $Q,P$.
The composite in the quasi-etale category is
\[
 \psi_F:\nu^*QF\xrightarrow{\sigma_{Ra_*F}}
           B a^*Ra_*F\xrightarrow{B\epsilon_{a,F}}BF.
\]
Its adjoint is $\tau_F:QF\to PF$.
The lax products on $Ra_*$ and $R\nu_*$ are the mates of their
strong-monoidal inverse images; explicitly they are adjoint to
the tensor products of the corresponding counits.
For a finite family $F_1,\ldots,F_t$, the two maps
\[
 \bigotimes_j\nu^*QF_j\longrightarrow B(\bigotimes_jF_j)
\]
obtained either by multiplying in $Q$ and then using $\psi$, or by
using the individual $\psi_{F_j}$ and multiplying in $B$, agree.
Indeed substitute the formula for $\psi$: the $\sigma$ portions
agree by its finite-arity product identity, and the remaining
identity is the defining counit formula for the lax product of
$Ra_*$.  The empty family checks the unit in the same way.
Adjunction with $\nu^*$ turns these equalities into exactly the
lax-monoidal squares for $\tau$.  Permutations and successive
substitutions give the associativity and symmetry cells because
their inverse-image identities were already the same restriction
maps.  This is also the monoidal calculus of mates of
\cite[Theorem 3.4.7]{HHLNMates}, which preserves composition of
these transformations.  It gives algebra and coefficient-module
compatibility, with no appeal to equality of cup products on cohomology.

\emph{Identification of the map to which comparison applies.}
For each etale neighbourhood $V\to X$, the map defining $\psi$
restricts a section on $V_{\bar\eta}$ to the generic fibre of its
completion along $V_Y$.  The last arrow is precisely evaluation
of the generic pullback/direct-image counit on that section.
More explicitly, choose a K-injective representative $I$ of $F$.
Since $a^*$ is exact, $a_*I$ computes $Ra_*F$.  The chain map
\[
 \nu^*i^*a_*I\xrightarrow{\sigma}B a^*a_*I
       \xrightarrow{B\epsilon}BI
\]
represents $\psi_F$.  Resolve $BI\to J$ by a K-injective complex
on the quasi-etale topos and take the ordinary $\nu^*\dashv\nu_*$
adjoint of this composite.  The resulting map
$i^*a_*I\to\nu_*J$ represents $\tau_F$.
No preservation of K-injectives by inverse image is needed.
On etale neighbourhoods it is exactly restriction to the completed
generic fibre followed by the resolution map.  These are the
restriction maps defining the algebraic-to-formal nearby comparison
in \cite[\S\S2--3, Theorem 3.1]{BerkovichFormalII}; the construction
uses the same generic pullback and specialization topoi.
Derived adjunction identifies their mate with $\tau$, independently
of resolutions.  Thus the comparison theorem applies to this specified
map, not to an isomorphism chosen from its cohomology groups.
That theorem supplies invertibility
on the stated constructible prime-to-residue-characteristic source;
the preceding construction supplies its multiplicative enhancement.

For a morphism of the supplied algebraized finite Kummer diagrams,
restriction to completion composes with the induced map of generic
fibres.  The corresponding $\sigma$ and counit squares therefore
paste to the same $\psi$, and taking mates gives naturality of $\tau$.
Coefficient extension carries free module restriction maps and
evaluation to their coefficient-changed versions.  The same argument
gives the coefficient-change squares, without declaring arbitrary
right-direct-image base-change maps invertible.  Neither an
invertible generic/closed base-change square nor an exceptional
specialization functor is asserted.
\end{proof}

\begin{lemma}[the precise passage from ordinary to tame nearby cycles]
\label{lem:frag-wild-inertia-criterion}
Let $P$ be wild inertia, a pro-$p$ group, and $\ell\ne p$.  On smooth
$P$-action sheaves of $R_r$-modules, invariants are exact.  For a
bounded complex $C$ with such an action, the canonical map
\[
 C^P\longrightarrow C
\]
is an equivalence if $P$ acts trivially on every cohomology sheaf of
$C$.  Accordingly the ordinary comparison in Proposition~
\ref{prop:frag-ordinary-nearby-boundary-disjoint} gives the desired
tame comparison only after this wild-triviality condition is proved
for the relevant nearby complex (or if no wild inertia is present).
\end{lemma}

\begin{proof}
To lift an invariant local section across an epimorphism, first choose
a local lift.  Smoothness makes its orbit factor locally through a
finite $p$-group.  Averaging that lift over this group is allowed in
$R_r$, since its order is invertible, and gives an invariant lift.
This proves exactness as a sheaf statement.  Thus
$H^j(C^P)=H^j(C)^P$, and the assumed triviality makes the canonical
map an isomorphism on every cohomology sheaf.  This proves the
criterion.  It does not prove wild-triviality from the word
``Kummer'' or from a coincidence of relative inertia groups.
\end{proof}

\begin{lemma}[Derived direct image under localization of topoi]
\label{lem:frag-slice-derived-base-change}
Work with the etale and Kummer-etale sheaf topoi and their limiting
localizations used here.  Let $f:\mathcal E\to\mathcal F$ be a
morphism of these topoi, and let $U$ be
an object of $\mathcal F$, and use constant coefficients $R_r$.
Write $p:\mathcal F/U\to\mathcal F$ and
$p':\mathcal E/f^*U\to\mathcal E$ for the localizations and
$f_U:\mathcal E/f^*U\to\mathcal F/U$ for the induced morphism.
For $K\in D^+(\mathcal E,R_r)$ the canonical base-change map is an
isomorphism
\[
 p^*Rf_*K\xrightarrow{\sim}Rf_{U,*}p'^*K.
\]
It identifies the pullback of the adjunction unit of $f$ with that
of $f_U$ and is compatible with further localization.
\end{lemma}

\begin{proof}
Evaluation on $V\to U$ identifies both underived direct images with
evaluation on $f^*V$.  The exact restriction $p'^*$ preserves
injectives because its left adjoint $p'_!$, the sheafified direct
sum over the slicing object, is exact.  Restricting a bounded-below
injective resolution therefore proves the displayed derived exchange.
The same evaluation identity identifies units; iterated slicing
gives compatibility with composition and its cells.
\end{proof}

\begin{lemma}[The derived analytic site model and its comparison map]
\label{lem:frag-derived-analytic-site-model}
Let $T$ be the analytic tube in the stated geometric source domain.
At precision $r$ use the stable infinity-categories
$D(\operatorname{Sh}(T_{\acute et},R_r))$ and
$D(\operatorname{Sh}(T_{q\acute et},R_r))$, obtained by localizing
complexes at quasi-isomorphisms.  K-injective representatives compute
right adjoints and derived Hom; K-flat representatives compute tensor.
For bounded constructible $F$ the canonical unit map
\[
 R\Gamma(T_{\acute et},F)\longrightarrow
 R\Gamma(T_{q\acute et},\mu^*F)
\]
is an equivalence in $D(R_r)$.  It respects evaluation for dualizable
coefficients and pullback by finite etale torsors.
\end{lemma}

\begin{proof}
The exact strong tensor functor $\mu^*$ induces a functor on the
stable localizations with its derived adjunction unit.  For an ordinary
etale abelian sheaf, the canonical comparison on
sections and cohomology is \cite[Theorem 3.3(i)--(ii)]{BerkovichFormalI},
applied to the identity quasi-etale map $T\to T$.
This published theorem applies to analytic spaces without a
compactness assumption.  Its degree-zero map is the inverse-image
unit; the higher maps are its derived cohomology maps.  Finite
truncation triangles therefore give the asserted equivalence for
bounded complexes.  Thus this bridge does not depend on the later
unpublished equivariant extension.  The tensor and mapping assertions
use the following argument in the selected derived enhancement.
For dualizable $E$, $R\mathcal Hom(E,F)=E^\vee\otimes^LF$, and
$R\operatorname{Hom}(E,F)=R\Gamma R\mathcal Hom(E,F)$ determines
the mapping spectrum.  Strong tensor inverse images preserve evaluation.
Torsor site maps, their units, and all Cech face and degeneracy maps
pass to these same localizations.  The exact tensor associated-sheaf
functors likewise derive and commute with inflation before any
full-faithfulness claim.  Finite-action descent defines their continuous
realization; Lemma~\ref{lem:frag-B-local-mapping-unit} checks its mapping map.
For cartesian adic systems, the universal property of the category
limit computes mapping spectra as limits of the finite-precision
mapping spectra.  Its selected full subcategory has the same maps.
\end{proof}

\section{Two source obstructions}
\label{sec:source-obstructions}

\begin{proposition}[An obstruction to the unrestricted common coefficient source]
\label{prop:frag-constructible-source-obstruction}
Let $X=\operatorname{Spec}R[t]$, $Y=X_s=\mathbf A^1_k$, with $k$
algebraically closed, and let $T$ be the geometric closed unit-disc
generic fibre of $\widehat X_{/Y}$.  Set
$L=i_{0,*}R_r$ on $Y_{\acute et}$.  This is a constructible coefficient
of finite Tor-amplitude, allowed by the earlier stratified coefficient
condition.  There is no algebraically constructible bounded complex
$F$ on $X_{\bar\eta}$ satisfying
\[
 \mu^*(F^{\mathrm{an}}|_T)\simeq\nu^*L
 \quad\text{on }T_{q\acute et}.
\]
Thus the unrestricted prescription ``pull back every constructible base
coefficient'' cannot also require that same analytic coefficient to
come from a common algebraically constructible generic source.
\end{proposition}

\begin{proof}
At a classical geometric point $t=a$ with $|a|\le1$, the right side
has degree-zero stalk $R_r$ if $\bar a=0$, and zero if $\bar a\ne0$.
These are the stalks of the actual inverse image under $\nu$.
If the displayed equivalence existed, exactness of the inverse images
would impose these stalk values on the constructible sheaf $H^0(F)$.
A constructible sheaf on the algebraic affine line is locally constant
on the complement of a finite set.  That complement is connected, so
the cardinality of its stalks is constant.  However, both sequences
\[
 a_m=\pi^m,\qquad b_m=1+\pi^m,\qquad m\ge1,
\]
lie in $T$ and contain infinitely many distinct points.  Outside any
finite exceptional set their required stalks have respectively
cardinality $|R_r|$ and $1$.  This is a contradiction.  The argument
uses only classical stalks as a necessary condition; it does not claim
that they alone classify all quasi-etale sheaves.

This is a source obstruction, not a counterexample to the already
proved star theorem on an algebraized lisse source.  Choosing instead
an algebraic skyscraper at $t=0$ changes the analytic coefficient: its
support is a single generic point rather than the whole residue disc.
Equal nearby outputs would not identify those two sources.  Nor does
the proposition assert that a comparison formulated entirely on
formal or quasi-etale coefficients is impossible; that is a different
realization datum which must be defined and proved separately.
\end{proof}

\begin{proposition}[The standard Artin fan is not a third realization]
\label{prop:frag-artin-obstruction}
The local comparison cannot be extended by replacing the relative Kummer face with
the ordinary unrigidified Artin fan.  For $P=\mathbf N^2$ and base direction
$\rho=e_1+e_2$, the closed face of
$[\mathbf A^2/\mathbf G_m^2]$ retains a residual $B\mathbf G_m$ after
stacky quotienting by the base diagonal.  At a strict geometric point
$\bar s\to S$, its lisse--\'{e}tale \(\ell\)-adic cohomology is
\[
 H^{2j}_{\mathrm{lis\acute et}}
 (B_{\bar s}\mathbf G_m,\mathbf Q_\ell)
 \simeq\mathbf Q_\ell(-j),
 \qquad H^{2j+1}_{\mathrm{lis\acute et}}=0,
\]
with the universal Chern class in
$H^2(B\mathbf G_m,\mathbf Q_\ell(1))$.  In contrast, the relative Kummer
cochain complex at $\bar s$ is finite:
\[
 \operatorname{Sym}\bigl((\mathbf Q_\ell(-1))[-1]\bigr)
 \simeq \mathbf Q_\ell\oplus\mathbf Q_\ell(-1)[-1].
\]
Thus the ordinary Artin-fan leg is not equivalent to the analytic or tame
nearby leg, even on the constant coefficient.
\end{proposition}

\begin{proof}
The quotient description leaves the residual diagonalizable stabilizer
$\mathbf G_m$.  Its ordinary classifying-stack cohomology is polynomial in
the universal Chern class, hence unbounded, while the Kummer cochain model
of the relative characteristic lattice has cohomological dimension one.
The graded dimensions differ, so no equivalence can exist.  Rigidifying the
base direction and retaining the Kummer site is exactly what removes this
obstruction in the relative root construction.
\end{proof}

\section*{Acknowledgments}

This mathematical paper was generated 100\% by artificial intelligence.  I
first acknowledge the textual data produced by humanity as a whole, which
provides the linguistic and cultural record from which mathematical expression
can be learned.  I then thank the accumulated scholarly achievements of
mathematicians throughout history.  I also thank the researchers whose work has
advanced and refined neural networks through successive generations.  Finally,
I thank the OpenAI team for the training and development of the models used in
the preparation of this paper.

\begin{thebibliography}{99}

\bibitem{StacksEtaleSpaceRealization}
The Stacks Project Authors,
\emph{The Stacks Project}, Lemma 66.27.3,
\url{https://stacks.math.columbia.edu/tag/0GF6}.


\bibitem{HHLNMates}
R.~Haugseng, F.~Hebestreit, S.~Linskens and J.~Nuiten,
\emph{Lax monoidal adjunctions, two-variable fibrations and the calculus
of mates}, arXiv:2011.08808, Theorem 3.4.7;
\url{https://arxiv.org/abs/2011.08808}.

\bibitem{RozenblyumFilteredColimits}
N.~Rozenblyum,
\emph{Filtered colimits of infinity-categories}, note, 30 December 2012,
Lemma 0.2.1 and \S0.5;
\url{https://www.math.toronto.edu/nick/notes/colimits.pdf}.

\bibitem{StacksSpaceLimitDescent}
The Stacks Project Authors,
\emph{The Stacks Project}, Lemmas 70.7.1, 70.6.2, and 70.6.4,
\url{https://stacks.math.columbia.edu/tag/07SK},
\url{https://stacks.math.columbia.edu/tag/07SL},
\url{https://stacks.math.columbia.edu/tag/07SN}.


\bibitem{StacksNoetherianTorsion}
The Stacks Project Authors,
\emph{The Stacks Project}, Lemma 47.10.1,
\url{https://stacks.math.columbia.edu/tag/0955}.


\bibitem{StacksConstructibleNoetherian}
The Stacks Project Authors,
\emph{The Stacks Project}, Sections 59.73--59.74:
Lemma 59.73.2 and Proposition 59.74.1, Lemma 59.74.2,
\url{https://stacks.math.columbia.edu/tag/03SA},
\url{https://stacks.math.columbia.edu/tag/03RY}.

\bibitem{StacksDerivedSerre}
The Stacks Project Authors,
\emph{The Stacks Project}, Lemma 13.17.4,
\url{https://stacks.math.columbia.edu/tag/06UP}.


\bibitem{BerkovichFormalII}
V.~G. Berkovich,
\emph{Vanishing cycles for formal schemes. II},
Invent. Math. \textbf{125} (1996), 367--390;
\url{https://www.wisdom.weizmann.ac.il/~vova/Inven_1996_125_formalII.pdf}.

\bibitem{BerkovichFormalI}
V.~G. Berkovich,
\emph{Vanishing cycles for formal schemes},
Invent. Math. \textbf{115} (1994), 539--571,
Theorem 3.3(i)--(ii);
\url{https://www.wisdom.weizmann.ac.il/~vova/Inven_1994_115_formal.pdf}.

\bibitem{LuZhengNearbyDuality}
Q.~Lu and W.~Zheng,
\emph{Categorical traces and a relative Lefschetz--Verdier formula},
Forum Math. Sigma \textbf{10} (2022), e10,
\S3.1, especially Corollaries 3.8--3.9;
\url{https://doi.org/10.1017/fms.2022.2}.

\bibitem{BerkovichEtale}
V.~G. Berkovich,
\emph{\'{E}tale cohomology for non-Archimedean analytic spaces},
Publ. Math. IH\'{E}S \textbf{78} (1993), 5--161.

\bibitem{NakayamaNearby}
C.~Nakayama,
\emph{Nearby cycles for log smooth families},
Compositio Math. \textbf{112} (1998), 45--75.

\bibitem{FujiwaraKatoNakayama}
L.~Illusie,
\emph{An overview of the work of K. Fujiwara, K. Kato, and C. Nakayama on
logarithmic \'{e}tale cohomology},
Ast\'{e}risque \textbf{279} (2002), 271--322;
\url{https://www.numdam.org/article/AST_2002__279__271_0.pdf}.

\bibitem{TalpoVistoliInfiniteRoot}
M.~Talpo and A.~Vistoli,
\emph{Infinite root stacks and quasi-coherent sheaves on logarithmic
schemes},
Proc. Lond. Math. Soc. (3) \textbf{116} (2018), no.~5, 1187--1243;
\url{https://arxiv.org/abs/1410.1164}.

\bibitem{TalpoVistoliKN}
M.~Talpo and A.~Vistoli,
\emph{The Kato--Nakayama space as a transcendental root stack},
arXiv:1611.04041 (2016);
\url{https://arxiv.org/abs/1611.04041}.

\bibitem{ScherotzkeSibillaTalpoMcKay}
S.~Scherotzke, N.~Sibilla, and M.~Talpo,
\emph{On a logarithmic version of the derived McKay correspondence},
Compos. Math. \textbf{154} (2018), no.~12, 2534--2585;
\url{https://arxiv.org/abs/1612.08961}.

\bibitem{LaszloOlssonFinite}
Y.~Laszlo and M.~Olsson,
\emph{The six operations for sheaves on Artin stacks I: finite coefficients},
Publ. Math. Inst. Hautes \'{E}tudes Sci. \textbf{107} (2008), 109--168;
\url{https://numdam.org/item/PMIHES_2008__107__109_0/}.

\bibitem{LaszloOlsson}
Y.~Laszlo and M.~Olsson,
\emph{The six operations for sheaves on Artin stacks II: adic coefficients},
Publ. Math. Inst. Hautes \'{E}tudes Sci. \textbf{107} (2008), 169--210.

\bibitem{LiuZhengEnhanced}
Y.~Liu and W.~Zheng,
\emph{Enhanced six operations and base change theorem for higher Artin
stacks}, arXiv:1211.5948, version of 17 December 2024;
\url{https://arxiv.org/abs/1211.5948}.

\bibitem{LiuZhengAdic}
Y.~Liu and W.~Zheng,
\emph{Enhanced adic formalism and perverse $t$-structures for higher Artin
stacks}, arXiv:1404.1128;
\url{https://arxiv.org/abs/1404.1128}.

\bibitem{StacksProject}
The Stacks Project Authors,
\emph{The Stacks Project},
chapters on limits of algebraic spaces and algebraic stacks, especially
Tags 07SK, 07SL, 07SN, 0ATT, and 0A9Z;
\url{https://stacks.math.columbia.edu}.

\bibitem{BrumerPseudocompact}
A.~Brumer,
\emph{Pseudocompact algebras, profinite groups and class formations},
J. Algebra \textbf{4} (1966), 442--470.

\bibitem{CarchediEtAlLogHomotopy}
D.~Carchedi, S.~Scherotzke, N.~Sibilla, and M.~Talpo,
\emph{On the profinite homotopy type of log schemes},
Adv. Math. \textbf{460} (2025), 110018;
\url{https://doi.org/10.1016/j.aim.2024.110018}.
The published version's Sections 5--6 compare actual and pro-root
models at the level of profinite shape.

\bibitem{EsnaultKerzPencils}
H.~Esnault and M.~Kerz,
\emph{Semi-stable Lefschetz pencils},
arXiv:2311.15886, Appendix~A.2.

\end{thebibliography}

\end{document}
